% ============================================================================
% Fichier     : SGSI_main.tex
% Titre       : Systèmes de Gestion de la Sécurité Informatique
% Auteur      : MINKA MI NGUIDJOÏ Thierry Emmanuel
% Institution : Chaire UNESCO d'Afrique Centrale Accès TIC : ENSP / LIMSI
% Version     : 0.1 -- octobre 2026 (squelette)
% ============================================================================

\documentclass[12pt, twoside, openright, a4paper]{book}

% --- Feuille de style unifiée ---
\usepackage{Style_SGSI}

% --- Bibliographie ---
\addbibresource{bibliographie/SGSI_references.bib}

% ============================================================================
% MÉTADONNÉES PDF
% ============================================================================
\hypersetup{
    pdftitle={Systèmes de Gestion de la Sécurité Informatique},
    pdfauthor={MINKA MI NGUIDJOÏ Thierry Emmanuel},
    pdfsubject={SIS498 · Master Professionnel en Sécurité des Systèmes Informatiques et des Données · Université de Garoua},
    pdfkeywords={SMSI, ISO 27001, ISO 27002, gouvernance, EBIOS, appréciation des risques, PCA, audit, SGSI}
}

\begin{document}

% --------------------------------------------------------------------------
% ÉLÉMENTS PRÉLIMINAIRES (pages romaines)
% --------------------------------------------------------------------------
\frontmatter

    % ==============================================================================
% Fichier : front/00_page_titre.tex
% Titre   : Page de titre - \titreCours
% Auteur  : MINKA MI NGUIDJOÏ Thierry Emmanuel
% Version : 1.1 -- novembre 2026 (charte des cours, modèle de HE2)
% ==============================================================================
% Logos : Université de Garoua à gauche et à droite.
% ==============================================================================

\begin{titlepage}
\thispagestyle{empty}

\begin{tikzpicture}[remember picture, overlay]
    \fill[CouleurPrimaire]
        (current page.north west) rectangle
        ([yshift=-3.6cm] current page.north east);
    \fill[CouleurAccent]
        ([yshift=-3.6cm] current page.north west) rectangle
        ([yshift=-4.0cm] current page.north east);
    \fill[CouleurPrimaire]
        (current page.south west) rectangle
        ([yshift=2.5cm] current page.south east);
    \fill[CouleurAccent]
        ([yshift=2.5cm] current page.south west) rectangle
        ([yshift=2.9cm] current page.south east);
    % En-tête institutionnel, posé dans le bandeau supérieur
    \node[align=center, text=white] at ([yshift=-1.8cm] current page.north) {%
        {\fontsize{11}{13}\selectfont
        \textsc{République du Cameroun}}\\[0.3ex]
        {\color{CouleurAccent!60!white}\rule{0.4\paperwidth}{0.4pt}}\\[0.6ex]
        {\fontsize{10}{12}\selectfont\textbf{\tutelle}}\\[0.2ex]
        {\fontsize{9}{11}\selectfont \faculte{} · \departement}\\[0.2ex]
        {\fontsize{8}{10}\selectfont \unitepedago}};
    % Logos : l'Université de Garoua à gauche et à droite (un seul établissement)
    \node[fill=white, rounded corners=4pt, inner sep=2pt]
        at ([xshift=1.75cm, yshift=-1.8cm] current page.north west)
        {\includegraphics[height=2.6cm]{figures/logos/garoua_logo.png}};
    \node[fill=white, rounded corners=4pt, inner sep=2pt]
        at ([xshift=-1.75cm, yshift=-1.8cm] current page.north east)
        {\includegraphics[height=2.6cm]{figures/logos/garoua_logo.png}};
    % Année et version, posées dans le bandeau inférieur
    \node[align=center] at ([yshift=1.25cm] current page.south) {%
        {\color{white}\bfseries\large Année académique \anneescolaire}\\[0.3ex]
        {\color{CouleurAccent!70!white}\small
        Version 1.0 | Document pédagogique de l'UE \codeUE}};
\end{tikzpicture}

\vspace*{0.6cm}

% --- Titre principal (sur deux lignes) ---
\begin{center}
    {\color{CouleurPrimaire}\rule{0.85\textwidth}{1.2pt}}\\[1.2ex]
    {\fontsize{24}{30}\selectfont\bfseries\color{CouleurPrimaire}
        Systèmes de Gestion de la\\[0.4ex]
        Sécurité Informatique}\\[0.8ex]
    {\color{CouleurAccent}\rule{0.85\textwidth}{2pt}}\\[0.5ex]
    {\color{CouleurPrimaire}\rule{0.85\textwidth}{0.5pt}}
\end{center}

\vspace{0.8cm}

\begin{center}
    {\fontsize{13}{16}\selectfont\itshape\color{CouleurSecondaire}
        Du risque à la certification~:\\
        piloter la sécurité de l'information}
\end{center}

\vspace{0.5cm}

% --- Cartouche sigle et code UE ---
\begin{center}
    \begin{tikzpicture}
        \node[draw=CouleurPrimaire, fill=CouleurPrimaire, rounded corners=6pt,
              inner xsep=14pt, inner ysep=6pt]{%
            \color{white}\bfseries\large \sigleCours{} - UE~\codeUE
        };
    \end{tikzpicture}
\end{center}

\vspace{1.0cm}

% --- Schéma de la discipline : le cycle d'amélioration continue du SMSI ---
\begin{center}
\begin{tikzpicture}[
    f/.style = {draw=CouleurBordDef, fill=CouleurDef, rounded corners=3pt,
        minimum width=2.2cm, minimum height=0.8cm, align=center,
        font=\small\bfseries\color{CouleurPrimaire}},
    arrow/.style = {-{Stealth[length=4pt]}, thick, color=CouleurSecondaire}]
\foreach \t [count=\i] in {Planifier, Déployer, Contrôler, Améliorer}
    \node[f] (n\i) at (3.3*\i - 3.3, 0) {\t};
\foreach \i [evaluate=\i as \j using int(\i+1)] in {1,2,3}
    \draw[arrow] (n\i) -- (n\j);
\draw[arrow] (n4.south) to[bend left=18] node[below, font=\footnotesize\itshape,
    text=CouleurSecondaire] {amélioration continue, guidée par les risques} (n1.south);
\end{tikzpicture}
\end{center}

\vspace{0.6cm}

\begin{center}
    {\large\bfseries\color{CouleurPrimaire}
        \auteurcours}\\[0.5ex]
    {\normalsize\color{CouleurSecondaire}
       Enseignant Chercheur | \tutelle}\\[0.2ex]
    {\small\color{CouleurBordInfo}
        \faculte{} · \departement}\\[0.2ex]
    {\small\color{CouleurBordInfo}
        \href{mailto:maletyon@proton.me}{maletyon@proton.me}}
\end{center}

\end{titlepage}

    \cleardoublepage

    % ==============================================================================
% Fichier : front/01_avant_propos.tex : SGSI (SIS498)
% ==============================================================================

\chapter*{Avant-propos}
\addcontentsline{toc}{chapter}{Avant-propos}
\markboth{Avant-propos}{Avant-propos}

Ce livre est le support de l'unité d'enseignement \textbf{SIS498, Systèmes de
Gestion de la Sécurité Informatique}, du Master Professionnel en Sécurité des
Systèmes Informatiques et des Données de l'Université de Garoua (Faculté des
Sciences, Département de Mathématiques et Informatiques).

\medskip
Là où d'autres enseignements apprennent à détecter, à investiguer ou à
attaquer, \textbf{SIS498} apprend à \textbf{décider} la sécurité et à la
\textbf{tenir dans la durée}~: gouverner, apprécier les risques, mettre en
\oe uvre les mesures, contrôler et améliorer. Le fil conducteur est unique~:
la \textbf{mairie d'Eséka}, dans le département du Nyong-et-Kéllé, décide,
après un incident, de bâtir un système de management de la sécurité de
l'information et de le conduire jusqu'à la certification. Le cours suit cette
démarche à travers quatre phases situées dans la région~: \textbf{NYONG}
(gouvernance), \textbf{KELLÉ} (risques), \textbf{MAKAK} (mise en \oe uvre) et
\textbf{MATOMB} (audit).

\medskip
Chaque notion est ancrée à la fois dans les standards internationaux (famille
\textbf{ISO/IEC 27000}, en particulier 27001, 27002 et 27005, ISO 31000, ISO
22301, ISO 19011, NIST CSF, EBIOS Risk Manager) et dans le droit camerounais
et africain (lois n°~2010/012 et 2024/017, Convention de Malabo). Un fil
directeur parcourt le livre~: le cycle \textbf{PDCA} (planifier, déployer,
contrôler, agir), par lequel un système de management se construit et
s'améliore.

\medskip
Les cas sont fictifs~: les noms, les dates et les montants le sont~; les
mécanismes de gestion et les références normatives sont exacts. Le livre
s'accompagne de l'environnement \textbf{NDAP SUKULU}, qui fait pratiquer
chaque chapitre.

\bigskip
\noindent \IfFileExists{../../Logos/signature_avant_propos.png}{\includegraphics[height=1.6cm]{../../Logos/signature_avant_propos.png}}{\textit{MINKA MI NGUIDJOÏ Thierry Emmanuel}}

    \cleardoublepage

    % ==============================================================================
% Fichier : front/02_mode_emploi.tex : SGSI (SIS498)
% ==============================================================================

\chapter*{Mode d'emploi}
\addcontentsline{toc}{chapter}{Mode d'emploi}
\markboth{Mode d'emploi}{Mode d'emploi}

\section*{À qui s'adresse ce cours}

Aux étudiants du Master Professionnel en Sécurité des Systèmes Informatiques
et des Données, et aux professionnels qui exercent ou se destinent aux
métiers de \textbf{RSSI}, de \textbf{DSI}, de \textbf{délégué à la protection
des données} (\sigle{DPO}), d'\textbf{auditeur de SMSI} et de
\textbf{responsable conformité}. Des notions de réseaux, de systèmes et de
droit numérique sont utiles.

\section*{Comment lire ce livre}

Le livre comprend un prologue et six parties. Chaque chapitre suit le même
plan~: une épigraphe, le niveau et les prérequis, les objectifs, une question
directrice, cinq sections, un encadré « à retenir », dix exercices gradués et
une transition. Les niveaux des exercices sont signalés par des étoiles~:
\niveaubase{} fondamental, \niveaumoyen{} intermédiaire, \niveauexpert{}
expert.

\section*{Les encadrés}

\begin{itemize}
    \item \textbf{Scénario opérationnel}~: un épisode des opérations NYONG,
          KELLÉ, MAKAK ou MATOMB, avec \emph{constatation} (faits) et
          \emph{analyse} (interprétation, niveau de confiance).
    \item \textbf{Actif}~: la fiche d'un actif informationnel (valeur, support,
          critères de sécurité prioritaires).
    \item \textbf{Mesure}~: une mesure de sécurité, son objectif et sa mise en
          \oe uvre.
    \item \textbf{Application métier}~: ce que le chapitre change pour un
          métier.
    \item \textbf{Cadre juridique}~: le texte de droit applicable.
\end{itemize}

\section*{Constatation et analyse}

Le cours distingue systématiquement ce que l'on \textbf{constate} (un fait
vérifiable, une preuve) de ce que l'on en \textbf{déduit} (une interprétation,
assortie d'un niveau de confiance). C'est la discipline première du management
de la sécurité et de l'audit.

\section*{NDAP SUKULU et évaluation}

L'environnement compagnon NDAP SUKULU (page suivante) propose, pour chaque
chapitre, une fiche d'auto-évaluation et, pour les chapitres pratiques, un
lab sur pièces. Fiches et labs se font hors des 30~heures de cours. Le code
du cours est \textbf{\codeUE}.

    \cleardoublepage

    % ==============================================================================
% Fichier : front/03_ndap_sukulu.tex : SAFI (SIS589)
% ==============================================================================

\chapter*{NDAP SUKULU, l'environnement compagnon}
\addcontentsline{toc}{chapter}{NDAP SUKULU, l'environnement compagnon}
\markboth{NDAP SUKULU}{NDAP SUKULU}

\begin{center}
    \includegraphics[width=0.52\textwidth]{figures/logos/logo_ndap_sukulu.png}
\end{center}

\medskip
\textit{Ndap sukulu} signifie « la salle de classe » en bassa, et
\textit{Malet Yon}, « ton enseignant ». \textbf{NDAP SUKULU} est
l'environnement compagnon du cours~: pour chaque chapitre, une fiche
d'auto-évaluation de 40~questions, corrigée automatiquement~; pour les
chapitres pratiques, un lab sur pièces (inventaires d'actifs, registres de
risques, extraits de politique, tableaux de bord) dans un cadre entièrement
fermé~; et un tableau de bord de progression.

\medskip
On y accède par le code de cours \textbf{\codeUE}. Les fiches et les labs se
font hors des 30~heures de cours. Les pièces des labs sont fictives et
établies pour l'enseignement.

% ==============================================================================
% Fichier : front/03b_labs.tex : SGSI (IS498)
% ==============================================================================

\section*{Les dix labs de NDAP SUKULU}

\begin{tableaucours}{L{2.4cm}L{8.4cm}L{1.1cm}}{%
    \tableauHeader{Lab} &
    \tableauHeader{Sujet} &
    \tableauHeader{Ch.}
}
SGSI-LAB-01 & Inventaire et classification des actifs & 1 \\
SGSI-LAB-02 & Rédiger une politique de sécurité & 5 \\
SGSI-LAB-03 & Analyse d'écart ISO 27001 (gap analysis) & 8 \\
SGSI-LAB-04 & Appréciation des risques (EBIOS Risk Manager) & 12 \\
SGSI-LAB-05 & Matrice et traitement des risques & 13 \\
SGSI-LAB-06 & Construire une déclaration d'applicabilité & 14 \\
SGSI-LAB-07 & Contrôle d'accès et moindre privilège (RBAC) & 16 \\
SGSI-LAB-08 & Plan de continuité et sauvegarde 3-2-1 & 17 \\
SGSI-LAB-09 & Tableau de bord d'indicateurs de sécurité & 18 \\
SGSI-LAB-10 & Audit interne et non-conformités & 19 \\
\end{tableaucours}

\medskip
Chaque lab dure environ une heure et se corrige automatiquement. Les pièces
se trouvent dans le dossier \texttt{preuves} (lecture seule) et une copie de
travail dans \texttt{travail}.


    \cleardoublepage

    \tableofcontents
    \cleardoublepage

    \listoffigures
    \listoftables
    \cleardoublepage

% --------------------------------------------------------------------------
% CORPS DU COURS (pages arabes)
% Numérotation des chapitres CONTINUE sur tout le cours : le prologue est le
% chapitre 0, puis C01 à C21, un numéro par fichier.
% --------------------------------------------------------------------------
\mainmatter

    \part*{Prologue Opérationnel}
    \addcontentsline{toc}{part}{Prologue Opérationnel}
        \setcounter{chapter}{-1}
        \IfFileExists{partie0/P00_decision_mairie.tex}{% ============================================================================
% Fichier  : partie0/P00_decision_mairie.tex
% Titre    : Prologue : la Décision de la Mairie
% Auteur   : MINKA MI NGUIDJOÏ Thierry Emmanuel
% ============================================================================
% Fil rouge : Mairie d'Eséka (Nyong-et-Kéllé). Cas fictif ; mécanismes et
% références normatives exacts.
% ============================================================================

\chapter{Prologue : la Décision de la Mairie}
\label{chap:prologue}
\chaptermark{Prologue}

\epigraph{%
    « La sécurité n'est pas un produit que l'on achète,\\
    c'est une décision que l'on tient dans la durée. »%
}{Principe du cours}

\niveauChapitre{Introductif}{Aucun prérequis. Les notions sont reprises en
détail dans les chapitres.}

\begin{boxobjectifs}
\begin{itemize}
    \item Comprendre qu'un incident de sécurité est d'abord un problème de
          management, non de seule technique.
    \item Distinguer la réaction ponctuelle d'une démarche durable et pilotée.
    \item Découvrir le fil rouge de la mairie d'Eséka et ses quatre phases.
    \item Situer le cycle \sigle{PDCA} au cœur du cours.
    \item Formuler les questions auxquelles le cours répond.
\end{itemize}
\end{boxobjectifs}

\bigskip

% ============================================================================
\section*{13 janvier 2026, 7 h 50, mairie d'Eséka}
% ============================================================================

L'agent d'état civil arrive avant l'ouverture au public. Il branche, comme
chaque matin, la clé \texttt{USB} sur laquelle il a travaillé la veille chez
lui. En quelques minutes, les fichiers du poste deviennent illisibles~: les
actes de naissance numérisés, le registre des mariages, le tableur des
redevances ne s'ouvrent plus. À 8~h~15, trois postes du service sont touchés,
reliés au même dossier partagé. À 9~h, une file d'usagers attend des copies
d'actes que la mairie ne peut plus délivrer.

\medskip
Ce que personne ne sait encore ce matin-là\footnote{Cas fictif établi pour le
cours. Les noms (mairie d'Eséka, cabinet \textsc{Matomb Conseil}, agents
cités), les dates et les montants sont fictifs~; les mécanismes de gestion et
les références normatives (ISO/IEC 27001, 27002, 27005, lois camerounaises)
sont exacts.}~: le mal n'est pas la clé \texttt{USB}. Le mal, c'est qu'un seul
poste suffise à paralyser un service, qu'aucune sauvegarde récente n'existe,
que le même mot de passe serve à tout le monde, et que personne, dans la
mairie, n'ait jamais décidé qui est responsable de quoi en matière de
sécurité.

% ============================================================================
\section{Réparer ne suffit pas}
\label{sec:p00-reparer}
% ============================================================================

Le \sigle{DSI} de la commune, M.~\textsc{Biloa}, rétablit les postes en deux
jours à partir d'une vieille copie. Le service rouvre. L'affaire pourrait
s'arrêter là~: on a réparé. Mais le maire, M.~\textsc{Essama}, pose la bonne
question en conseil~: « Qu'est-ce qui nous garantit que cela ne
recommencera pas la semaine prochaine, en pire~? »

\medskip
La réponse est désagréable~: rien. Réparer traite l'incident du jour~; cela ne
traite pas la \termecle{cause}. La cause est organisationnelle~: pas de
politique, pas de rôles, pas d'appréciation des risques, pas de règles de
sauvegarde, pas de contrôle. Tant que ces manques demeurent, le prochain
incident est une question de temps. La mairie n'a pas un problème technique~:
elle a un problème de \termecle{management de la sécurité}.

\begin{boxmetier}{RSSI}
\textbf{Du pompier au pilote.} Réagir à chaque incident, c'est être pompier.
Un \sigle{RSSI} construit un système qui réduit le nombre d'incendies et
limite leurs dégâts. Il ne promet pas le risque zéro~: il promet un risque
\emph{connu}, \emph{décidé} et \emph{maîtrisé}.
\end{boxmetier}

% ============================================================================
\section{La décision : un système, pas un achat}
\label{sec:p00-decision}
% ============================================================================

Le maire ne débloque pas un budget pour « acheter de la sécurité ». Il prend
une décision de gouvernance~: doter la mairie d'un \termecle{système de
management de la sécurité de l'information} (\sigle{SMSI}), c'est-à-dire un
ensemble de règles, de rôles et d'activités qui, ensemble, maintiennent la
sécurité dans la durée. Il nomme Mme~\textsc{Abena} \sigle{RSSI} de la commune
et mandate le cabinet \textsc{Matomb Conseil} pour l'accompagner vers une
certification \norme{ISO/IEC 27001}.

\begin{boxdef}
\termecle{Système de management de la sécurité de l'information} (\sigle{SMSI})~:
partie du système de management global qui, sur la base d'une appréciation des
risques, établit, met en \oe uvre, exploite, surveille et améliore la sécurité
de l'information~\cite{iso27001}. Ce n'est pas un logiciel~: c'est une manière
organisée de décider et de tenir ses décisions.
\end{boxdef}

% ============================================================================
\section{Le fil rouge : quatre phases}
\label{sec:p00-filrouge}
% ============================================================================

Le cours suit la mairie d'Eséka, dans le département du Nyong-et-Kéllé, tout au
long de sa démarche. Quatre phases, nommées d'après la géographie locale,
structurent le fil rouge et reviennent dans les chapitres.

\begin{tableaucours}{L{3.2cm}L{9.0cm}}{%
    \tableauHeader{Phase} &
    \tableauHeader{Ce qu'elle installe}
}
Opération \textsc{Nyong} & La gouvernance et la politique de sécurité (parties I et II) \\
Opération \textsc{Kellé} & L'appréciation et le traitement des risques (partie IV) \\
Opération \textsc{Makak} & La mise en \oe uvre des mesures (partie V) \\
Opération \textsc{Matomb} & L'audit interne et la certification (partie VI) \\
\end{tableaucours}
\vspace{-0.8em}
\begin{center}{\small\itshape Tableau~0.1. Les quatre phases du fil rouge}\end{center}

% ============================================================================
\section{Le cœur du cours : le cycle PDCA}
\label{sec:p00-pdca}
% ============================================================================

Un \sigle{SMSI} n'est pas un projet qui se termine~: c'est une \termecle{roue}
qui tourne. Le cycle \sigle{PDCA} (\emph{Plan, Do, Check, Act}), ou roue de
Deming, en est le moteur et revient à chaque partie du cours.

\begin{tableaucours}{L{2.6cm}L{9.6cm}}{%
    \tableauHeader{Étape} &
    \tableauHeader{Ce que fait la mairie}
}
Planifier (\emph{Plan}) & Apprécier les risques, fixer la politique et les objectifs \\
Déployer (\emph{Do}) & Mettre en \oe uvre les mesures et former les agents \\
Contrôler (\emph{Check}) & Mesurer, auditer, passer en revue de direction \\
Agir (\emph{Act}) & Corriger, améliorer, décider le cycle suivant \\
\end{tableaucours}
\vspace{-0.8em}
\begin{center}{\small\itshape Tableau~0.2. Le cycle PDCA appliqué à la mairie}\end{center}

\begin{boxinfo}
\textbf{Les questions auxquelles ce cours répond.}
\begin{itemize}
    \item Qui décide de la sécurité, et au nom de quoi~?
    \item Comment connaître ses risques avant qu'ils ne deviennent des
          incidents~?
    \item Quelles mesures mettre en \oe uvre, et comment prouver qu'elles
          fonctionnent~?
    \item Comment faire reconnaître, par un audit, que la sécurité est
          maîtrisée~?
\end{itemize}
\end{boxinfo}

\separateur

\begin{boxresume}
\textbf{Ce qu'il faut retenir de ce prologue~:}
\begin{itemize}
    \item Un incident est d'abord un révélateur de \textbf{causes
          organisationnelles}, pas seulement un problème technique.
    \item \textbf{Réparer} traite le symptôme~; \textbf{manager} traite la
          cause.
    \item Un \textbf{SMSI} est un système durable de décisions tenues, non un
          achat ponctuel.
    \item Le fil rouge suit la \textbf{mairie d'Eséka} en quatre phases.
    \item Le \textbf{cycle PDCA} est le moteur du cours.
\end{itemize}
\end{boxresume}

\begin{boxinfo}
\textbf{Transition.} Avant de manager la sécurité, il faut la définir~: le
Chapitre~\ref{chap:cia} pose ce qu'est la sécurité de l'information.
\end{boxinfo}
}{}

    \part{Fondements du Management de la Sécurité}
    \label{part:fondements}

        \IfFileExists{partie1/P01_C01_securite_information.tex}{% ============================================================================
% Fichier  : partie1/P01_C01_securite_information.tex
% Titre    : Qu'est-ce que la Sécurité de l'Information
% Auteur   : MINKA MI NGUIDJOÏ Thierry Emmanuel
% ============================================================================
% Sources : ISO/IEC 27000:2018 ; ISO/IEC 27001:2022. Mairie d'Eséka : cas fictif.
% ============================================================================

\chapter{Qu'est-ce que la Sécurité de l'Information}
\label{chap:cia}
\chaptermark{Sécurité de l'information}

\epigraph{%
    « Protéger l'information, ce n'est pas la cacher~:\\
    c'est la rendre sûre, exacte et disponible à qui de droit. »%
}{Principe du cours}

\niveauChapitre{Fondamental}{Le prologue. Aucune connaissance technique
préalable.}

\begin{boxobjectifs}
\begin{itemize}
    \item Définir la sécurité de l'information et la distinguer de la sécurité
          informatique.
    \item Énoncer la triade \sigle{DIC} et l'illustrer sur des actifs réels.
    \item Ajouter les critères d'authenticité, de non-répudiation et de
          traçabilité.
    \item Identifier l'information comme un actif à valeur, quel que soit son
          support.
    \item Relier chaque critère à une conséquence concrète pour la mairie.
\end{itemize}
\end{boxobjectifs}

\bigskip

% ============================================================================
\section*{Question directrice : que protège-t-on, au juste, quand on « sécurise » ?}
% ============================================================================

La clé \texttt{USB} a chiffré des fichiers~: l'information est devenue
indisponible. Mais si elle avait été copiée et revendue, le problème aurait
été autre. Ce chapitre nomme précisément ce que l'on protège et selon quels
critères.

\medskip
\textbf{Plan du chapitre.}
\begin{enumerate}
    \item L'information, un actif (section~\ref{sec:c01-actif}).
    \item La triade \sigle{DIC} (section~\ref{sec:c01-dic}).
    \item Les critères complémentaires (section~\ref{sec:c01-criteres}).
    \item Sécurité de l'information ou sécurité informatique (section~\ref{sec:c01-perimetre}).
    \item Lire un incident par ses critères (section~\ref{sec:c01-lecture}).
\end{enumerate}

\separateur

% ============================================================================
\section{L'Information, un Actif}
\label{sec:c01-actif}
% ============================================================================

\begin{boxdef}
\termecle{Information}~: donnée qui a un sens et une valeur pour
l'organisation, quel que soit son support (fichier, papier, parole, savoir
d'un agent). \termecle{Actif}~: tout élément ayant de la valeur et qu'il faut
donc protéger.
\end{boxdef}

On sécurise une information parce qu'elle a de la valeur~: les actes de
naissance font foi, le fichier des contribuables conditionne les recettes, la
messagerie engage la commune. La valeur n'est pas dans le disque dur~: elle est
dans l'information. Protéger le support sans penser l'information, c'est
protéger le coffre en oubliant ce qu'il contient.

\begin{boxobjet}{Le registre d'état civil numérisé}
\textbf{Constatation.} Le registre numérisé des naissances est consultable par
tous les agents du service, sans mot de passe individuel, et n'est copié nulle
part ailleurs.

\medskip
\textbf{Analyse.} C'est un actif de valeur majeure (il fait foi) mais sans
protection~: ni confidentialité (accès non maîtrisé), ni disponibilité (copie
unique). Sa valeur et sa vulnérabilité en font une cible prioritaire.
\end{boxobjet}

% ============================================================================
\section{La Triade DIC}
\label{sec:c01-dic}
% ============================================================================

\begin{boxdef}
La \termecle{triade \sigle{DIC}} définit trois critères de sécurité~:
\begin{itemize}
    \item \termecle{Disponibilité}~: l'information est accessible à qui de
          droit quand il en a besoin.
    \item \termecle{Intégrité}~: l'information n'est ni altérée ni détruite de
          façon non autorisée.
    \item \termecle{Confidentialité}~: l'information n'est accessible qu'aux
          personnes autorisées.
\end{itemize}
En anglais~: \emph{CIA} (Confidentiality, Integrity, Availability).
\end{boxdef}

Les trois critères se tiennent et s'arbitrent. Verrouiller un registre
(confidentialité) sans en garder de copie met en péril sa disponibilité. La
sécurité consiste à tenir les trois, à hauteur de la valeur de l'actif.

\begin{tableaucours}{L{3.0cm}L{4.6cm}L{4.6cm}}{%
    \tableauHeader{Critère} &
    \tableauHeader{Atteinte à Eséka} &
    \tableauHeader{Conséquence}
}
Disponibilité & Registre chiffré, illisible & Service à l'arrêt, usagers bloqués \\
Intégrité & Date de naissance modifiée & Acte faux, responsabilité engagée \\
Confidentialité & Fichier contribuables copié & Vie privée exposée, loi 2024/017 \\
\end{tableaucours}
\vspace{-0.8em}
\begin{center}{\small\itshape Tableau~1.1. La triade DIC sur les actifs de la mairie}\end{center}

% ============================================================================
\section{Les Critères Complémentaires}
\label{sec:c01-criteres}
% ============================================================================

À la triade s'ajoutent, selon les besoins, trois critères~:

\begin{itemize}
    \item \termecle{Authenticité}~: l'origine de l'information est vérifiée
          (l'acte vient bien de la mairie).
    \item \termecle{Non-répudiation}~: on ne peut nier avoir effectué une
          action (l'agent qui a validé un acte ne peut le contester).
    \item \termecle{Traçabilité}~: les actions sont enregistrées et
          rattribuables.
\end{itemize}

Ces critères soutiennent la confiance~: un acte authentique, non répudiable et
tracé a une valeur juridique qu'un simple fichier n'a pas.

% ============================================================================
\section{Sécurité de l'Information ou Sécurité Informatique}
\label{sec:c01-perimetre}
% ============================================================================

\begin{boxdef}
La \termecle{sécurité de l'information} protège l'information sous toutes ses
formes (numérique, papier, orale). La \termecle{sécurité informatique} ne
protège que les systèmes qui la traitent. La première englobe la seconde.
\end{boxdef}

La distinction n'est pas scolaire. À la mairie, un acte lu à voix haute au
guichet, un dossier papier oublié sur un bureau ou une confidence échappent à
toute mesure informatique mais relèvent de la sécurité de l'information. Un
\sigle{SMSI} couvre l'information, pas seulement les ordinateurs.

\begin{boxmetier}{RSSI}
\textbf{Penser actif, pas machine.} Le \sigle{RSSI} d'Eséka commence par
lister les informations de valeur et leur support, avant de parler pare-feu.
On protège ce qui a de la valeur, pas ce qui a une prise électrique.
\end{boxmetier}

% ============================================================================
\section{Lire un Incident par ses Critères}
\label{sec:c01-lecture}
% ============================================================================

Nommer les critères donne une grille de lecture. Tout incident s'analyse en
demandant quel critère a été atteint, sur quel actif, et avec quelle
conséquence.

\begin{boxprocedure}
\textbf{Constatation.} Le 13~janvier, les fichiers du service d'état civil
sont illisibles.

\medskip
\textbf{Analyse} (confiance élevée). Critère atteint~: \emph{disponibilité}.
Actif~: registres et tableurs du service. L'intégrité et la confidentialité
ne sont pas, à ce stade, mises en cause~: distinguer les trois évite de
conclure trop vite à un vol de données.
\end{boxprocedure}

\separateur

\begin{boxresume}
\textbf{Ce qu'il faut retenir de ce chapitre~:}
\begin{itemize}
    \item L'\textbf{information} est un actif à valeur~; on la protège pour
          cette valeur.
    \item La \textbf{triade DIC} (disponibilité, intégrité, confidentialité)
          définit ce que « sécuriser » veut dire.
    \item \textbf{Authenticité, non-répudiation, traçabilité} soutiennent la
          confiance.
    \item La \textbf{sécurité de l'information} dépasse la sécurité
          informatique.
    \item Tout incident se lit par le \textbf{critère atteint} sur un actif.
\end{itemize}
\end{boxresume}

\begin{boxinfo}
\textbf{Points de contrôle de l'élève.} Je sais~: définir information et
actif~; énoncer la triade \sigle{DIC}~; citer les trois critères
complémentaires~; distinguer sécurité de l'information et sécurité
informatique~; analyser un incident par ses critères.
\end{boxinfo}

\bigskip

\begin{boxexo}[title={\ding{45}\quad Exercice~1.1. Trois critères \niveaubase}]
Pour chacun des cas suivants, dites quel critère \sigle{DIC} est atteint~:
registre chiffré~; montant d'une taxe modifié~; liste d'électeurs copiée.
\end{boxexo}

\begin{boxexo}[title={\ding{45}\quad Exercice~1.2. Actif ou support \niveaubase}]
Parmi~: le disque dur, l'acte de naissance, l'ordinateur, le savoir de l'agent,
lesquels sont des actifs informationnels~? Justifiez.
\end{boxexo}

\begin{boxexo}[title={\ding{45}\quad Exercice~1.3. Information ou informatique \niveaubase}]
Donnez un exemple, à la mairie, d'atteinte à la sécurité de l'information qui
ne passe par aucun ordinateur.
\end{boxexo}

\begin{boxexo}[title={\ding{45}\quad Exercice~1.4. Les compléments \niveaumoyen}]
Expliquez, sur l'exemple d'un acte de naissance, l'authenticité, la
non-répudiation et la traçabilité.
\end{boxexo}

\begin{boxexo}[title={\ding{45}\quad Exercice~1.5. Arbitrage \niveaumoyen}]
Verrouiller le registre sans copie renforce un critère et en menace un autre.
Lesquels~? Comment tenir les deux~?
\end{boxexo}

\begin{boxexo}[title={\ding{45}\quad Exercice~1.6. Valeur d'un actif \niveaumoyen}]
Classez trois actifs de la mairie par valeur décroissante et justifiez le
classement.
\end{boxexo}

\begin{boxexo}[title={\ding{45}\quad Exercice~1.7. Grille de lecture \niveaumoyen}]
Appliquez la grille « critère / actif / conséquence » au vol du fichier des
contribuables.
\end{boxexo}

\begin{boxexo}[title={\ding{45}\quad Exercice~1.8. Confidentialité et loi \niveaumoyen}]
En quoi une atteinte à la confidentialité du fichier des contribuables
engage-t-elle la mairie au regard de la loi 2024/017~?
\end{boxexo}

\begin{boxexo}[title={\ding{45}\quad Exercice~1.9. Cartographie \niveauexpert}]
Dressez un tableau des cinq principaux actifs informationnels de la mairie
d'Eséka, avec leur support et les critères prioritaires pour chacun.
\end{boxexo}

\begin{boxexo}[title={\ding{45}\quad Exercice~1.10. Note de cadrage \niveauexpert}]
Rédigez une note d'une page au maire définissant « sécurité de l'information »
pour la mairie et justifiant pourquoi elle dépasse l'achat d'un antivirus.
\end{boxexo}

\bigskip

\begin{boxinfo}
\textbf{Pour aller plus loin.}
\begin{itemize}
    \item ISO/IEC 27000~\cite{iso27000}~; ISO/IEC 27001~\cite{iso27001}.
\end{itemize}
\end{boxinfo}

\begin{boxinfo}
\textbf{Transition.} Savoir ce que l'on protège ne dit pas encore pourquoi ni
pour qui~: le Chapitre~\ref{chap:enjeux} relie le risque à la décision.
\end{boxinfo}
}{}
        \IfFileExists{partie1/P01_C02_risque_decision.tex}{% ============================================================================
% Fichier  : partie1/P01_C02_risque_decision.tex
% Titre    : Du Risque à la Décision
% Auteur   : MINKA MI NGUIDJOÏ Thierry Emmanuel
% ============================================================================
% Sources : ISO/IEC 27001:2022 ; ISO 31000:2018. Mairie d'Eséka : cas fictif.
% ============================================================================

\chapter{Du Risque à la Décision}
\label{chap:enjeux}
\chaptermark{Du risque à la décision}

\epigraph{%
    « Sécuriser, c'est décider~: quel risque accepter, lequel réduire,\\
    et avec quelles ressources. »%
}{Principe du cours}

\niveauChapitre{Fondamental}{Le Chapitre~\ref{chap:cia}.}

\begin{boxobjectifs}
\begin{itemize}
    \item Relier la sécurité aux enjeux économiques, réglementaires et
          stratégiques de l'organisation.
    \item Identifier les acteurs de la sécurité et leurs responsabilités.
    \item Définir le patrimoine informationnel à protéger.
    \item Comprendre que la sécurité est une décision fondée sur le risque.
    \item Situer la mairie d'Eséka dans son écosystème (ANTIC, parties
          prenantes).
\end{itemize}
\end{boxobjectifs}

\bigskip

% ============================================================================
\section*{Question directrice : au nom de quoi la mairie investit-elle dans la sécurité ?}
% ============================================================================

Dépenser pour la sécurité, c'est renoncer à dépenser ailleurs. Ce chapitre
montre que cette dépense se justifie par des enjeux, s'organise autour
d'acteurs et se décide en fonction du risque.

\medskip
\textbf{Plan du chapitre.}
\begin{enumerate}
    \item Les enjeux (section~\ref{sec:c02-enjeux}).
    \item Les acteurs (section~\ref{sec:c02-acteurs}).
    \item Le patrimoine informationnel (section~\ref{sec:c02-patrimoine}).
    \item La sécurité, une décision (section~\ref{sec:c02-decision}).
    \item L'écosystème de la mairie (section~\ref{sec:c02-ecosysteme}).
\end{enumerate}

\separateur

% ============================================================================
\section{Les Enjeux}
\label{sec:c02-enjeux}
% ============================================================================

\begin{tableaucours}{L{3.4cm}L{9.0cm}}{%
    \tableauHeader{Enjeu} &
    \tableauHeader{Pour la mairie d'Eséka}
}
Économique & Service interrompu, recettes de taxes retardées, remédiation coûteuse \\
Réglementaire & Loi 2024/017 sur les données personnelles, obligations de l'ANTIC \\
Stratégique & Confiance des usagers, image de la commune, continuité du service public \\
\end{tableaucours}
\vspace{-0.8em}
\begin{center}{\small\itshape Tableau~2.1. Les enjeux de la sécurité pour la mairie}\end{center}

Pour une collectivité, l'enjeu premier n'est pas le chiffre d'affaires mais la
\termecle{continuité du service public} et la confiance des administrés. Un
état civil à l'arrêt, c'est un citoyen qui ne peut ni se marier, ni inscrire
un enfant, ni prouver son identité.

% ============================================================================
\section{Les Acteurs}
\label{sec:c02-acteurs}
% ============================================================================

\begin{tableaucours}{L{3.6cm}L{8.8cm}}{%
    \tableauHeader{Acteur} &
    \tableauHeader{Responsabilité}
}
Le maire (direction) & Impulse, valide la politique, alloue les ressources, assume \\
RSSI & Pilote la sécurité, apprécie les risques, propose les mesures \\
DSI & Met en \oe uvre les solutions techniques \\
DPO / conseil & Veille à la conformité des traitements de données personnelles \\
Agents (utilisateurs) & Respectent les règles, signalent les incidents \\
ANTIC & Autorité nationale : appui, contrôle, réponse aux incidents \\
\end{tableaucours}
\vspace{-0.8em}
\begin{center}{\small\itshape Tableau~2.2. Les acteurs et leurs responsabilités}\end{center}

La sécurité n'est l'affaire de personne si elle est l'affaire de tous sans
responsabilité nommée. La première décision du maire, nommer une \sigle{RSSI},
crée le point où la responsabilité se concentre.

\begin{boxmetier}{Direction générale}
\textbf{Le rôle que nul ne peut déléguer.} Le maire peut déléguer la mise en
\oe uvre, jamais la responsabilité finale. C'est lui qui arbitre entre le
risque et le coût, parce que lui seul engage la commune.
\end{boxmetier}

% ============================================================================
\section{Le Patrimoine Informationnel}
\label{sec:c02-patrimoine}
% ============================================================================

\begin{boxdef}
\termecle{Patrimoine informationnel}~: ensemble des informations de valeur
d'une organisation et des moyens qui les portent. Le protéger suppose d'abord
de le connaître.
\end{boxdef}

\begin{boxobjet}{Le patrimoine de la mairie}
\textbf{Constatation.} Avant la démarche, la mairie n'a aucun inventaire de
ses informations sensibles.

\medskip
\textbf{Analyse.} On ne protège bien que ce que l'on a recensé. L'opération
\textsc{Nyong} commencera par un inventaire~: état civil, fichier des
contribuables, délibérations du conseil, messagerie, site web. Sans cet
inventaire, l'appréciation des risques n'a pas d'objet.
\end{boxobjet}

% ============================================================================
\section{La Sécurité, une Décision}
\label{sec:c02-decision}
% ============================================================================

La sécurité absolue n'existe pas et coûterait trop cher. Sécuriser, c'est
donc \termecle{décider}~: pour chaque risque, choisir de le réduire, de le
transférer, de l'éviter ou de l'accepter (Chapitre~\ref{chap:traitement}).
Cette décision appartient à la direction, sur proposition du \sigle{RSSI} et au
vu des risques appréciés.

\begin{boxinfo}
\textbf{Le fondement du SMSI.} Tout le cours découle de cette idée~: la
sécurité se \emph{décide} sur la base du \emph{risque}. C'est pourquoi
l'appréciation des risques (partie IV) est le cœur de la norme
\norme{ISO/IEC 27001}.
\end{boxinfo}

% ============================================================================
\section{L'Écosystème de la Mairie}
\label{sec:c02-ecosysteme}
% ============================================================================

La mairie n'est pas seule. Elle dépend de prestataires (hébergeur du site,
éditeur du logiciel d'état civil), relève d'autorités (ANTIC, tutelle) et sert
des usagers. Ces \termecle{parties intéressées} ont des attentes de sécurité
qui s'imposent au \sigle{SMSI} et seront recensées à la clause~4 de la norme
(Chapitre~\ref{chap:clauses46}).

\begin{boxjuris}
La loi camerounaise n°~2024/017 du 23~décembre~2024 sur la protection des
données à caractère personnel impose au responsable de traitement (ici, la
mairie) de garantir la sécurité des données qu'il détient. La sécurité n'est
donc pas seulement prudente~: elle est une obligation légale.
\end{boxjuris}

\separateur

\begin{boxresume}
\textbf{Ce qu'il faut retenir de ce chapitre~:}
\begin{itemize}
    \item La sécurité répond à des \textbf{enjeux} économiques,
          réglementaires et stratégiques~; pour une mairie, la continuité du
          service public prime.
    \item Chaque \textbf{acteur} a une responsabilité nommée~; la direction
          assume, le \sigle{RSSI} pilote.
    \item On protège un \textbf{patrimoine informationnel} qu'il faut d'abord
          inventorier.
    \item Sécuriser est une \textbf{décision} fondée sur le risque.
    \item Les \textbf{parties intéressées} imposent des attentes au
          \sigle{SMSI}.
\end{itemize}
\end{boxresume}

\begin{boxinfo}
\textbf{Points de contrôle de l'élève.} Je sais~: citer les trois types
d'enjeux~; nommer les acteurs et leurs rôles~; définir le patrimoine
informationnel~; expliquer pourquoi la sécurité est une décision~; citer une
obligation légale.
\end{boxinfo}

\bigskip

\begin{boxexo}[title={\ding{45}\quad Exercice~2.1. Trois enjeux \niveaubase}]
Donnez, pour la mairie, un exemple concret d'enjeu économique, réglementaire
et stratégique.
\end{boxexo}

\begin{boxexo}[title={\ding{45}\quad Exercice~2.2. Qui fait quoi \niveaubase}]
Reliez chaque acteur (maire, RSSI, DSI, agent) à sa responsabilité principale.
\end{boxexo}

\begin{boxexo}[title={\ding{45}\quad Exercice~2.3. Inventaire \niveaubase}]
Listez cinq informations de valeur que détient une mairie.
\end{boxexo}

\begin{boxexo}[title={\ding{45}\quad Exercice~2.4. Responsabilité finale \niveaumoyen}]
Pourquoi le maire ne peut-il pas déléguer la responsabilité finale de la
sécurité~?
\end{boxexo}

\begin{boxexo}[title={\ding{45}\quad Exercice~2.5. Décider \niveaumoyen}]
« La mairie n'a qu'à tout sécuriser au maximum. » Montrez les limites de cette
phrase.
\end{boxexo}

\begin{boxexo}[title={\ding{45}\quad Exercice~2.6. Parties intéressées \niveaumoyen}]
Identifiez quatre parties intéressées de la mairie et une attente de sécurité
pour chacune.
\end{boxexo}

\begin{boxexo}[title={\ding{45}\quad Exercice~2.7. Continuité publique \niveaumoyen}]
Pourquoi la continuité du service public pèse-t-elle plus, pour une mairie,
que pour une entreprise commerciale~?
\end{boxexo}

\begin{boxexo}[title={\ding{45}\quad Exercice~2.8. Obligation légale \niveaumoyen}]
En quoi la loi 2024/017 transforme-t-elle la sécurité des données de la mairie
en obligation~?
\end{boxexo}

\begin{boxexo}[title={\ding{45}\quad Exercice~2.9. Cartographie des acteurs \niveauexpert}]
Dressez l'organigramme de sécurité que la mairie d'Eséka devrait adopter, avec
les responsabilités de chacun.
\end{boxexo}

\begin{boxexo}[title={\ding{45}\quad Exercice~2.10. Argumentaire budgétaire \niveauexpert}]
Rédigez l'argumentaire par lequel la RSSI justifie, devant le conseil, un
budget sécurité, en s'appuyant sur les enjeux et le risque.
\end{boxexo}

\bigskip

\begin{boxinfo}
\textbf{Pour aller plus loin.}
\begin{itemize}
    \item ISO/IEC 27001~\cite{iso27001}~; ISO 31000~\cite{iso31000}.
\end{itemize}
\end{boxinfo}

\begin{boxinfo}
\textbf{Transition.} Décider une fois ne suffit pas~: il faut un cycle qui
entretienne la décision. Le Chapitre~\ref{chap:pdca} présente le \sigle{PDCA}.
\end{boxinfo}
}{}
        \IfFileExists{partie1/P01_C03_pdca.tex}{% ============================================================================
% Fichier  : partie1/P01_C03_pdca.tex
% Titre    : Le Cycle PDCA et l'Amélioration Continue
% Auteur   : MINKA MI NGUIDJOÏ Thierry Emmanuel
% ============================================================================
% Sources : ISO/IEC 27001:2022 (approche processus). Mairie d'Eséka : cas fictif.
% ============================================================================

\chapter{Le Cycle PDCA et l'Amélioration Continue}
\label{chap:pdca}
\chaptermark{Le cycle PDCA}

\epigraph{%
    « Un système de sécurité qui ne tourne pas se dégrade~:\\
    ce qui était sûr hier ne l'est plus demain. »%
}{Principe du cours}

\niveauChapitre{Fondamental}{Les Chapitres~\ref{chap:cia} et~\ref{chap:enjeux}.}

\begin{boxobjectifs}
\begin{itemize}
    \item Énoncer les quatre étapes du cycle \sigle{PDCA}.
    \item Comprendre pourquoi la sécurité est un cycle et non un projet.
    \item Relier chaque étape aux parties du cours et aux phases du fil rouge.
    \item Distinguer correction, action corrective et amélioration.
    \item Situer l'amélioration continue comme exigence de la norme.
\end{itemize}
\end{boxobjectifs}

\bigskip

% ============================================================================
\section*{Question directrice : pourquoi la sécurité n'est-elle jamais « finie » ?}
% ============================================================================

La mairie pourrait croire qu'après avoir installé des mesures, le travail est
accompli. Ce chapitre montre que la sécurité se dégrade si on ne l'entretient
pas, et que le cycle \sigle{PDCA} est l'outil de cet entretien.

\medskip
\textbf{Plan du chapitre.}
\begin{enumerate}
    \item Les quatre étapes (section~\ref{sec:c03-etapes}).
    \item Un cycle, pas un projet (section~\ref{sec:c03-cycle}).
    \item PDCA et plan du cours (section~\ref{sec:c03-plan}).
    \item Corriger, améliorer (section~\ref{sec:c03-corriger}).
    \item L'amélioration continue, une exigence (section~\ref{sec:c03-exigence}).
\end{enumerate}

\separateur

% ============================================================================
\section{Les Quatre Étapes}
\label{sec:c03-etapes}
% ============================================================================

\begin{boxdef}
Le \termecle{cycle PDCA} (\emph{Plan, Do, Check, Act}), ou roue de Deming,
organise l'amélioration continue en quatre étapes répétées~:
\termecle{planifier}, \termecle{déployer}, \termecle{contrôler}, \termecle{agir}.
\end{boxdef}

\begin{tableaucours}{L{2.6cm}L{4.4cm}L{5.2cm}}{%
    \tableauHeader{Étape} &
    \tableauHeader{But} &
    \tableauHeader{À la mairie}
}
Planifier & Fixer objectifs et mesures & Apprécier les risques, écrire la politique \\
Déployer & Mettre en \oe uvre & Déployer les mesures, former les agents \\
Contrôler & Vérifier l'efficacité & Mesurer, auditer, revue de direction \\
Agir & Améliorer & Corriger les écarts, décider le cycle suivant \\
\end{tableaucours}
\vspace{-0.8em}
\begin{center}{\small\itshape Tableau~3.1. Le cycle PDCA, étape par étape}\end{center}

% ============================================================================
\section{Un Cycle, pas un Projet}
\label{sec:c03-cycle}
% ============================================================================

Un projet a une fin~; un cycle recommence. La menace évolue, les agents
changent, les systèmes vieillissent~: une mesure efficace en janvier peut ne
plus l'être en décembre. Le \sigle{PDCA} rattrape cette dérive en imposant de
revenir régulièrement sur ce qui a été décidé.

\begin{boxscenario}{Opération NYONG. Le piège du « c'est fait »}
\textbf{Constatation.} Six mois après avoir imposé des mots de passe
individuels, la mairie découvre que trois agents les ont de nouveau partagés.

\medskip
\textbf{Analyse} (confiance élevée). La mesure avait été déployée (\emph{Do})
mais jamais contrôlée (\emph{Check}). Sans l'étape de contrôle, une mesure se
relâche. Le cycle n'est pas une formalité~: c'est ce qui empêche le retour en
arrière.
\end{boxscenario}

% ============================================================================
\section{PDCA et Plan du Cours}
\label{sec:c03-plan}
% ============================================================================

Le cours épouse le cycle. Les parties II et III \emph{planifient} (gouvernance,
politique, cadre \norme{ISO/IEC 27001})~; la partie IV \emph{planifie} par le
risque~; la partie V \emph{déploie} les mesures~; la partie VI \emph{contrôle}
et \emph{agit} (audit, certification, incidents, amélioration).

\begin{boxmetier}{RSSI}
\textbf{Tenir la roue.} Le \sigle{RSSI} n'a pas pour mission de « finir » la
sécurité, mais de faire tourner la roue~: chaque année, replanifier au vu des
nouveaux risques et des écarts constatés. Son tableau de bord suit le cycle.
\end{boxmetier}

% ============================================================================
\section{Corriger, Améliorer}
\label{sec:c03-corriger}
% ============================================================================

\begin{boxdef}
\termecle{Correction}~: action immédiate qui traite un écart constaté (rétablir
un service). \termecle{Action corrective}~: action qui traite la \emph{cause}
de l'écart pour qu'il ne se reproduise pas. \termecle{Amélioration}~:
renforcement au-delà de la simple conformité.
\end{boxdef}

Rétablir les postes après la clé \texttt{USB} était une \emph{correction}.
Interdire les clés personnelles et imposer des sauvegardes est une \emph{action
corrective}. Automatiser la détection des clés non autorisées serait une
\emph{amélioration}. La norme attend les trois, dans cet ordre de maturité.

% ============================================================================
\section{L'Amélioration Continue, une Exigence}
\label{sec:c03-exigence}
% ============================================================================

L'amélioration continue n'est pas une bonne intention~: c'est la clause~10 de
la norme \norme{ISO/IEC 27001} (Chapitre~\ref{chap:clauses710}). Un
\sigle{SMSI} qui ne s'améliore pas n'est pas conforme. Le \sigle{PDCA} est donc
à la fois une philosophie et une exigence vérifiée en audit.

\begin{boxinfo}
\textbf{À retenir pour tout le cours.} Chaque fois qu'une partie s'ouvre,
demandez-vous à quelle étape du \sigle{PDCA} elle correspond. La structure du
cours \emph{est} le cycle.
\end{boxinfo}

\separateur

\begin{boxresume}
\textbf{Ce qu'il faut retenir de ce chapitre~:}
\begin{itemize}
    \item Le \textbf{PDCA} (planifier, déployer, contrôler, agir) organise
          l'amélioration continue.
    \item La sécurité est un \textbf{cycle}, pas un projet~: sans entretien,
          elle se dégrade.
    \item Le \textbf{cours} suit les étapes du cycle.
    \item \textbf{Correction}, \textbf{action corrective} et
          \textbf{amélioration} se distinguent par ce qu'elles traitent.
    \item L'\textbf{amélioration continue} est une exigence de la norme
          (clause~10).
\end{itemize}
\end{boxresume}

\begin{boxinfo}
\textbf{Points de contrôle de l'élève.} Je sais~: énoncer les quatre étapes~;
expliquer pourquoi la sécurité est un cycle~; relier le PDCA au plan du cours~;
distinguer correction, action corrective et amélioration~; rattacher
l'amélioration continue à la norme.
\end{boxinfo}

\bigskip

\begin{boxexo}[title={\ding{45}\quad Exercice~3.1. Les quatre étapes \niveaubase}]
Nommez les quatre étapes du \sigle{PDCA} et donnez une action de la mairie pour
chacune.
\end{boxexo}

\begin{boxexo}[title={\ding{45}\quad Exercice~3.2. Cycle ou projet \niveaubase}]
Expliquez en deux phrases pourquoi la sécurité n'est pas un projet qui se
termine.
\end{boxexo}

\begin{boxexo}[title={\ding{45}\quad Exercice~3.3. Trois mots \niveaubase}]
Distinguez correction, action corrective et amélioration par un exemple de
chacune.
\end{boxexo}

\begin{boxexo}[title={\ding{45}\quad Exercice~3.4. L'étape manquante \niveaumoyen}]
Dans le cas des mots de passe repartagés, quelle étape du cycle a manqué~?
Justifiez.
\end{boxexo}

\begin{boxexo}[title={\ding{45}\quad Exercice~3.5. Rattacher au cycle \niveaumoyen}]
À quelle étape du PDCA rattachez-vous~: un audit interne~; une formation~; une
appréciation des risques~; une revue de direction~?
\end{boxexo}

\begin{boxexo}[title={\ding{45}\quad Exercice~3.6. De la correction à la cause \niveaumoyen}]
Pour l'incident de la clé \texttt{USB}, proposez une correction, une action
corrective et une amélioration.
\end{boxexo}

\begin{boxexo}[title={\ding{45}\quad Exercice~3.7. Dérive \niveaumoyen}]
Donnez deux raisons pour lesquelles une mesure efficace peut cesser de l'être
avec le temps.
\end{boxexo}

\begin{boxexo}[title={\ding{45}\quad Exercice~3.8. Exigence \niveaumoyen}]
Pourquoi un SMSI qui ne s'améliore jamais n'est-il pas conforme à
\norme{ISO/IEC 27001}~?
\end{boxexo}

\begin{boxexo}[title={\ding{45}\quad Exercice~3.9. Planifier un tour \niveauexpert}]
Décrivez un cycle PDCA annuel complet pour la mairie d'Eséka, de la
planification à l'amélioration.
\end{boxexo}

\begin{boxexo}[title={\ding{45}\quad Exercice~3.10. Note au conseil \niveauexpert}]
Rédigez une note expliquant au conseil municipal pourquoi la sécurité
nécessite un budget récurrent, et non une dépense unique.
\end{boxexo}

\bigskip

\begin{boxinfo}
\textbf{Pour aller plus loin.}
\begin{itemize}
    \item ISO/IEC 27001, approche processus et clause~10~\cite{iso27001}.
\end{itemize}
\end{boxinfo}

\begin{boxinfo}
\textbf{Transition.} Faire tourner la roue suppose une organisation qui la
porte~: le Chapitre~\ref{chap:gouvernance} ouvre la gouvernance de la sécurité.
\end{boxinfo}
}{}

    \part{Gouvernance et Politique de Sécurité}
    \label{part:gouvernance}

        \IfFileExists{partie2/P02_C04_gouvernance.tex}{% ============================================================================
% Fichier  : partie2/P02_C04_gouvernance.tex
% Titre    : La Gouvernance de la Sécurité
% Auteur   : MINKA MI NGUIDJOÏ Thierry Emmanuel
% ============================================================================
% Sources : ISO/IEC 27001:2022 (clause 5) ; COBIT 2019. Mairie d'Eséka : fictif.
% ============================================================================

\chapter{La Gouvernance de la Sécurité}
\label{chap:gouvernance}
\chaptermark{La gouvernance}

\epigraph{%
    « Sans responsable nommé, la sécurité est l'affaire de personne. »%
}{Principe du cours}

\niveauChapitre{Fondamental}{La partie I.}

\begin{boxobjectifs}
\begin{itemize}
    \item Définir la gouvernance de la sécurité et la distinguer du management
          opérationnel.
    \item Préciser le rôle du \sigle{RSSI}, de la direction et du comité de
          sécurité.
    \item Comprendre le rattachement et l'indépendance du \sigle{RSSI}.
    \item Établir une matrice de responsabilités (\sigle{RACI}).
    \item Relier la gouvernance à l'engagement de la direction exigé par la
          norme.
\end{itemize}
\end{boxobjectifs}

\bigskip

% ============================================================================
\section*{Question directrice : qui décide, qui exécute, qui contrôle ?}
% ============================================================================

Nommer une \sigle{RSSI} ne suffit pas~: encore faut-il définir ses pouvoirs,
son rattachement et ses relations avec la direction et les services. Ce
chapitre organise la gouvernance de la sécurité à la mairie.

\medskip
\textbf{Plan du chapitre.}
\begin{enumerate}
    \item Gouverner et manager (section~\ref{sec:c04-gouverner}).
    \item Le RSSI (section~\ref{sec:c04-rssi}).
    \item La direction et le comité de sécurité (section~\ref{sec:c04-comite}).
    \item La matrice de responsabilités (section~\ref{sec:c04-raci}).
    \item L'engagement de la direction (section~\ref{sec:c04-engagement}).
\end{enumerate}

\separateur

% ============================================================================
\section{Gouverner et Manager}
\label{sec:c04-gouverner}
% ============================================================================

\begin{boxdef}
\termecle{Gouvernance de la sécurité}~: ensemble des responsabilités et
pratiques par lesquelles la direction oriente la sécurité, fixe les objectifs
et s'assure que les risques sont maîtrisés. Elle décide~; le management
exécute.
\end{boxdef}

La gouvernance répond à « quoi et pourquoi »~; le management opérationnel à
« comment ». Confondre les deux conduit soit à une direction qui bricole de la
technique, soit à une technique sans mandat. La norme sépare nettement le rôle
d'orientation (direction) du rôle de pilotage (\sigle{RSSI}).

% ============================================================================
\section{Le RSSI}
\label{sec:c04-rssi}
% ============================================================================

\begin{tableaucours}{L{3.0cm}L{9.2cm}}{%
    \tableauHeader{Mission} &
    \tableauHeader{Contenu}
}
Définir & Proposer la stratégie et la politique de sécurité \\
Piloter & Conduire l'appréciation des risques et le plan de traitement \\
Coordonner & Mettre en relation direction, DSI, agents, prestataires \\
Contrôler & Vérifier le respect des règles, suivre les indicateurs \\
Conseiller & Éclairer la direction pour ses arbitrages \\
\end{tableaucours}
\vspace{-0.8em}
\begin{center}{\small\itshape Tableau~4.1. Les missions du RSSI}\end{center}

\begin{boxprocedure}
\textbf{Constatation.} À la mairie, Mme~\textsc{Abena} (RSSI) est aussi chef du
service informatique, sous l'autorité du \sigle{DSI}.

\medskip
\textbf{Analyse} (confiance moyenne). Ce rattachement crée un conflit~: juger
la sécurité des systèmes que l'on exploite soi-même, c'est être juge et
partie. Il est préférable que la \sigle{RSSI} rende compte directement au
maire, pour garder son indépendance de contrôle.
\end{boxprocedure}

% ============================================================================
\section{La Direction et le Comité de Sécurité}
\label{sec:c04-comite}
% ============================================================================

La direction (le maire) impulse, valide la politique, alloue les ressources et
assume la responsabilité finale. Pour instruire ses décisions, elle s'appuie
sur un \termecle{comité de sécurité}~: instance périodique qui réunit maire,
\sigle{RSSI}, \sigle{DSI} et responsables de services, arbitre les risques et
suit le plan d'action.

\begin{boxmetier}{Direction générale}
\textbf{Présider, pas assister.} Un comité de sécurité que le maire délègue
systématiquement perd son autorité. La présence de la direction signale que la
sécurité est une priorité, non une affaire de techniciens.
\end{boxmetier}

% ============================================================================
\section{La Matrice de Responsabilités}
\label{sec:c04-raci}
% ============================================================================

\begin{boxdef}
Une matrice \termecle{RACI} attribue, pour chaque activité, qui est
\emph{Responsable} (fait), \emph{Approbateur} (décide), \emph{Consulté} et
\emph{Informé}. Elle supprime les zones où « on croyait que l'autre le
faisait ».
\end{boxdef}

\begin{tableaucours}{L{4.0cm}L{2.0cm}L{2.0cm}L{2.0cm}}{%
    \tableauHeader{Activité} &
    \tableauHeader{Maire} &
    \tableauHeader{RSSI} &
    \tableauHeader{DSI}
}
Valider la politique & A & R & C \\
Apprécier les risques & I & R & C \\
Déployer une mesure & I & A & R \\
Suivre les indicateurs & I & R & C \\
\end{tableaucours}
\vspace{-0.8em}
\begin{center}{\small\itshape Tableau~4.2. Extrait de la matrice RACI de la mairie}\end{center}

% ============================================================================
\section{L'Engagement de la Direction}
\label{sec:c04-engagement}
% ============================================================================

La gouvernance n'existe que si la direction s'engage réellement. La clause~5
de \norme{ISO/IEC 27001} l'exige~: la direction doit démontrer son
\termecle{leadership} en établissant la politique, en attribuant les rôles et
en fournissant les ressources (Chapitre~\ref{chap:clauses46}).

\begin{boxjuris}
En droit camerounais comme dans la norme, la responsabilité de la sécurité des
données remonte au responsable de traitement. Le maire ne peut se retrancher
derrière ses services~: la gouvernance formalise cette responsabilité, elle ne
la dilue pas.
\end{boxjuris}

\separateur

\begin{boxresume}
\textbf{Ce qu'il faut retenir de ce chapitre~:}
\begin{itemize}
    \item La \textbf{gouvernance} oriente et décide~; le management exécute.
    \item Le \textbf{RSSI} définit, pilote, coordonne, contrôle, conseille~;
          son indépendance suppose un bon rattachement.
    \item La \textbf{direction} impulse et assume~; le \textbf{comité de
          sécurité} instruit ses décisions.
    \item La \textbf{matrice RACI} supprime les responsabilités floues.
    \item La norme exige l'\textbf{engagement} de la direction (clause~5).
\end{itemize}
\end{boxresume}

\begin{boxinfo}
\textbf{Points de contrôle de l'élève.} Je sais~: distinguer gouverner et
manager~; énoncer les missions du RSSI~; expliquer le rôle du comité~; lire une
matrice RACI~; relier la gouvernance à la clause~5.
\end{boxinfo}

\bigskip

\begin{boxexo}[title={\ding{45}\quad Exercice~4.1. Gouverner ou manager \niveaubase}]
Classez~: fixer la politique~; configurer un pare-feu~; arbitrer un budget~;
appliquer un correctif. Gouvernance ou management~?
\end{boxexo}

\begin{boxexo}[title={\ding{45}\quad Exercice~4.2. Missions du RSSI \niveaubase}]
Citez cinq missions du \sigle{RSSI} et donnez un exemple pour l'une d'elles.
\end{boxexo}

\begin{boxexo}[title={\ding{45}\quad Exercice~4.3. Le comité \niveaubase}]
Qui siège au comité de sécurité de la mairie, et à quoi sert-il~?
\end{boxexo}

\begin{boxexo}[title={\ding{45}\quad Exercice~4.4. Juge et partie \niveaumoyen}]
Pourquoi rattacher la RSSI au DSI pose-t-il un problème d'indépendance~?
Proposez un meilleur rattachement.
\end{boxexo}

\begin{boxexo}[title={\ding{45}\quad Exercice~4.5. RACI \niveaumoyen}]
Complétez une ligne RACI pour l'activité « gérer les accès des agents ».
\end{boxexo}

\begin{boxexo}[title={\ding{45}\quad Exercice~4.6. Leadership \niveaumoyen}]
Donnez trois actes concrets par lesquels le maire démontre son engagement.
\end{boxexo}

\begin{boxexo}[title={\ding{45}\quad Exercice~4.7. Délégation \niveaumoyen}]
Que peut déléguer le maire, que ne peut-il pas déléguer~? Justifiez.
\end{boxexo}

\begin{boxexo}[title={\ding{45}\quad Exercice~4.8. Zone floue \niveaumoyen}]
Donnez un exemple de responsabilité floue à la mairie et montrez comment la
RACI la lève.
\end{boxexo}

\begin{boxexo}[title={\ding{45}\quad Exercice~4.9. Charte de comité \niveauexpert}]
Rédigez la charte du comité de sécurité de la mairie (composition, fréquence,
missions, décisions).
\end{boxexo}

\begin{boxexo}[title={\ding{45}\quad Exercice~4.10. Fiche de poste \niveauexpert}]
Rédigez la fiche de poste de la RSSI de la mairie d'Eséka, avec rattachement et
pouvoirs.
\end{boxexo}

\bigskip

\begin{boxinfo}
\textbf{Pour aller plus loin.}
\begin{itemize}
    \item ISO/IEC 27001, clause~5~\cite{iso27001}~; COBIT 2019~\cite{cobit2019}.
\end{itemize}
\end{boxinfo}

\begin{boxinfo}
\textbf{Transition.} La gouvernance s'exprime d'abord dans un document~: le
Chapitre~\ref{chap:politique} construit la politique de sécurité.
\end{boxinfo}
}{}
        \IfFileExists{partie2/P02_C05_politique.tex}{% ============================================================================
% Fichier  : partie2/P02_C05_politique.tex
% Titre    : La Politique de Sécurité du Système d'Information
% Auteur   : MINKA MI NGUIDJOÏ Thierry Emmanuel
% ============================================================================
% Sources : ISO/IEC 27001:2022 (clause 5.2) ; ISO/IEC 27002:2022. Fictif : Eséka.
% ============================================================================

\chapter{La Politique de Sécurité du Système d'Information}
\label{chap:politique}
\chaptermark{La politique de sécurité}

\epigraph{%
    « Une politique qui ne s'applique pas est pire que pas de politique~:\\
    elle donne l'illusion d'être protégé. »%
}{Principe du cours}

\niveauChapitre{Fondamental}{Le Chapitre~\ref{chap:gouvernance}.}

\begin{boxobjectifs}
\begin{itemize}
    \item Définir la politique de sécurité (\sigle{PSSI}) et ses objectifs.
    \item Décrire la structure type d'une politique de sécurité.
    \item Distinguer politique, procédure et charte.
    \item Comprendre le cycle de vie d'une politique (validation, diffusion,
          révision).
    \item Relier la politique à l'exigence de la clause~5.2 de la norme.
\end{itemize}
\end{boxobjectifs}

\bigskip

% ============================================================================
\section*{Question directrice : comment transformer des intentions en règles opposables ?}
% ============================================================================

Le maire veut « plus de sécurité ». Cette intention ne protège rien tant
qu'elle n'est pas écrite, validée, diffusée et appliquée. La politique de
sécurité est le document qui opère cette transformation.

\medskip
\textbf{Plan du chapitre.}
\begin{enumerate}
    \item Qu'est-ce qu'une politique (section~\ref{sec:c05-definir}).
    \item Objectifs et portée (section~\ref{sec:c05-objectifs}).
    \item La structure type (section~\ref{sec:c05-structure}).
    \item Politique, procédure, charte (section~\ref{sec:c05-niveaux}).
    \item Le cycle de vie du document (section~\ref{sec:c05-cycle}).
\end{enumerate}

\separateur

% ============================================================================
\section{Qu'est-ce qu'une Politique}
\label{sec:c05-definir}
% ============================================================================

\begin{boxdef}
\termecle{Politique de sécurité du système d'information} (\sigle{PSSI})~:
document formel, approuvé par la direction, qui fixe les objectifs de sécurité,
les règles à respecter et les responsabilités. Elle est le cadre de référence
de toutes les décisions de sécurité.
\end{boxdef}

La \sigle{PSSI} est au \sigle{SMSI} ce que la constitution est à un État~: le
texte fondateur dont découlent les règles particulières. Elle doit être courte,
claire, et tenue par la direction, sans quoi elle reste lettre morte.

% ============================================================================
\section{Objectifs et Portée}
\label{sec:c05-objectifs}
% ============================================================================

\begin{tableaucours}{L{4.0cm}L{8.2cm}}{%
    \tableauHeader{Objectif} &
    \tableauHeader{Ce qu'il apporte}
}
Établir un cadre & Une référence commune pour arbitrer \\
Sensibiliser & Des règles connues de tous les agents \\
Définir les responsabilités & Chacun sait ce qu'il doit faire \\
Assurer la conformité & Le respect des obligations légales \\
Faciliter la décision & Un critère pour trancher les cas particuliers \\
\end{tableaucours}
\vspace{-0.8em}
\begin{center}{\small\itshape Tableau~5.1. Les objectifs d'une politique de sécurité}\end{center}

La \termecle{portée} délimite ce que couvre la politique~: quels services,
quels sites, quelles informations. À Eséka, la politique couvre tous les
services de la mairie et l'ensemble du patrimoine informationnel recensé.

% ============================================================================
\section{La Structure Type}
\label{sec:c05-structure}
% ============================================================================

\begin{tableaucours}{L{4.2cm}L{8.0cm}}{%
    \tableauHeader{Section} &
    \tableauHeader{Contenu}
}
Introduction et contexte & Objectifs, portée, références réglementaires \\
Gouvernance & Organisation, rôles, comité de sécurité \\
Classification de l'information & Niveaux et règles de manipulation \\
Contrôle d'accès & Création des comptes, authentification, privilèges \\
Sécurité des opérations & Sauvegarde, journalisation, gestion des incidents \\
Conformité & Audits, revue, sanctions \\
\end{tableaucours}
\vspace{-0.8em}
\begin{center}{\small\itshape Tableau~5.2. Structure type d'une politique de sécurité}\end{center}

\begin{boxmetier}{RSSI}
\textbf{Écrire pour être lu.} La RSSI d'Eséka proscrit le jargon~: une
politique que les agents ne comprennent pas n'est pas appliquée. Chaque règle
dit qui, quoi, quand, et la conséquence du non-respect.
\end{boxmetier}

% ============================================================================
\section{Politique, Procédure, Charte}
\label{sec:c05-niveaux}
% ============================================================================

\begin{boxdef}
Trois niveaux se complètent~: la \termecle{politique} fixe les principes (le
« quoi »)~; la \termecle{procédure} décrit les étapes d'exécution (le
« comment »)~; la \termecle{charte} engage l'utilisateur sur ses droits et
devoirs.
\end{boxdef}

La politique dit « les accès sont nominatifs »~; la procédure décrit comment
créer et supprimer un compte~; la charte informatique, signée par l'agent,
l'engage à ne pas partager son mot de passe (Chapitre~\ref{chap:conformite}).

% ============================================================================
\section{Le Cycle de Vie du Document}
\label{sec:c05-cycle}
% ============================================================================

Une politique vit~: elle est \termecle{rédigée}, \termecle{validée} par la
direction, \termecle{diffusée} à tous, \termecle{appliquée}, puis
\termecle{révisée} périodiquement. Sans validation de la direction, elle n'a
pas d'autorité~; sans révision, elle se périme.

\begin{boxprocedure}
\textbf{Constatation.} La mairie a rédigé une politique en 2022, jamais signée
par le maire ni communiquée aux agents.

\medskip
\textbf{Analyse} (confiance élevée). Un document non validé et non diffusé
n'est pas une politique~: c'est un brouillon. L'opération \textsc{Nyong}
reprendra le cycle complet, de la validation à la diffusion contre signature.
\end{boxprocedure}

\begin{boxjuris}
La clause~5.2 de \norme{ISO/IEC 27001} exige une politique de sécurité
établie, approuvée par la direction, communiquée et disponible. L'absence de
politique formelle est un écart majeur en audit de certification
(Chapitre~\ref{chap:audit}).
\end{boxjuris}

\separateur

\begin{boxresume}
\textbf{Ce qu'il faut retenir de ce chapitre~:}
\begin{itemize}
    \item La \textbf{PSSI} est le texte fondateur du \sigle{SMSI}, approuvé par
          la direction.
    \item Elle poursuit cinq \textbf{objectifs} et définit une \textbf{portée}.
    \item Sa \textbf{structure type} couvre gouvernance, accès, opérations,
          conformité.
    \item \textbf{Politique, procédure, charte} se distinguent par leur niveau.
    \item Une politique a un \textbf{cycle de vie}~: validée, diffusée,
          révisée~; la clause~5.2 l'exige.
\end{itemize}
\end{boxresume}

\begin{boxinfo}
\textbf{Points de contrôle de l'élève.} Je sais~: définir la \sigle{PSSI}~;
citer ses objectifs~; décrire sa structure~; distinguer politique, procédure et
charte~; énoncer son cycle de vie.
\end{boxinfo}

\bigskip

\begin{boxexo}[title={\ding{45}\quad Exercice~5.1. Définition \niveaubase}]
Qu'est-ce qu'une politique de sécurité, et qui doit l'approuver~?
\end{boxexo}

\begin{boxexo}[title={\ding{45}\quad Exercice~5.2. Objectifs \niveaubase}]
Citez trois objectifs d'une \sigle{PSSI} et illustrez-en un pour la mairie.
\end{boxexo}

\begin{boxexo}[title={\ding{45}\quad Exercice~5.3. Trois niveaux \niveaubase}]
Classez~: « les accès sont nominatifs »~; « voici comment créer un compte »~;
« je m'engage à ne pas partager mon mot de passe ». Politique, procédure ou
charte~?
\end{boxexo}

\begin{boxexo}[title={\ding{45}\quad Exercice~5.4. Portée \niveaumoyen}]
Rédigez la phrase de portée de la politique de la mairie d'Eséka.
\end{boxexo}

\begin{boxexo}[title={\ding{45}\quad Exercice~5.5. Brouillon ou politique \niveaumoyen}]
Pourquoi le document de 2022 n'est-il pas une politique~? Que manque-t-il~?
\end{boxexo}

\begin{boxexo}[title={\ding{45}\quad Exercice~5.6. Une règle complète \niveaumoyen}]
Rédigez une règle de politique sur l'usage des clés \texttt{USB}, avec la
conséquence du non-respect.
\end{boxexo}

\begin{boxexo}[title={\ding{45}\quad Exercice~5.7. Cycle de vie \niveaumoyen}]
Énoncez les cinq étapes du cycle de vie d'une politique et le risque d'en
oublier une.
\end{boxexo}

\begin{boxexo}[title={\ding{45}\quad Exercice~5.8. Écart d'audit \niveaumoyen}]
Pourquoi l'absence de politique validée est-elle un écart majeur en audit~?
\end{boxexo}

\begin{boxexo}[title={\ding{45}\quad Exercice~5.9. Table des matières \niveauexpert}]
Dressez la table des matières complète de la politique de sécurité de la mairie
d'Eséka.
\end{boxexo}

\begin{boxexo}[title={\ding{45}\quad Exercice~5.10. Section rédigée \niveauexpert}]
Rédigez la section « Contrôle d'accès » de la politique de la mairie (une
page).
\end{boxexo}

\bigskip

\begin{boxinfo}
\textbf{Pour aller plus loin.}
\begin{itemize}
    \item ISO/IEC 27001, clause~5.2~\cite{iso27001}~; ISO/IEC 27002~\cite{iso27002}.
\end{itemize}
\end{boxinfo}

\begin{boxinfo}
\textbf{Transition.} La politique s'adosse à des obligations légales et à des
engagements individuels~: le Chapitre~\ref{chap:conformite} traite chartes,
conformité et données personnelles.
\end{boxinfo}
}{}
        \IfFileExists{partie2/P02_C06_chartes_rgpd.tex}{% ============================================================================
% Fichier  : partie2/P02_C06_chartes_rgpd.tex
% Titre    : Chartes, Conformité et Protection des Données
% Auteur   : MINKA MI NGUIDJOÏ Thierry Emmanuel
% ============================================================================
% Sources : Loi camerounaise 2024/017 ; Convention de Malabo 2014 ;
% ISO/IEC 27701. Mairie d'Eséka : cas fictif.
% ============================================================================

\chapter{Chartes, Conformité et Protection des Données}
\label{chap:conformite}
\chaptermark{Chartes et conformité}

\epigraph{%
    « Les données personnelles ne sont pas la propriété de la mairie~:\\
    elles lui sont confiées. »%
}{Principe du cours}

\niveauChapitre{Fondamental}{Le Chapitre~\ref{chap:politique}.}

\begin{boxobjectifs}
\begin{itemize}
    \item Rédiger et faire vivre une charte informatique.
    \item Énoncer les principes de protection des données personnelles.
    \item Situer la loi camerounaise 2024/017 et la Convention de Malabo.
    \item Définir le rôle du délégué à la protection des données (\sigle{DPO}).
    \item Distinguer sécurité et conformité, et montrer leur recouvrement.
\end{itemize}
\end{boxobjectifs}

\bigskip

% ============================================================================
\section*{Question directrice : la mairie a-t-elle le droit de faire ce qu'elle fait de ses données ?}
% ============================================================================

La mairie détient des données très sensibles~: naissances, mariages, décès,
contribuables. En détenir ne vaut pas en faire ce que l'on veut. Ce chapitre
pose le cadre de conformité et l'engagement des agents.

\medskip
\textbf{Plan du chapitre.}
\begin{enumerate}
    \item La charte informatique (section~\ref{sec:c06-charte}).
    \item Les principes de protection des données (section~\ref{sec:c06-principes}).
    \item Le cadre juridique camerounais (section~\ref{sec:c06-droit}).
    \item Le délégué à la protection des données (section~\ref{sec:c06-dpo}).
    \item Sécurité et conformité (section~\ref{sec:c06-recouvrement}).
\end{enumerate}

\separateur

% ============================================================================
\section{La Charte Informatique}
\label{sec:c06-charte}
% ============================================================================

\begin{boxdef}
\termecle{Charte informatique}~: document opérationnel qui précise les règles
d'utilisation du système d'information, les droits et devoirs des utilisateurs,
et les sanctions en cas de non-respect. Signée, elle engage l'agent.
\end{boxdef}

La charte traduit la politique en obligations individuelles~: usage des postes,
gestion des mots de passe, messagerie, supports amovibles, télétravail. Sa
force vient de la signature~: un agent qui a signé ne peut plaider
l'ignorance.

% ============================================================================
\section{Les Principes de Protection des Données}
\label{sec:c06-principes}
% ============================================================================

\begin{tableaucours}{L{3.6cm}L{8.6cm}}{%
    \tableauHeader{Principe} &
    \tableauHeader{Exigence}
}
Licéité et loyauté & Une base légale pour chaque traitement \\
Finalité & Données collectées pour un but déterminé \\
Minimisation & Données pertinentes et limitées au nécessaire \\
Exactitude & Données exactes et tenues à jour \\
Conservation limitée & Durée de conservation appropriée \\
Sécurité & Mesures techniques et organisationnelles \\
Responsabilité & Démontrer la conformité \\
\end{tableaucours}
\vspace{-0.8em}
\begin{center}{\small\itshape Tableau~6.1. Les principes de protection des données}\end{center}

\begin{boxobjet}{Le fichier des contribuables}
\textbf{Constatation.} Le fichier des contribuables contient des informations
qui ne servent plus (anciens redevables partis depuis dix ans).

\medskip
\textbf{Analyse.} Deux principes sont en défaut~: la \emph{minimisation} et la
\emph{conservation limitée}. Conserver sans fin expose inutilement des données
personnelles. La conformité impose une politique de purge.
\end{boxobjet}

% ============================================================================
\section{Le Cadre Juridique Camerounais}
\label{sec:c06-droit}
% ============================================================================

\begin{boxjuris}
La loi camerounaise n°~2024/017 du 23~décembre~2024 régit la protection des
données à caractère personnel~: elle impose une base légale, la sécurité des
données et des droits aux personnes concernées. La Convention de l'Union
africaine sur la cybersécurité et la protection des données (Convention de
\textbf{Malabo}, 2014) fixe le cadre continental. La mairie, responsable de
traitement, doit s'y conformer.
\end{boxjuris}

La conformité n'est pas optionnelle pour une administration~: elle conditionne
la légalité même du service rendu.

% ============================================================================
\section{Le Délégué à la Protection des Données}
\label{sec:c06-dpo}
% ============================================================================

\begin{boxdef}
\termecle{Délégué à la protection des données} (\sigle{DPO})~: personne chargée
de veiller à la conformité des traitements de données personnelles, de
conseiller l'organisation et de faire le lien avec l'autorité de protection.
\end{boxdef}

À Eséka, faute de \sigle{DPO} interne, la mairie confie ce rôle à
Me~\textsc{Essomba}, conseil extérieur. Le \sigle{DPO} tient le registre des
traitements, évalue les risques pour la vie privée et instruit les demandes des
usagers.

\begin{boxmetier}{Responsable conformité}
\textbf{Conformité n'est pas sécurité.} Le responsable conformité vérifie que
la mairie respecte ses obligations~; le \sigle{RSSI} réduit les risques. Les
deux se recoupent mais ne se confondent pas~: on peut être sécurisé et non
conforme (données gardées trop longtemps), ou conforme sur le papier et peu
sécurisé.
\end{boxmetier}

% ============================================================================
\section{Sécurité et Conformité}
\label{sec:c06-recouvrement}
% ============================================================================

La sécurité (\sigle{DIC}) et la conformité (respect des règles) se recouvrent
largement~: la loi exige la sécurité, et la sécurité sert la conformité. Mais
elles diffèrent de logique~: la sécurité part du \emph{risque}, la conformité
part de l'\emph{obligation}. Un bon \sigle{SMSI} tient les deux.

\begin{boxinfo}
\textbf{Un pont vers le reste du cours.} La conformité donne des exigences
d'entrée à l'appréciation des risques (partie IV)~: certaines mesures seront
obligatoires quel que soit le niveau de risque, parce que la loi les impose.
\end{boxinfo}

\separateur

\begin{boxresume}
\textbf{Ce qu'il faut retenir de ce chapitre~:}
\begin{itemize}
    \item La \textbf{charte informatique} engage l'agent par sa signature.
    \item Les \textbf{principes} de protection des données (finalité,
          minimisation, conservation limitée\ldots) encadrent les traitements.
    \item La \textbf{loi 2024/017} et la \textbf{Convention de Malabo} fixent
          le cadre camerounais et africain.
    \item Le \textbf{DPO} veille à la conformité des traitements.
    \item \textbf{Sécurité et conformité} se recoupent sans se confondre.
\end{itemize}
\end{boxresume}

\begin{boxinfo}
\textbf{Points de contrôle de l'élève.} Je sais~: rédiger une règle de charte~;
citer les principes de protection des données~; situer la loi 2024/017 et
Malabo~; définir le rôle du \sigle{DPO}~; distinguer sécurité et conformité.
\end{boxinfo}

\bigskip

\begin{boxexo}[title={\ding{45}\quad Exercice~6.1. La charte \niveaubase}]
Citez trois règles que doit contenir la charte informatique de la mairie.
\end{boxexo}

\begin{boxexo}[title={\ding{45}\quad Exercice~6.2. Les principes \niveaubase}]
Énoncez quatre principes de protection des données personnelles.
\end{boxexo}

\begin{boxexo}[title={\ding{45}\quad Exercice~6.3. La loi \niveaubase}]
Quelle loi camerounaise régit la protection des données personnelles, et de
quand date-t-elle~?
\end{boxexo}

\begin{boxexo}[title={\ding{45}\quad Exercice~6.4. Deux principes en défaut \niveaumoyen}]
Pour le fichier des contribuables conservé sans fin, quels principes sont
violés~? Que faire~?
\end{boxexo}

\begin{boxexo}[title={\ding{45}\quad Exercice~6.5. Rôle du DPO \niveaumoyen}]
Décrivez trois missions du \sigle{DPO} à la mairie d'Eséka.
\end{boxexo}

\begin{boxexo}[title={\ding{45}\quad Exercice~6.6. Sécurisé mais non conforme \niveaumoyen}]
Donnez un exemple d'une mairie bien sécurisée mais non conforme à la loi.
\end{boxexo}

\begin{boxexo}[title={\ding{45}\quad Exercice~6.7. Base légale \niveaumoyen}]
Quelle base légale justifie que la mairie traite les actes d'état civil~?
\end{boxexo}

\begin{boxexo}[title={\ding{45}\quad Exercice~6.8. Signature \niveaumoyen}]
Pourquoi la signature de la charte est-elle juridiquement importante~?
\end{boxexo}

\begin{boxexo}[title={\ding{45}\quad Exercice~6.9. Charte rédigée \niveauexpert}]
Rédigez une charte informatique d'une page pour les agents de la mairie
d'Eséka.
\end{boxexo}

\begin{boxexo}[title={\ding{45}\quad Exercice~6.10. Registre des traitements \niveauexpert}]
Dressez un extrait du registre des traitements de la mairie (trois traitements,
finalité, données, durée de conservation).
\end{boxexo}

\bigskip

\begin{boxinfo}
\textbf{Pour aller plus loin.}
\begin{itemize}
    \item Loi 2024/017~\cite{loi2024017}~; Convention de Malabo~\cite{malabo2014}.
\end{itemize}
\end{boxinfo}

\begin{boxinfo}
\textbf{Transition.} Gouvernance, politique et conformité posées, le cours
entre dans la norme qui les structure~: le Chapitre~\ref{chap:iso27000}
présente la famille \norme{ISO/IEC 27000}.
\end{boxinfo}
}{}

    \part{Le SMSI selon ISO/IEC 27001}
    \label{part:smsi}

        \IfFileExists{partie3/P03_C07_famille_27000.tex}{% ============================================================================
% Fichier  : partie3/P03_C07_famille_27000.tex
% Titre    : La Famille ISO/IEC 27000 et le SMSI
% Auteur   : MINKA MI NGUIDJOÏ Thierry Emmanuel
% ============================================================================
% Sources : ISO/IEC 27000:2018 ; 27001:2022 ; 27002, 27005. Fictif : Eséka.
% ============================================================================

\chapter{La Famille ISO/IEC 27000 et le SMSI}
\label{chap:iso27000}
\chaptermark{La famille ISO 27000}

\epigraph{%
    « Une norme ne protège rien par elle-même~:\\
    elle organise la manière dont on se protège. »%
}{Principe du cours}

\niveauChapitre{Fondamental}{La partie II.}

\begin{boxobjectifs}
\begin{itemize}
    \item Situer la famille \norme{ISO/IEC 27000} et ses principales normes.
    \item Distinguer la norme d'exigences (27001) des normes de lignes
          directrices.
    \item Comprendre ce qu'est un \sigle{SMSI} au sens de la norme.
    \item Expliquer le principe de la certification.
    \item Replacer la démarche de la mairie dans ce cadre.
\end{itemize}
\end{boxobjectifs}

\bigskip

% ============================================================================
\section*{Question directrice : pourquoi s'appuyer sur une norme plutôt que bricoler ?}
% ============================================================================

La mairie pourrait inventer sa propre méthode. S'appuyer sur
\norme{ISO/IEC 27001}, c'est hériter d'un cadre éprouvé, reconnu et auditable.
Ce chapitre présente la famille de normes et le \sigle{SMSI} qu'elle définit.

\medskip
\textbf{Plan du chapitre.}
\begin{enumerate}
    \item La famille 27000 (section~\ref{sec:c07-famille}).
    \item Exigences et lignes directrices (section~\ref{sec:c07-exigences}).
    \item Le SMSI selon la norme (section~\ref{sec:c07-smsi}).
    \item Le principe de la certification (section~\ref{sec:c07-certif}).
    \item La mairie dans le cadre (section~\ref{sec:c07-mairie}).
\end{enumerate}

\separateur

% ============================================================================
\section{La Famille 27000}
\label{sec:c07-famille}
% ============================================================================

\begin{tableaucours}{L{2.6cm}L{9.6cm}}{%
    \tableauHeader{Norme} &
    \tableauHeader{Objet}
}
ISO/IEC 27000 & Vue d'ensemble et vocabulaire \\
ISO/IEC 27001 & Exigences du SMSI (la seule certifiable) \\
ISO/IEC 27002 & Mesures de sécurité (code de bonne pratique) \\
ISO/IEC 27005 & Gestion des risques de sécurité de l'information \\
ISO/IEC 27004 & Mesure, indicateurs et évaluation \\
ISO/IEC 27006 & Exigences des organismes de certification \\
\end{tableaucours}
\vspace{-0.8em}
\begin{center}{\small\itshape Tableau~7.1. Quelques normes de la famille ISO/IEC 27000}\end{center}

La famille forme un système cohérent~: la 27001 \emph{exige}, les autres
\emph{aident}. On ne se certifie que sur la 27001~; les autres soutiennent sa
mise en \oe uvre.

% ============================================================================
\section{Exigences et Lignes Directrices}
\label{sec:c07-exigences}
% ============================================================================

\begin{boxdef}
Une norme d'\termecle{exigences} (\norme{ISO/IEC 27001}) énonce ce qui
\emph{doit} être fait pour être conforme~; ses clauses emploient « doit ». Une
norme de \termecle{lignes directrices} (27002, 27005) propose des bonnes
pratiques~; elle emploie « il convient de ».
\end{boxdef}

La distinction est capitale en audit~: un auditeur vérifie le respect des
\emph{exigences} de la 27001. Les lignes directrices éclairent le « comment »,
sans être, elles, directement certifiables.

% ============================================================================
\section{Le SMSI selon la Norme}
\label{sec:c07-smsi}
% ============================================================================

Le \sigle{SMSI} de la 27001 s'articule en sept clauses d'exigences
(clauses~4 à~10), qui suivent le cycle \sigle{PDCA}~:

\begin{tableaucours}{L{2.0cm}L{6.6cm}L{3.2cm}}{%
    \tableauHeader{Clause} &
    \tableauHeader{Exigence} &
    \tableauHeader{Étape PDCA}
}
4 & Contexte de l'organisme & Plan \\
5 & Leadership & Plan \\
6 & Planification & Plan \\
7 & Support & Do \\
8 & Fonctionnement & Do \\
9 & Évaluation des performances & Check \\
10 & Amélioration & Act \\
\end{tableaucours}
\vspace{-0.8em}
\begin{center}{\small\itshape Tableau~7.2. Les clauses d'exigences et le PDCA}\end{center}

Les Chapitres~\ref{chap:clauses46} et~\ref{chap:clauses710} détaillent ces
clauses~; le Chapitre~\ref{chap:iso27002} traite les mesures de l'Annexe~A.

% ============================================================================
\section{Le Principe de la Certification}
\label{sec:c07-certif}
% ============================================================================

\begin{boxdef}
\termecle{Certification \norme{ISO/IEC 27001}}~: attestation, délivrée par un
organisme accrédité, qu'un \sigle{SMSI} satisfait aux exigences de la norme sur
un périmètre défini. Elle se gagne par un audit et se maintient par des audits
de surveillance.
\end{boxdef}

La certification ne prouve pas l'absence d'incident~: elle atteste que
l'organisation \emph{gère} sa sécurité selon la norme. C'est une reconnaissance
de maturité, non une garantie d'invulnérabilité.

\begin{boxmetier}{Auditeur SMSI}
\textbf{Auditer l'exigence, pas l'intention.} L'auditeur ne note pas la bonne
volonté~: il vérifie, preuve à l'appui, que chaque « doit » de la norme est
satisfait. Une mairie motivée mais sans preuves échoue.
\end{boxmetier}

% ============================================================================
\section{La Mairie dans le Cadre}
\label{sec:c07-mairie}
% ============================================================================

\begin{boxprocedure}
\textbf{Constatation.} La mairie d'Eséka vise la certification \norme{ISO/IEC
27001} de son \sigle{SMSI} sur le périmètre « services aux usagers » (état
civil, taxes).

\medskip
\textbf{Analyse.} Le choix du périmètre est stratégique~: commencer par un
périmètre maîtrisable, quitte à l'étendre ensuite. Les quatre phases du fil
rouge (\textsc{Nyong}, \textsc{Kellé}, \textsc{Makak}, \textsc{Matomb})
couvrent les clauses~4 à~10.
\end{boxprocedure}

\separateur

\begin{boxresume}
\textbf{Ce qu'il faut retenir de ce chapitre~:}
\begin{itemize}
    \item La \textbf{famille 27000} forme un système~: la 27001 exige, les
          autres aident.
    \item Seule \textbf{ISO/IEC 27001} est \textbf{certifiable}~; elle emploie
          « doit ».
    \item Le \textbf{SMSI} de la norme tient en sept clauses (4 à~10) alignées
          sur le \sigle{PDCA}.
    \item La \textbf{certification} atteste la maîtrise, non l'absence
          d'incident.
    \item La mairie vise un \textbf{périmètre} défini et maîtrisable.
\end{itemize}
\end{boxresume}

\begin{boxinfo}
\textbf{Points de contrôle de l'élève.} Je sais~: citer les normes de la
famille 27000~; distinguer exigences et lignes directrices~; énoncer les sept
clauses~; expliquer la certification~; définir un périmètre.
\end{boxinfo}

\bigskip

\begin{boxexo}[title={\ding{45}\quad Exercice~7.1. La famille \niveaubase}]
Associez chaque norme à son objet~: 27001, 27002, 27005.
\end{boxexo}

\begin{boxexo}[title={\ding{45}\quad Exercice~7.2. Doit ou il convient \niveaubase}]
Laquelle de ces normes emploie « doit »~: 27001 ou 27002~? Pourquoi est-ce
important~?
\end{boxexo}

\begin{boxexo}[title={\ding{45}\quad Exercice~7.3. Les clauses \niveaubase}]
Citez les sept clauses d'exigences de la 27001 et leur étape PDCA.
\end{boxexo}

\begin{boxexo}[title={\ding{45}\quad Exercice~7.4. Ce que prouve la certification \niveaumoyen}]
« Une mairie certifiée ne peut pas être victime d'un incident. » Corrigez cette
affirmation.
\end{boxexo}

\begin{boxexo}[title={\ding{45}\quad Exercice~7.5. Rôle des lignes directrices \niveaumoyen}]
À quoi servent la 27002 et la 27005 si l'on se certifie sur la 27001~?
\end{boxexo}

\begin{boxexo}[title={\ding{45}\quad Exercice~7.6. Périmètre \niveaumoyen}]
Pourquoi la mairie commence-t-elle par un périmètre restreint~?
\end{boxexo}

\begin{boxexo}[title={\ding{45}\quad Exercice~7.7. Preuve \niveaumoyen}]
Pourquoi une mairie motivée mais sans preuves échoue-t-elle à l'audit~?
\end{boxexo}

\begin{boxexo}[title={\ding{45}\quad Exercice~7.8. PDCA et clauses \niveaumoyen}]
Reliez les clauses 9 et 10 aux étapes Check et Act, avec un exemple.
\end{boxexo}

\begin{boxexo}[title={\ding{45}\quad Exercice~7.9. Carte des normes \niveauexpert}]
Dressez une carte des normes 27000 utiles à la mairie, en indiquant pour chaque
étape du projet laquelle mobiliser.
\end{boxexo}

\begin{boxexo}[title={\ding{45}\quad Exercice~7.10. Note de périmètre \niveauexpert}]
Rédigez la note qui définit et justifie le périmètre de certification de la
mairie d'Eséka.
\end{boxexo}

\bigskip

\begin{boxinfo}
\textbf{Pour aller plus loin.}
\begin{itemize}
    \item ISO/IEC 27000~\cite{iso27000}~; ISO/IEC 27001~\cite{iso27001}.
\end{itemize}
\end{boxinfo}

\begin{boxinfo}
\textbf{Transition.} Le Chapitre~\ref{chap:clauses46} entre dans les premières
clauses~: contexte, leadership et planification.
\end{boxinfo}
}{}
        \IfFileExists{partie3/P03_C08_clauses_4_6.tex}{% ============================================================================
% Fichier  : partie3/P03_C08_clauses_4_6.tex
% Titre    : Contexte, Leadership, Planification
% Auteur   : MINKA MI NGUIDJOÏ Thierry Emmanuel
% ============================================================================
% Sources : ISO/IEC 27001:2022, clauses 4 à 6. Mairie d'Eséka : cas fictif.
% ============================================================================

\chapter{Contexte, Leadership, Planification}
\label{chap:clauses46}
\chaptermark{Clauses 4 à 6}

\epigraph{%
    « Avant de planifier la sécurité, il faut savoir pour qui,\\
    contre quoi et dans quelles limites. »%
}{Principe du cours}

\niveauChapitre{Intermédiaire}{Le Chapitre~\ref{chap:iso27000}.}

\begin{boxobjectifs}
\begin{itemize}
    \item Déterminer le contexte et les parties intéressées (clause~4).
    \item Préciser le domaine d'application du \sigle{SMSI}.
    \item Caractériser le leadership exigé de la direction (clause~5).
    \item Planifier les actions face aux risques et les objectifs (clause~6).
    \item Comprendre le lien entre objectifs de sécurité et appréciation des
          risques.
\end{itemize}
\end{boxobjectifs}

\bigskip

% ============================================================================
\section*{Question directrice : comment poser les fondations d'un SMSI conforme ?}
% ============================================================================

Les trois premières clauses d'exigences fixent le cadre~: qui, quoi, pourquoi,
et avec quels objectifs. Ce chapitre les traduit pour la mairie d'Eséka.

\medskip
\textbf{Plan du chapitre.}
\begin{enumerate}
    \item Le contexte (clause~4) (section~\ref{sec:c08-contexte}).
    \item Le domaine d'application (section~\ref{sec:c08-domaine}).
    \item Le leadership (clause~5) (section~\ref{sec:c08-leadership}).
    \item La planification (clause~6) (section~\ref{sec:c08-planif}).
    \item Les objectifs de sécurité (section~\ref{sec:c08-objectifs}).
\end{enumerate}

\separateur

% ============================================================================
\section{Le Contexte (Clause 4)}
\label{sec:c08-contexte}
% ============================================================================

\begin{boxdef}
La clause~4 exige de déterminer les \termecle{enjeux internes et externes}
pertinents, les \termecle{parties intéressées} et leurs attentes, puis de
fixer le \termecle{domaine d'application} du \sigle{SMSI}.
\end{boxdef}

\begin{tableaucours}{L{3.2cm}L{9.0cm}}{%
    \tableauHeader{Élément} &
    \tableauHeader{Pour la mairie d'Eséka}
}
Enjeux externes & Loi 2024/017, attentes des usagers, menace cyber \\
Enjeux internes & Budget contraint, compétences limitées, matériel vieillissant \\
Parties intéressées & Usagers, tutelle, ANTIC, prestataires, agents \\
\end{tableaucours}
\vspace{-0.8em}
\begin{center}{\small\itshape Tableau~8.1. Le contexte de la mairie (clause 4)}\end{center}

% ============================================================================
\section{Le Domaine d'Application}
\label{sec:c08-domaine}
% ============================================================================

Le \termecle{domaine d'application} (\emph{scope}) dit ce que couvre le
\sigle{SMSI}~: quels services, sites, actifs et exclusions. Il doit être écrit
et justifié~: un périmètre flou rend l'audit impossible.

\begin{boxprocedure}
\textbf{Constatation.} La mairie fixe le domaine~: « le système d'information
des services aux usagers de la mairie d'Eséka (état civil et taxes), à
l'exclusion du réseau technique de la voirie ».

\medskip
\textbf{Analyse.} L'exclusion est explicite et justifiée (la voirie ne traite
pas de données sensibles). Un domaine clair protège contre l'extension
incontrôlée et cadre l'audit.
\end{boxprocedure}

% ============================================================================
\section{Le Leadership (Clause 5)}
\label{sec:c08-leadership}
% ============================================================================

La clause~5 exige que la direction \termecle{démontre} son leadership~: établir
la politique, attribuer rôles et responsabilités, fournir les ressources,
promouvoir l'amélioration continue. Ce n'est pas une déclaration~: ce sont des
actes vérifiables (Chapitre~\ref{chap:gouvernance}).

\begin{boxmetier}{Direction générale}
\textbf{Des preuves de leadership.} Le maire prouve son engagement par une
politique signée, un budget voté, une RSSI nommée et sa présence au comité.
L'auditeur demandera ces preuves~; les bonnes intentions ne comptent pas.
\end{boxmetier}

% ============================================================================
\section{La Planification (Clause 6)}
\label{sec:c08-planif}
% ============================================================================

\begin{boxdef}
La clause~6 exige de planifier les \termecle{actions face aux risques et
opportunités}, de définir le \termecle{processus d'appréciation et de
traitement des risques}, et de fixer des \termecle{objectifs de sécurité}
mesurables.
\end{boxdef}

La planification est le cœur « Plan » du \sigle{PDCA}. C'est ici que la norme
renvoie à l'appréciation des risques (partie IV)~: on ne planifie pas au
hasard, mais au vu des risques appréciés.

% ============================================================================
\section{Les Objectifs de Sécurité}
\label{sec:c08-objectifs}
% ============================================================================

\begin{boxdef}
Un \termecle{objectif de sécurité} est un résultat visé, cohérent avec la
politique, \emph{mesurable}, suivi et mis à jour. « Améliorer la sécurité »
n'est pas un objectif~; « sauvegarder quotidiennement l'état civil et le
vérifier chaque mois » en est un.
\end{boxdef}

\begin{tableaucours}{L{5.4cm}L{3.0cm}L{3.4cm}}{%
    \tableauHeader{Objectif} &
    \tableauHeader{Indicateur} &
    \tableauHeader{Cible}
}
Sauvegarder l'état civil & Fréquence & Quotidienne \\
Former les agents & Taux de formation & 100 \% en 6 mois \\
Réduire les accès partagés & Comptes nominatifs & 100 \% \\
\end{tableaucours}
\vspace{-0.8em}
\begin{center}{\small\itshape Tableau~8.2. Des objectifs de sécurité mesurables}\end{center}

\separateur

\begin{boxresume}
\textbf{Ce qu'il faut retenir de ce chapitre~:}
\begin{itemize}
    \item La \textbf{clause~4} détermine contexte, parties intéressées et
          domaine d'application.
    \item Le \textbf{domaine d'application} doit être écrit, justifié, avec ses
          exclusions.
    \item La \textbf{clause~5} exige un leadership \emph{prouvé} de la
          direction.
    \item La \textbf{clause~6} planifie les actions face aux risques et fixe
          les objectifs.
    \item Un \textbf{objectif} de sécurité est mesurable et suivi.
\end{itemize}
\end{boxresume}

\begin{boxinfo}
\textbf{Points de contrôle de l'élève.} Je sais~: déterminer le contexte et les
parties intéressées~; rédiger un domaine d'application~; citer les preuves de
leadership~; expliquer la planification~; formuler un objectif mesurable.
\end{boxinfo}

\bigskip

\begin{boxexo}[title={\ding{45}\quad Exercice~8.1. Contexte \niveaubase}]
Donnez deux enjeux externes et deux enjeux internes de la mairie d'Eséka.
\end{boxexo}

\begin{boxexo}[title={\ding{45}\quad Exercice~8.2. Parties intéressées \niveaubase}]
Citez quatre parties intéressées et une attente de sécurité pour chacune.
\end{boxexo}

\begin{boxexo}[title={\ding{45}\quad Exercice~8.3. Objectif mesurable \niveaubase}]
Transformez « mieux gérer les mots de passe » en objectif mesurable.
\end{boxexo}

\begin{boxexo}[title={\ding{45}\quad Exercice~8.4. Domaine d'application \niveaumoyen}]
Rédigez le domaine d'application du \sigle{SMSI} de la mairie, avec une
exclusion justifiée.
\end{boxexo}

\begin{boxexo}[title={\ding{45}\quad Exercice~8.5. Preuves de leadership \niveaumoyen}]
Citez quatre preuves par lesquelles le maire démontre son leadership.
\end{boxexo}

\begin{boxexo}[title={\ding{45}\quad Exercice~8.6. Périmètre flou \niveaumoyen}]
Pourquoi un domaine d'application flou rend-il l'audit impossible~?
\end{boxexo}

\begin{boxexo}[title={\ding{45}\quad Exercice~8.7. Planifier par le risque \niveaumoyen}]
Expliquez pourquoi la clause~6 renvoie à l'appréciation des risques.
\end{boxexo}

\begin{boxexo}[title={\ding{45}\quad Exercice~8.8. Trois objectifs \niveaumoyen}]
Proposez trois objectifs de sécurité mesurables pour la mairie, avec cible.
\end{boxexo}

\begin{boxexo}[title={\ding{45}\quad Exercice~8.9. Dossier de contexte \niveauexpert}]
Rédigez le dossier de contexte (clause~4) complet de la mairie d'Eséka.
\end{boxexo}

\begin{boxexo}[title={\ding{45}\quad Exercice~8.10. Plan d'objectifs \niveauexpert}]
Construisez un tableau de cinq objectifs de sécurité, avec indicateur, cible,
responsable et échéance.
\end{boxexo}

\bigskip

\begin{boxinfo}
\textbf{Pour aller plus loin.}
\begin{itemize}
    \item ISO/IEC 27001, clauses~4 à~6~\cite{iso27001}.
\end{itemize}
\end{boxinfo}

\begin{boxinfo}
\textbf{Transition.} Le Chapitre~\ref{chap:clauses710} poursuit avec le support,
le fonctionnement, l'évaluation et l'amélioration (clauses~7 à~10).
\end{boxinfo}
}{}
        \IfFileExists{partie3/P03_C09_clauses_7_10.tex}{% ============================================================================
% Fichier  : partie3/P03_C09_clauses_7_10.tex
% Titre    : Support, Fonctionnement, Évaluation, Amélioration
% Auteur   : MINKA MI NGUIDJOÏ Thierry Emmanuel
% ============================================================================
% Sources : ISO/IEC 27001:2022, clauses 7 à 10. Mairie d'Eséka : cas fictif.
% ============================================================================

\chapter{Support, Fonctionnement, Évaluation, Amélioration}
\label{chap:clauses710}
\chaptermark{Clauses 7 à 10}

\epigraph{%
    « Un système se juge à ce qu'il produit comme preuves,\\
    pas à ce qu'il déclare comme intentions. »%
}{Principe du cours}

\niveauChapitre{Intermédiaire}{Le Chapitre~\ref{chap:clauses46}.}

\begin{boxobjectifs}
\begin{itemize}
    \item Identifier les ressources, compétences et communications exigées
          (clause~7).
    \item Comprendre l'exigence d'informations documentées.
    \item Décrire le fonctionnement et la maîtrise opérationnelle (clause~8).
    \item Expliquer l'évaluation des performances (clause~9).
    \item Caractériser l'amélioration et les non-conformités (clause~10).
\end{itemize}
\end{boxobjectifs}

\bigskip

% ============================================================================
\section*{Question directrice : comment faire vivre et prouver le SMSI ?}
% ============================================================================

Les clauses~7 à~10 font fonctionner le \sigle{SMSI}, le mesurent et
l'améliorent. Elles correspondent aux étapes Do, Check et Act du
\sigle{PDCA}. Ce chapitre les parcourt pour la mairie.

\medskip
\textbf{Plan du chapitre.}
\begin{enumerate}
    \item Le support (clause~7) (section~\ref{sec:c09-support}).
    \item Les informations documentées (section~\ref{sec:c09-doc}).
    \item Le fonctionnement (clause~8) (section~\ref{sec:c09-fonction}).
    \item L'évaluation des performances (clause~9) (section~\ref{sec:c09-eval}).
    \item L'amélioration (clause~10) (section~\ref{sec:c09-amelioration}).
\end{enumerate}

\separateur

% ============================================================================
\section{Le Support (Clause 7)}
\label{sec:c09-support}
% ============================================================================

\begin{boxdef}
La clause~7 exige de fournir les \termecle{ressources}, d'assurer les
\termecle{compétences}, la \termecle{sensibilisation}, la \termecle{communication}
et la maîtrise des \termecle{informations documentées}.
\end{boxdef}

Un \sigle{SMSI} sans moyens ni agents formés reste un document. La mairie doit
budgéter la sécurité, former ses agents (Chapitre~\ref{chap:orga}) et
communiquer ses règles.

% ============================================================================
\section{Les Informations Documentées}
\label{sec:c09-doc}
% ============================================================================

\begin{boxdef}
\termecle{Information documentée}~: document ou enregistrement que la norme
exige de tenir à jour et de maîtriser (politique, procédures, registres,
preuves d'audit). « Documenté » signifie tracé et retrouvable.
\end{boxdef}

La maîtrise documentaire est l'épine dorsale de l'audit~: ce qui n'est pas
documenté est réputé ne pas exister. Chaque exigence satisfaite laisse une
preuve écrite.

\begin{boxprocedure}
\textbf{Constatation.} La mairie sauvegarde l'état civil mais n'en conserve
aucune trace écrite (ni journal, ni test de restauration).

\medskip
\textbf{Analyse.} Même réelle, la sauvegarde est invérifiable~: sans
information documentée, l'auditeur ne peut la constater. La mesure existe sans
preuve~: en audit, elle ne compte pas.
\end{boxprocedure}

% ============================================================================
\section{Le Fonctionnement (Clause 8)}
\label{sec:c09-fonction}
% ============================================================================

La clause~8 exige de \termecle{planifier et maîtriser} les opérations, de
réaliser l'\termecle{appréciation des risques} à intervalles planifiés et de
mettre en \oe uvre le \termecle{traitement des risques}. C'est le « Do » du
cycle~: on exécute ce que la planification a décidé (parties IV et V).

% ============================================================================
\section{L'Évaluation des Performances (Clause 9)}
\label{sec:c09-eval}
% ============================================================================

\begin{tableaucours}{L{4.2cm}L{8.0cm}}{%
    \tableauHeader{Exigence} &
    \tableauHeader{Contenu}
}
Surveillance et mesure & Suivre les indicateurs de sécurité \\
Audit interne & Vérifier la conformité du SMSI \\
Revue de direction & La direction examine et décide \\
\end{tableaucours}
\vspace{-0.8em}
\begin{center}{\small\itshape Tableau~9.1. L'évaluation des performances (clause 9)}\end{center}

La clause~9 est le « Check »~: mesurer (Chapitre~\ref{chap:indicateurs}),
auditer et passer en revue de direction (Chapitre~\ref{chap:audit}). Sans
mesure, on ne sait pas si le \sigle{SMSI} fonctionne.

% ============================================================================
\section{L'Amélioration (Clause 10)}
\label{sec:c09-amelioration}
% ============================================================================

\begin{boxdef}
La clause~10 exige de traiter les \termecle{non-conformités} par des actions
correctives qui en éliminent la cause, et d'assurer l'\termecle{amélioration
continue} du \sigle{SMSI}.
\end{boxdef}

C'est le « Act »~: corriger les écarts à leur cause (Chapitre~\ref{chap:pdca})
et relancer le cycle. La boucle est bouclée~: de la planification (clause~6) à
l'amélioration (clause~10), le \sigle{SMSI} tourne.

\begin{boxmetier}{RSSI}
\textbf{Documenter pour durer.} La RSSI tient un dossier vivant~: politique,
appréciation des risques, plan de traitement, journaux, rapports d'audit,
revues de direction. Ce dossier \emph{est} la preuve que le \sigle{SMSI}
fonctionne.
\end{boxmetier}

\separateur

\begin{boxresume}
\textbf{Ce qu'il faut retenir de ce chapitre~:}
\begin{itemize}
    \item La \textbf{clause~7} fournit ressources, compétences, communication
          et informations documentées.
    \item Ce qui n'est pas \textbf{documenté} est réputé inexistant en audit.
    \item La \textbf{clause~8} exécute l'appréciation et le traitement des
          risques.
    \item La \textbf{clause~9} mesure, audite et passe en revue de direction.
    \item La \textbf{clause~10} corrige à la cause et améliore en continu.
\end{itemize}
\end{boxresume}

\begin{boxinfo}
\textbf{Points de contrôle de l'élève.} Je sais~: citer les exigences de
support~; définir une information documentée~; expliquer le fonctionnement~;
décrire l'évaluation des performances~; caractériser l'amélioration.
\end{boxinfo}

\bigskip

\begin{boxexo}[title={\ding{45}\quad Exercice~9.1. La clause 7 \niveaubase}]
Citez trois éléments de support exigés par la clause~7.
\end{boxexo}

\begin{boxexo}[title={\ding{45}\quad Exercice~9.2. Documenté \niveaubase}]
Qu'entend la norme par « information documentée »~? Donnez deux exemples.
\end{boxexo}

\begin{boxexo}[title={\ding{45}\quad Exercice~9.3. Les trois exigences \niveaubase}]
Citez les trois exigences de l'évaluation des performances (clause~9).
\end{boxexo}

\begin{boxexo}[title={\ding{45}\quad Exercice~9.4. Mesure sans preuve \niveaumoyen}]
Pourquoi une sauvegarde réelle mais non tracée ne compte-t-elle pas en audit~?
\end{boxexo}

\begin{boxexo}[title={\ding{45}\quad Exercice~9.5. PDCA et clauses \niveaumoyen}]
Reliez les clauses~7 à~10 aux étapes Do, Check et Act.
\end{boxexo}

\begin{boxexo}[title={\ding{45}\quad Exercice~9.6. Non-conformité \niveaumoyen}]
Distinguez corriger une non-conformité et en traiter la cause, par un exemple.
\end{boxexo}

\begin{boxexo}[title={\ding{45}\quad Exercice~9.7. Dossier de preuves \niveaumoyen}]
Listez cinq documents du dossier du \sigle{SMSI} de la mairie.
\end{boxexo}

\begin{boxexo}[title={\ding{45}\quad Exercice~9.8. Revue de direction \niveaumoyen}]
À quoi sert la revue de direction, et qui y participe~?
\end{boxexo}

\begin{boxexo}[title={\ding{45}\quad Exercice~9.9. Procédure documentaire \niveauexpert}]
Rédigez la procédure de maîtrise des informations documentées de la mairie
(création, validation, diffusion, archivage).
\end{boxexo}

\begin{boxexo}[title={\ding{45}\quad Exercice~9.10. Tableau clauses-preuves \niveauexpert}]
Pour chaque clause 7 à 10, indiquez une preuve documentée que la mairie devra
produire en audit.
\end{boxexo}

\bigskip

\begin{boxinfo}
\textbf{Pour aller plus loin.}
\begin{itemize}
    \item ISO/IEC 27001, clauses~7 à~10~\cite{iso27001}.
\end{itemize}
\end{boxinfo}

\begin{boxinfo}
\textbf{Transition.} Les exigences posées, reste le catalogue des mesures~: le
Chapitre~\ref{chap:iso27002} présente \norme{ISO/IEC 27002} et l'Annexe~A.
\end{boxinfo}
}{}
        \IfFileExists{partie3/P03_C10_iso27002.tex}{% ============================================================================
% Fichier  : partie3/P03_C10_iso27002.tex
% Titre    : ISO/IEC 27002 et les Mesures de Sécurité
% Auteur   : MINKA MI NGUIDJOÏ Thierry Emmanuel
% ============================================================================
% Sources : ISO/IEC 27002:2022 ; ISO/IEC 27001:2022, Annexe A. Fictif : Eséka.
% ============================================================================

\chapter{ISO/IEC 27002 et les Mesures de Sécurité}
\label{chap:iso27002}
\chaptermark{ISO 27002 et l'Annexe A}

\epigraph{%
    « Les mesures ne se choisissent pas dans un catalogue par goût~:\\
    elles se justifient par le risque. »%
}{Principe du cours}

\niveauChapitre{Intermédiaire}{Le Chapitre~\ref{chap:iso27000}.}

\begin{boxobjectifs}
\begin{itemize}
    \item Situer \norme{ISO/IEC 27002} et l'Annexe~A de la 27001.
    \item Énoncer les quatre thèmes de mesures de la version 2022.
    \item Comprendre la structure d'une mesure (attributs, objectif).
    \item Relier mesures et appréciation des risques.
    \item Introduire la déclaration d'applicabilité.
\end{itemize}
\end{boxobjectifs}

\bigskip

% ============================================================================
\section*{Question directrice : parmi toutes les mesures possibles, lesquelles retenir ?}
% ============================================================================

La 27002 décrit des dizaines de mesures. La mairie ne les appliquera pas
toutes~: elle retiendra celles que son risque justifie. Ce chapitre présente le
catalogue et prépare le choix.

\medskip
\textbf{Plan du chapitre.}
\begin{enumerate}
    \item 27002 et Annexe~A (section~\ref{sec:c10-annexea}).
    \item Les quatre thèmes (section~\ref{sec:c10-themes}).
    \item Structure d'une mesure (section~\ref{sec:c10-structure}).
    \item Mesures et risques (section~\ref{sec:c10-risques}).
    \item Vers la déclaration d'applicabilité (section~\ref{sec:c10-soa}).
\end{enumerate}

\separateur

% ============================================================================
\section{27002 et Annexe A}
\label{sec:c10-annexea}
% ============================================================================

\begin{boxdef}
L'\termecle{Annexe~A} de \norme{ISO/IEC 27001} liste les mesures de sécurité de
référence. La norme \norme{ISO/IEC 27002} en donne le \termecle{guide de mise
en \oe uvre}~: pour chaque mesure, son objectif et des conseils d'application.
\end{boxdef}

L'Annexe~A fournit la liste de contrôle~; la 27002 explique comment faire. La
version 2022 réorganise les mesures en quatre thèmes et introduit des
\termecle{attributs} (type de mesure, propriété de sécurité, fonction).

% ============================================================================
\section{Les Quatre Thèmes}
\label{sec:c10-themes}
% ============================================================================

\begin{tableaucours}{L{3.4cm}L{2.0cm}L{6.6cm}}{%
    \tableauHeader{Thème} &
    \tableauHeader{Nombre} &
    \tableauHeader{Exemples}
}
Organisationnelles & 37 & Politiques, rôles, fournisseurs, incidents \\
Liées aux personnes & 8 & Sensibilisation, filtrage, télétravail \\
Physiques & 14 & Accès aux locaux, supports, équipements \\
Technologiques & 34 & Accès, chiffrement, journalisation, sauvegarde \\
\end{tableaucours}
\vspace{-0.8em}
\begin{center}{\small\itshape Tableau~10.1. Les quatre thèmes de l'ISO/IEC 27002:2022}\end{center}

Au total, 93~mesures. Les parties V du cours détaillent les mesures
organisationnelles et humaines (Chapitre~\ref{chap:orga}), techniques
(Chapitre~\ref{chap:technique}) et de continuité (Chapitre~\ref{chap:continuite}).

% ============================================================================
\section{Structure d'une Mesure}
\label{sec:c10-structure}
% ============================================================================

\begin{boxprocedure}
\textbf{Constatation.} La mesure « sauvegarde de l'information » (27002) a pour
objectif de garantir la \emph{disponibilité}~; ses attributs la classent en
mesure préventive et corrective.

\medskip
\textbf{Analyse.} Chaque mesure a un objectif (ce qu'elle protège) et un guide
(comment l'appliquer). À la mairie, cette mesure impose une règle de
sauvegarde quotidienne et un test de restauration (Chapitre~\ref{chap:continuite}).
\end{boxprocedure}

Comprendre la structure évite l'application mécanique~: on applique une mesure
pour son \emph{objectif}, pas pour cocher une case.

% ============================================================================
\section{Mesures et Risques}
\label{sec:c10-risques}
% ============================================================================

La norme est claire~: on ne retient pas une mesure parce qu'elle figure dans
l'Annexe~A, mais parce que l'\termecle{appréciation des risques} (partie IV) la
justifie. Les mesures sont la réponse au risque, pas une liste à appliquer en
bloc.

\begin{boxmetier}{Auditeur SMSI}
\textbf{La traçabilité risque-mesure.} L'auditeur vérifie que chaque mesure
retenue répond à un risque identifié, et que chaque risque significatif reçoit
une mesure. Une mesure sans risque et un risque sans mesure sont deux écarts.
\end{boxmetier}

% ============================================================================
\section{Vers la Déclaration d'Applicabilité}
\label{sec:c10-soa}
% ============================================================================

\begin{boxdef}
La \termecle{déclaration d'applicabilité} (\sigle{DdA}, \emph{Statement of
Applicability})~: document qui, pour chaque mesure de l'Annexe~A, indique si
elle est retenue, sa justification, et son état de mise en \oe uvre.
\end{boxdef}

La \sigle{DdA} est la pièce maîtresse qui relie les risques aux mesures. Elle
est détaillée au Chapitre~\ref{chap:dda}, après l'étude des risques. Retenir
ou écarter une mesure s'y justifie explicitement.

\separateur

\begin{boxresume}
\textbf{Ce qu'il faut retenir de ce chapitre~:}
\begin{itemize}
    \item L'\textbf{Annexe~A} liste les mesures~; la \textbf{27002} guide leur
          mise en \oe uvre.
    \item La version 2022 compte \textbf{93 mesures} en quatre thèmes.
    \item Chaque mesure a un \textbf{objectif}~; on l'applique pour cet
          objectif.
    \item Les mesures se justifient par le \textbf{risque}, non par le
          catalogue.
    \item La \textbf{déclaration d'applicabilité} relie risques et mesures.
    \end{itemize}
\end{boxresume}

\begin{boxinfo}
\textbf{Points de contrôle de l'élève.} Je sais~: distinguer Annexe~A et
27002~; citer les quatre thèmes~; décrire la structure d'une mesure~; relier
mesure et risque~; définir la \sigle{DdA}.
\end{boxinfo}

\bigskip

\begin{boxexo}[title={\ding{45}\quad Exercice~10.1. Annexe A ou 27002 \niveaubase}]
Laquelle liste les mesures, laquelle en guide la mise en \oe uvre~?
\end{boxexo}

\begin{boxexo}[title={\ding{45}\quad Exercice~10.2. Les quatre thèmes \niveaubase}]
Citez les quatre thèmes de mesures et un exemple pour chacun.
\end{boxexo}

\begin{boxexo}[title={\ding{45}\quad Exercice~10.3. Objectif d'une mesure \niveaubase}]
Quel est l'objectif de la mesure « sauvegarde de l'information »~?
\end{boxexo}

\begin{boxexo}[title={\ding{45}\quad Exercice~10.4. Mesure et risque \niveaumoyen}]
Pourquoi n'applique-t-on pas toutes les mesures de l'Annexe~A~?
\end{boxexo}

\begin{boxexo}[title={\ding{45}\quad Exercice~10.5. Classer des mesures \niveaumoyen}]
Classez par thème~: chiffrement, sensibilisation, contrôle d'accès aux locaux,
gestion des fournisseurs.
\end{boxexo}

\begin{boxexo}[title={\ding{45}\quad Exercice~10.6. Deux écarts \niveaumoyen}]
Qu'est-ce qu'une « mesure sans risque » et un « risque sans mesure »~?
\end{boxexo}

\begin{boxexo}[title={\ding{45}\quad Exercice~10.7. Mesures pour la mairie \niveaumoyen}]
Citez trois mesures de l'Annexe~A particulièrement utiles à la mairie d'Eséka.
\end{boxexo}

\begin{boxexo}[title={\ding{45}\quad Exercice~10.8. Rôle de la DdA \niveaumoyen}]
À quoi sert la déclaration d'applicabilité~?
\end{boxexo}

\begin{boxexo}[title={\ding{45}\quad Exercice~10.9. Fiche de mesure \niveauexpert}]
Rédigez une fiche de mesure pour « contrôle d'accès » (objectif, mise en
\oe uvre à la mairie, preuve attendue).
\end{boxexo}

\begin{boxexo}[title={\ding{45}\quad Exercice~10.10. Extrait de DdA \niveauexpert}]
Rédigez un extrait de \sigle{DdA} pour cinq mesures (retenue~? justification~?
état~?).
\end{boxexo}

\bigskip

\begin{boxinfo}
\textbf{Pour aller plus loin.}
\begin{itemize}
    \item ISO/IEC 27002~\cite{iso27002}~; ISO/IEC 27001, Annexe~A~\cite{iso27001}.
\end{itemize}
\end{boxinfo}

\begin{boxinfo}
\textbf{Transition.} Choisir des mesures suppose de connaître ses risques~: la
partie IV s'ouvre au Chapitre~\ref{chap:vocabrisque} sur le vocabulaire du
risque.
\end{boxinfo}
}{}

    \part{L'Appréciation et le Traitement des Risques}
    \label{part:risques}

        \IfFileExists{partie4/P04_C11_vocabulaire_risque.tex}{% ============================================================================
% Fichier  : partie4/P04_C11_vocabulaire_risque.tex
% Titre    : Concepts et Vocabulaire du Risque
% Auteur   : MINKA MI NGUIDJOÏ Thierry Emmanuel
% ============================================================================
% Sources : ISO/IEC 27005:2022 ; ISO 31000:2018. Mairie d'Eséka : cas fictif.
% ============================================================================

\chapter{Concepts et Vocabulaire du Risque}
\label{chap:vocabrisque}
\chaptermark{Vocabulaire du risque}

\epigraph{%
    « On ne traite bien que les risques que l'on a su nommer. »%
}{Principe du cours}

\niveauChapitre{Fondamental}{La partie III.}

\begin{boxobjectifs}
\begin{itemize}
    \item Définir actif, menace, vulnérabilité, impact et probabilité.
    \item Énoncer les formules du risque et les interpréter.
    \item Distinguer risque inhérent et risque résiduel.
    \item Comprendre l'appétence et les critères d'acceptation du risque.
    \item Appliquer le vocabulaire à un cas de la mairie.
\end{itemize}
\end{boxobjectifs}

\bigskip

% ============================================================================
\section*{Question directrice : qu'est-ce, au juste, qu'un risque ?}
% ============================================================================

Le mot « risque » est employé à tort et à travers. Pour l'apprécier, il faut
le décomposer en éléments précis. Ce chapitre fixe le vocabulaire commun à
toutes les méthodes.

\medskip
\textbf{Plan du chapitre.}
\begin{enumerate}
    \item Les cinq notions (section~\ref{sec:c11-notions}).
    \item La formule du risque (section~\ref{sec:c11-formule}).
    \item Inhérent et résiduel (section~\ref{sec:c11-residuel}).
    \item Appétence et acceptation (section~\ref{sec:c11-appetence}).
    \item Un cas appliqué (section~\ref{sec:c11-cas}).
\end{enumerate}

\separateur

% ============================================================================
\section{Les Cinq Notions}
\label{sec:c11-notions}
% ============================================================================

\begin{boxdef}
\begin{itemize}
    \item \termecle{Actif}~: élément de valeur à protéger.
    \item \termecle{Menace}~: cause potentielle d'un incident indésirable.
    \item \termecle{Vulnérabilité}~: faiblesse exploitable par une menace.
    \item \termecle{Impact}~: conséquence d'un incident.
    \item \termecle{Probabilité} (ou vraisemblance)~: chance qu'un incident se
          produise.
\end{itemize}
\end{boxdef}

Ces notions s'enchaînent~: une \emph{menace} exploite une \emph{vulnérabilité}
d'un \emph{actif}, avec une certaine \emph{probabilité}, et produit un
\emph{impact}. Supprimer l'un des maillons, c'est réduire le risque.

% ============================================================================
\section{La Formule du Risque}
\label{sec:c11-formule}
% ============================================================================

\begin{boxdef}
Le \termecle{risque} se formule, selon le niveau de détail, par~:
\[
\text{Risque} = \text{Menace} \times \text{Vulnérabilité} \times \text{Impact}
\quad\text{ou}\quad
\text{Risque} = \text{Probabilité} \times \text{Impact}.
\]
\end{boxdef}

Ces formules ne sont pas des calculs exacts mais des \emph{grilles de
raisonnement}~: elles rappellent que le risque croît avec la probabilité et
avec l'impact, et qu'un impact nul (actif sans valeur) annule le risque.

% ============================================================================
\section{Inhérent et Résiduel}
\label{sec:c11-residuel}
% ============================================================================

\begin{boxdef}
\termecle{Risque inhérent}~: niveau de risque avant toute mesure.
\termecle{Risque résiduel}~: niveau qui subsiste après traitement. Le risque
résiduel n'est jamais nul~: il est accepté, explicitement, par la direction.
\end{boxdef}

Distinguer les deux évite deux illusions~: croire qu'on part de zéro (il y a
toujours un risque inhérent) et croire qu'on arrive à zéro (il reste toujours
un résiduel). Le management consiste à ramener le résiduel sous un seuil
accepté.

% ============================================================================
\section{Appétence et Acceptation}
\label{sec:c11-appetence}
% ============================================================================

\begin{boxdef}
\termecle{Appétence au risque}~: niveau de risque que l'organisation est prête
à accepter. Les \termecle{critères d'acceptation} traduisent cette appétence en
seuils concrets.
\end{boxdef}

Pour une mairie, l'appétence est faible sur l'état civil (qui fait foi) et plus
élevée sur le site web d'information. Fixer l'appétence est une décision de
gouvernance (le maire), pas un choix technique.

\begin{boxmetier}{RSSI}
\textbf{Chiffrer sans faux-semblant.} La RSSI sait que « Risque = P × I » ne
produit pas un nombre scientifique, mais un classement défendable. L'important
n'est pas la précision du chiffre~: c'est la cohérence du classement et la
traçabilité du raisonnement.
\end{boxmetier}

% ============================================================================
\section{Un Cas Appliqué}
\label{sec:c11-cas}
% ============================================================================

\begin{boxobjet}{Le registre d'état civil}
\textbf{Constatation.} Actif~: registre d'état civil numérisé. Menace~:
rançongiciel apporté par clé \texttt{USB}. Vulnérabilité~: pas de sauvegarde,
accès partagé. Impact~: service à l'arrêt, actes indisponibles.

\medskip
\textbf{Analyse} (confiance élevée). Probabilité élevée (la menace s'est déjà
réalisée), impact majeur~: risque inhérent très élevé. Deux mesures,
la sauvegarde quotidienne et l'interdiction des clés personnelles, réduisent la
probabilité et l'impact, abaissant le risque résiduel à un niveau acceptable.
\end{boxobjet}

\separateur

\begin{boxresume}
\textbf{Ce qu'il faut retenir de ce chapitre~:}
\begin{itemize}
    \item Un risque enchaîne \textbf{actif, menace, vulnérabilité,
          probabilité, impact}.
    \item Les \textbf{formules} du risque sont des grilles de raisonnement, non
          des calculs exacts.
    \item Le \textbf{risque résiduel} subsiste toujours et s'accepte
          explicitement.
    \item L'\textbf{appétence} au risque est une décision de gouvernance.
    \item Le vocabulaire s'applique à chaque \textbf{actif} de la mairie.
\end{itemize}
\end{boxresume}

\begin{boxinfo}
\textbf{Points de contrôle de l'élève.} Je sais~: définir les cinq notions~;
énoncer les formules du risque~; distinguer inhérent et résiduel~; expliquer
l'appétence~; analyser un cas.
\end{boxinfo}

\bigskip

\begin{boxexo}[title={\ding{45}\quad Exercice~11.1. Les cinq notions \niveaubase}]
Définissez actif, menace, vulnérabilité, impact et probabilité.
\end{boxexo}

\begin{boxexo}[title={\ding{45}\quad Exercice~11.2. La chaîne \niveaubase}]
Pour le vol du fichier des contribuables, nommez la menace, la vulnérabilité et
l'impact.
\end{boxexo}

\begin{boxexo}[title={\ding{45}\quad Exercice~11.3. Inhérent ou résiduel \niveaubase}]
Un risque « avant toute mesure » est-il inhérent ou résiduel~?
\end{boxexo}

\begin{boxexo}[title={\ding{45}\quad Exercice~11.4. Formule \niveaumoyen}]
Pourquoi un impact nul annule-t-il le risque, même si la probabilité est
élevée~?
\end{boxexo}

\begin{boxexo}[title={\ding{45}\quad Exercice~11.5. Résiduel jamais nul \niveaumoyen}]
Expliquez pourquoi le risque résiduel n'est jamais nul.
\end{boxexo}

\begin{boxexo}[title={\ding{45}\quad Exercice~11.6. Appétence \niveaumoyen}]
Pourquoi l'appétence au risque de la mairie diffère-t-elle entre l'état civil
et le site web~?
\end{boxexo}

\begin{boxexo}[title={\ding{45}\quad Exercice~11.7. Casser un maillon \niveaumoyen}]
Montrez qu'agir sur la vulnérabilité réduit le risque, sur l'exemple de la
mairie.
\end{boxexo}

\begin{boxexo}[title={\ding{45}\quad Exercice~11.8. Chiffre défendable \niveaumoyen}]
« Risque = P × I donne un nombre scientifique. » Nuancez.
\end{boxexo}

\begin{boxexo}[title={\ding{45}\quad Exercice~11.9. Analyse d'actif \niveauexpert}]
Analysez le risque sur la messagerie de la mairie (actif, menace,
vulnérabilité, probabilité, impact, inhérent, résiduel).
\end{boxexo}

\begin{boxexo}[title={\ding{45}\quad Exercice~11.10. Critères d'acceptation \niveauexpert}]
Proposez des critères d'acceptation du risque pour la mairie d'Eséka.
\end{boxexo}

\bigskip

\begin{boxinfo}
\textbf{Pour aller plus loin.}
\begin{itemize}
    \item ISO/IEC 27005~\cite{iso27005}~; ISO 31000~\cite{iso31000}.
\end{itemize}
\end{boxinfo}

\begin{boxinfo}
\textbf{Transition.} Le vocabulaire posé, le Chapitre~\ref{chap:methodes}
présente les méthodes qui l'appliquent~: EBIOS RM, OCTAVE, MEHARI.
\end{boxinfo}
}{}
        \IfFileExists{partie4/P04_C12_methodes_risque.tex}{% ============================================================================
% Fichier  : partie4/P04_C12_methodes_risque.tex
% Titre    : Les Méthodes d'Appréciation des Risques
% Auteur   : MINKA MI NGUIDJOÏ Thierry Emmanuel
% ============================================================================
% Sources : EBIOS Risk Manager (ANSSI, 2018) ; OCTAVE ; MEHARI ; ISO 27005.
% Mairie d'Eséka : cas fictif.
% ============================================================================

\chapter{Les Méthodes d'Appréciation des Risques}
\label{chap:methodes}
\chaptermark{Les méthodes de risque}

\epigraph{%
    « Une méthode ne supprime pas le jugement~:\\
    elle le rend traçable et partageable. »%
}{Principe du cours}

\niveauChapitre{Intermédiaire}{Le Chapitre~\ref{chap:vocabrisque}.}

\begin{boxobjectifs}
\begin{itemize}
    \item Présenter \sigle{EBIOS} Risk Manager et ses cinq ateliers.
    \item Décrire la méthode \sigle{OCTAVE} et ses trois phases.
    \item Situer \sigle{MEHARI} et son approche par bases de connaissances.
    \item Comparer les méthodes et choisir selon le contexte.
    \item Appliquer une démarche d'appréciation à la mairie.
\end{itemize}
\end{boxobjectifs}

\bigskip

% ============================================================================
\section*{Question directrice : comment conduire une appréciation des risques sans se perdre ?}
% ============================================================================

Apprécier les risques « à l'intuition » produit des oublis et des injustices.
Les méthodes structurent la démarche. Ce chapitre présente les principales et
aide à choisir.

\medskip
\textbf{Plan du chapitre.}
\begin{enumerate}
    \item EBIOS Risk Manager (section~\ref{sec:c12-ebios}).
    \item OCTAVE (section~\ref{sec:c12-octave}).
    \item MEHARI (section~\ref{sec:c12-mehari}).
    \item Comparer et choisir (section~\ref{sec:c12-choisir}).
    \item Appliquer à la mairie (section~\ref{sec:c12-mairie}).
\end{enumerate}

\separateur

% ============================================================================
\section{EBIOS Risk Manager}
\label{sec:c12-ebios}
% ============================================================================

\begin{boxdef}
\termecle{\sigle{EBIOS} Risk Manager}~: méthode française d'appréciation et de
traitement des risques numériques, publiée par l'\sigle{ANSSI}, organisée en
cinq ateliers.
\end{boxdef}

\begin{tableaucours}{L{1.4cm}L{4.6cm}L{6.2cm}}{%
    \tableauHeader{Atelier} &
    \tableauHeader{Nom} &
    \tableauHeader{Produit}
}
1 & Cadrage et socle & Périmètre, valeurs métier, socle de sécurité \\
2 & Sources de risque & Qui menace, pour quels objectifs \\
3 & Scénarios stratégiques & Chemins d'attaque de haut niveau \\
4 & Scénarios opérationnels & Chemins techniques détaillés \\
5 & Traitement du risque & Mesures, risques résiduels \\
\end{tableaucours}
\vspace{-0.8em}
\begin{center}{\small\itshape Tableau~12.1. Les cinq ateliers d'EBIOS Risk Manager}\end{center}

\sigle{EBIOS} \sigle{RM} met l'accent sur les \emph{sources de risque} et les
\emph{scénarios}~: elle part de qui veut nuire et pourquoi, plutôt que d'une
liste de vulnérabilités.

% ============================================================================
\section{OCTAVE}
\label{sec:c12-octave}
% ============================================================================

\begin{boxdef}
\termecle{\sigle{OCTAVE}} (\emph{Operationally Critical Threat, Asset and
Vulnerability Evaluation})~: méthode centrée sur les actifs critiques,
organisée en trois phases.
\end{boxdef}

\begin{enumerate}
    \item \textbf{Profils de menaces fondés sur les actifs}~: identifier les
          actifs critiques et les menaces qui les visent.
    \item \textbf{Vulnérabilités de l'infrastructure}~: analyser les
          faiblesses techniques.
    \item \textbf{Stratégie et plan de sécurité}~: évaluer impact et
          probabilité, bâtir le plan.
\end{enumerate}

\sigle{OCTAVE} est conduite par l'organisation elle-même, ce qui favorise
l'appropriation.

% ============================================================================
\section{MEHARI}
\label{sec:c12-mehari}
% ============================================================================

\begin{boxdef}
\termecle{\sigle{MEHARI}} (Méthode Harmonisée d'Analyse de Risques)~: méthode
du \sigle{CLUSIF} s'appuyant sur des bases de connaissances pour évaluer
l'exposition et l'efficacité des mesures.
\end{boxdef}

\sigle{MEHARI} propose une analyse fine en situation normale, en cas de
dysfonctionnement majeur, et par la vulnérabilité. Son atout est la richesse de
ses bases~; son coût, leur lourdeur pour une petite structure.

% ============================================================================
\section{Comparer et Choisir}
\label{sec:c12-choisir}
% ============================================================================

\begin{tableaucours}{L{2.6cm}L{4.4cm}L{5.2cm}}{%
    \tableauHeader{Méthode} &
    \tableauHeader{Force} &
    \tableauHeader{Convient à}
}
EBIOS RM & Scénarios, sources de risque & Projets structurés, lien ISO 27005 \\
OCTAVE & Appropriation interne & Organisations qui s'auto-évaluent \\
MEHARI & Bases de connaissances & Structures matures et outillées \\
\end{tableaucours}
\vspace{-0.8em}
\begin{center}{\small\itshape Tableau~12.2. Choisir une méthode selon le contexte}\end{center}

Aucune méthode n'est « la meilleure »~: la bonne est celle que l'organisation
peut réellement conduire. Toutes sont compatibles avec \norme{ISO/IEC 27005}.

% ============================================================================
\section{Appliquer à la Mairie}
\label{sec:c12-mairie}
% ============================================================================

\begin{boxprocedure}
\textbf{Constatation.} La mairie d'Eséka, aux moyens limités, choisit
\sigle{EBIOS} \sigle{RM} pour l'opération \textsc{Kellé}, en version allégée.

\medskip
\textbf{Analyse.} Le choix est cohérent~: \sigle{EBIOS} \sigle{RM} est
publique, gratuite, outillée, et son approche par scénarios parle aux élus (qui
veut nuire à la mairie, et comment). Une version allégée évite la lourdeur pour
une petite structure.
\end{boxprocedure}

\begin{boxmetier}{RSSI}
\textbf{Adapter, pas appliquer à la lettre.} La RSSI adapte la méthode à la
taille de la mairie~: cinq ateliers menés en une journée d'atelier collectif,
plutôt qu'une étude de plusieurs mois. La méthode sert le projet, pas
l'inverse.
\end{boxmetier}

\separateur

\begin{boxresume}
\textbf{Ce qu'il faut retenir de ce chapitre~:}
\begin{itemize}
    \item \textbf{EBIOS RM} (ANSSI) structure l'appréciation en cinq ateliers,
          centrés sur les scénarios.
    \item \textbf{OCTAVE} part des actifs critiques en trois phases, par
          auto-évaluation.
    \item \textbf{MEHARI} s'appuie sur des bases de connaissances, pour
          structures matures.
    \item On \textbf{choisit} la méthode que l'organisation peut conduire~;
          toutes sont compatibles ISO 27005.
    \item La mairie adopte \textbf{EBIOS RM allégée} pour l'opération
          \textsc{Kellé}.
\end{itemize}
\end{boxresume}

\begin{boxinfo}
\textbf{Points de contrôle de l'élève.} Je sais~: citer les cinq ateliers
d'EBIOS RM~; décrire les trois phases d'OCTAVE~; situer MEHARI~; comparer les
méthodes~; justifier un choix.
\end{boxinfo}

\bigskip

\begin{boxexo}[title={\ding{45}\quad Exercice~12.1. Les ateliers \niveaubase}]
Citez les cinq ateliers d'\sigle{EBIOS} Risk Manager.
\end{boxexo}

\begin{boxexo}[title={\ding{45}\quad Exercice~12.2. Les phases \niveaubase}]
Énoncez les trois phases d'\sigle{OCTAVE}.
\end{boxexo}

\begin{boxexo}[title={\ding{45}\quad Exercice~12.3. MEHARI \niveaubase}]
Quel est le principe de \sigle{MEHARI} et son principal atout~?
\end{boxexo}

\begin{boxexo}[title={\ding{45}\quad Exercice~12.4. Source de risque \niveaumoyen}]
En quoi EBIOS RM se distingue-t-elle en partant des sources de risque~?
\end{boxexo}

\begin{boxexo}[title={\ding{45}\quad Exercice~12.5. Choisir \niveaumoyen}]
Pour une petite mairie, quelle méthode recommandez-vous~? Justifiez.
\end{boxexo}

\begin{boxexo}[title={\ding{45}\quad Exercice~12.6. Adapter \niveaumoyen}]
Pourquoi faut-il adapter la méthode à la taille de l'organisation~?
\end{boxexo}

\begin{boxexo}[title={\ding{45}\quad Exercice~12.7. Atelier 1 \niveaumoyen}]
Que produit l'atelier~1 d'EBIOS RM pour la mairie d'Eséka~?
\end{boxexo}

\begin{boxexo}[title={\ding{45}\quad Exercice~12.8. Compatibilité ISO \niveaumoyen}]
Pourquoi ces méthodes sont-elles toutes compatibles avec \norme{ISO/IEC 27005}~?
\end{boxexo}

\begin{boxexo}[title={\ding{45}\quad Exercice~12.9. Mini-EBIOS \niveauexpert}]
Conduisez un mini-atelier EBIOS RM sur l'état civil~: deux sources de risque,
deux scénarios stratégiques.
\end{boxexo}

\begin{boxexo}[title={\ding{45}\quad Exercice~12.10. Note de méthode \niveauexpert}]
Rédigez la note qui justifie, auprès du maire, le choix de méthode pour
l'opération \textsc{Kellé}.
\end{boxexo}

\bigskip

\begin{boxinfo}
\textbf{Pour aller plus loin.}
\begin{itemize}
    \item EBIOS Risk Manager~\cite{ebios2018}~; ISO/IEC 27005~\cite{iso27005}.
\end{itemize}
\end{boxinfo}

\begin{boxinfo}
\textbf{Transition.} Les risques appréciés, il faut les cartographier et
décider~: le Chapitre~\ref{chap:traitement} traite la cartographie et le
traitement.
\end{boxinfo}
}{}
        \IfFileExists{partie4/P04_C13_cartographie_traitement.tex}{% ============================================================================
% Fichier  : partie4/P04_C13_cartographie_traitement.tex
% Titre    : Cartographie et Traitement des Risques
% Auteur   : MINKA MI NGUIDJOÏ Thierry Emmanuel
% ============================================================================
% Sources : ISO/IEC 27005:2022 ; ISO 31000:2018. Mairie d'Eséka : cas fictif.
% ============================================================================

\chapter{Cartographie et Traitement des Risques}
\label{chap:traitement}
\chaptermark{Cartographie et traitement}

\epigraph{%
    « Cartographier, c'est transformer une inquiétude diffuse\\
    en une liste de décisions à prendre. »%
}{Principe du cours}

\niveauChapitre{Intermédiaire}{Les Chapitres~\ref{chap:vocabrisque} et~\ref{chap:methodes}.}

\begin{boxobjectifs}
\begin{itemize}
    \item Construire et lire une matrice des risques.
    \item Hiérarchiser les risques par niveau.
    \item Énoncer les quatre stratégies de traitement.
    \item Choisir une stratégie selon le risque et le coût.
    \item Formaliser le risque résiduel et son acceptation.
\end{itemize}
\end{boxobjectifs}

\bigskip

% ============================================================================
\section*{Question directrice : une fois les risques connus, que décide-t-on ?}
% ============================================================================

Connaître ses risques ne sert à rien si l'on n'en fait rien. Ce chapitre
transforme la liste des risques en décisions~: les classer, puis, pour chacun,
choisir un traitement.

\medskip
\textbf{Plan du chapitre.}
\begin{enumerate}
    \item La matrice des risques (section~\ref{sec:c13-matrice}).
    \item Hiérarchiser (section~\ref{sec:c13-hierarchie}).
    \item Les quatre stratégies (section~\ref{sec:c13-strategies}).
    \item Choisir un traitement (section~\ref{sec:c13-choisir}).
    \item Le risque résiduel (section~\ref{sec:c13-residuel}).
\end{enumerate}

\separateur

% ============================================================================
\section{La Matrice des Risques}
\label{sec:c13-matrice}
% ============================================================================

\begin{boxdef}
La \termecle{matrice des risques} croise la \emph{probabilité} et l'\emph{impact}
pour situer chaque risque dans une grille et en déduire un niveau.
\end{boxdef}

\begin{tableaucours}{L{2.4cm}L{2.0cm}L{2.0cm}L{2.0cm}L{2.0cm}}{%
    \tableauHeader{Impact / Prob.} &
    \tableauHeader{Faible} &
    \tableauHeader{Moyenne} &
    \tableauHeader{Élevée} &
    \tableauHeader{Très él.}
}
Mineur & Faible & Faible & Moyen & Moyen \\
Modéré & Faible & Moyen & Moyen & Élevé \\
Majeur & Moyen & Moyen & Élevé & Très él. \\
Critique & Moyen & Élevé & Très él. & Très él. \\
\end{tableaucours}
\vspace{-0.8em}
\begin{center}{\small\itshape Tableau~13.1. Matrice probabilité $\times$ impact}\end{center}

% ============================================================================
\section{Hiérarchiser}
\label{sec:c13-hierarchie}
% ============================================================================

La matrice hiérarchise~: on traite d'abord les risques « très élevés » et
« élevés », on surveille les « moyens », on documente l'acceptation des
« faibles ». Cette hiérarchie fonde le \termecle{plan de traitement}~: on ne
peut pas tout traiter en même temps, donc on priorise.

\begin{boxprocedure}
\textbf{Constatation.} À Eséka, le risque « perte de l'état civil » est très
élevé~; le risque « défiguration du site web » est moyen.

\medskip
\textbf{Analyse.} La priorité va à l'état civil~: impact critique, probabilité
démontrée. Le site web, moins critique, sera traité ensuite. Hiérarchiser,
c'est choisir où mettre ses moyens limités.
\end{boxprocedure}

% ============================================================================
\section{Les Quatre Stratégies}
\label{sec:c13-strategies}
% ============================================================================

\begin{boxdef}
Pour chaque risque, quatre \termecle{stratégies de traitement}~:
\begin{itemize}
    \item \termecle{Éviter}~: supprimer la cause (renoncer à l'activité
          risquée).
    \item \termecle{Réduire}~: diminuer la probabilité ou l'impact par des
          mesures.
    \item \termecle{Transférer}~: reporter le risque (assurance,
          sous-traitance).
    \item \termecle{Accepter}~: assumer le risque résiduel en connaissance de
          cause.
\end{itemize}
\end{boxdef}

% ============================================================================
\section{Choisir un Traitement}
\label{sec:c13-choisir}
% ============================================================================

\begin{tableaucours}{L{4.2cm}L{3.0cm}L{4.4cm}}{%
    \tableauHeader{Risque} &
    \tableauHeader{Stratégie} &
    \tableauHeader{Mesure}
}
Perte de l'état civil & Réduire & Sauvegarde 3-2-1, test de restauration \\
Fuite du fichier taxes & Réduire & Chiffrement, accès nominatif \\
Panne électrique & Transférer & Contrat de maintenance, onduleur loué \\
Site web décoratif & Accepter & Risque résiduel faible, documenté \\
\end{tableaucours}
\vspace{-0.8em}
\begin{center}{\small\itshape Tableau~13.2. Choix de traitement à la mairie}\end{center}

Le choix met en balance le \emph{coût} du traitement et le \emph{gain} de
réduction du risque~: réduire un risque faible à grands frais est un mauvais
arbitrage.

\begin{boxmetier}{Direction générale}
\textbf{Accepter est une décision, pas un oubli.} Accepter un risque est
légitime, à condition que ce soit explicite, documenté et décidé par la
direction. Un risque « accepté » par simple inaction n'est pas accepté~: il est
ignoré.
\end{boxmetier}

% ============================================================================
\section{Le Risque Résiduel}
\label{sec:c13-residuel}
% ============================================================================

Après traitement, il reste un \termecle{risque résiduel}. La norme exige que la
direction l'\termecle{accepte formellement} (Chapitre~\ref{chap:dda}). Cette
acceptation est tracée~: c'est la direction qui porte, en connaissance de
cause, ce qui n'a pas été éliminé.

\separateur

\begin{boxresume}
\textbf{Ce qu'il faut retenir de ce chapitre~:}
\begin{itemize}
    \item La \textbf{matrice} situe chaque risque et en déduit un niveau.
    \item On \textbf{hiérarchise} pour traiter d'abord les risques les plus
          forts.
    \item Quatre \textbf{stratégies}~: éviter, réduire, transférer, accepter.
    \item On \textbf{choisit} en pesant coût et gain.
    \item Le \textbf{risque résiduel} est accepté formellement par la
          direction.
\end{itemize}
\end{boxresume}

\begin{boxinfo}
\textbf{Points de contrôle de l'élève.} Je sais~: lire une matrice des
risques~; hiérarchiser~; énoncer les quatre stratégies~; choisir un
traitement~; expliquer l'acceptation du risque résiduel.
\end{boxinfo}

\bigskip

\begin{boxexo}[title={\ding{45}\quad Exercice~13.1. Lire la matrice \niveaubase}]
Un risque de probabilité élevée et d'impact critique~: quel niveau~?
\end{boxexo}

\begin{boxexo}[title={\ding{45}\quad Exercice~13.2. Les stratégies \niveaubase}]
Citez les quatre stratégies de traitement du risque.
\end{boxexo}

\begin{boxexo}[title={\ding{45}\quad Exercice~13.3. Hiérarchiser \niveaubase}]
Entre un risque « très élevé » et un « moyen », lequel traite-t-on d'abord~?
\end{boxexo}

\begin{boxexo}[title={\ding{45}\quad Exercice~13.4. Associer \niveaumoyen}]
Associez une stratégie à chacun~: perte de l'état civil~; panne électrique~;
site web décoratif.
\end{boxexo}

\begin{boxexo}[title={\ding{45}\quad Exercice~13.5. Coût-gain \niveaumoyen}]
Pourquoi réduire un risque faible à grands frais est-il un mauvais arbitrage~?
\end{boxexo}

\begin{boxexo}[title={\ding{45}\quad Exercice~13.6. Transférer \niveaumoyen}]
Donnez un exemple de risque que la mairie peut transférer, et comment.
\end{boxexo}

\begin{boxexo}[title={\ding{45}\quad Exercice~13.7. Accepter ou ignorer \niveaumoyen}]
Quelle différence entre un risque accepté et un risque ignoré~?
\end{boxexo}

\begin{boxexo}[title={\ding{45}\quad Exercice~13.8. Résiduel \niveaumoyen}]
Qui accepte le risque résiduel, et pourquoi doit-ce être tracé~?
\end{boxexo}

\begin{boxexo}[title={\ding{45}\quad Exercice~13.9. Cartographie \niveauexpert}]
Dressez la cartographie de cinq risques de la mairie dans la matrice, avec leur
niveau.
\end{boxexo}

\begin{boxexo}[title={\ding{45}\quad Exercice~13.10. Plan de traitement \niveauexpert}]
Construisez un plan de traitement pour ces cinq risques (stratégie, mesure,
responsable, échéance).
\end{boxexo}

\bigskip

\begin{boxinfo}
\textbf{Pour aller plus loin.}
\begin{itemize}
    \item ISO/IEC 27005~\cite{iso27005}~; ISO 31000~\cite{iso31000}.
\end{itemize}
\end{boxinfo}

\begin{boxinfo}
\textbf{Transition.} Le plan de traitement se formalise dans un document
central~: le Chapitre~\ref{chap:dda} présente la déclaration d'applicabilité.
\end{boxinfo}
}{}
        \IfFileExists{partie4/P04_C14_declaration_applicabilite.tex}{% ============================================================================
% Fichier  : partie4/P04_C14_declaration_applicabilite.tex
% Titre    : La Déclaration d'Applicabilité et le Plan de Traitement
% Auteur   : MINKA MI NGUIDJOÏ Thierry Emmanuel
% ============================================================================
% Sources : ISO/IEC 27001:2022 (clause 6.1.3) ; ISO/IEC 27005. Fictif : Eséka.
% ============================================================================

\chapter{La Déclaration d'Applicabilité et le Plan de Traitement}
\label{chap:dda}
\chaptermark{La déclaration d'applicabilité}

\epigraph{%
    « La déclaration d'applicabilité est la carte d'identité du SMSI~:\\
    elle dit ce qu'on fait, ce qu'on ne fait pas, et pourquoi. »%
}{Principe du cours}

\niveauChapitre{Intermédiaire}{Le Chapitre~\ref{chap:traitement}.}

\begin{boxobjectifs}
\begin{itemize}
    \item Définir la déclaration d'applicabilité (\sigle{DdA}) et son rôle.
    \item Construire une \sigle{DdA} à partir des risques et de l'Annexe~A.
    \item Distinguer \sigle{DdA} et plan de traitement des risques.
    \item Justifier l'inclusion et l'exclusion d'une mesure.
    \item Faire acter le risque résiduel par la direction.
\end{itemize}
\end{boxobjectifs}

\bigskip

% ============================================================================
\section*{Question directrice : comment prouver que les mesures retenues sont les bonnes ?}
% ============================================================================

L'auditeur ne croit pas sur parole~: il demande à voir le lien entre risques et
mesures. La déclaration d'applicabilité est le document qui établit ce lien.
Ce chapitre la construit.

\medskip
\textbf{Plan du chapitre.}
\begin{enumerate}
    \item Rôle de la DdA (section~\ref{sec:c14-role}).
    \item Construire la DdA (section~\ref{sec:c14-construire}).
    \item DdA et plan de traitement (section~\ref{sec:c14-plan}).
    \item Inclure, exclure (section~\ref{sec:c14-justifier}).
    \item Acter le risque résiduel (section~\ref{sec:c14-acter}).
\end{enumerate}

\separateur

% ============================================================================
\section{Rôle de la DdA}
\label{sec:c14-role}
% ============================================================================

\begin{boxdef}
La \termecle{déclaration d'applicabilité} (\sigle{DdA})~: document exigé par la
clause~6.1.3 de \norme{ISO/IEC 27001} qui, pour chaque mesure de l'Annexe~A,
indique si elle est \emph{retenue}, sa \emph{justification} (inclusion ou
exclusion) et son \emph{état} de mise en \oe uvre.
\end{boxdef}

La \sigle{DdA} est le document que tout auditeur demande en premier~: elle
résume le \sigle{SMSI} en une page de synthèse et une grille. Elle fait le pont
entre la partie IV (risques) et la partie V (mesures).

% ============================================================================
\section{Construire la DdA}
\label{sec:c14-construire}
% ============================================================================

\begin{tableaucours}{L{3.4cm}L{1.6cm}L{3.2cm}L{3.0cm}}{%
    \tableauHeader{Mesure (Annexe A)} &
    \tableauHeader{Retenue} &
    \tableauHeader{Justification} &
    \tableauHeader{État}
}
Sauvegarde & Oui & Risque R1 (état civil) & En place \\
Contrôle d'accès & Oui & Risque R2 (taxes) & En place \\
Chiffrement & Oui & Risque R2 & En cours \\
Développement sécurisé & Non & Pas de développement interne & Exclue \\
\end{tableaucours}
\vspace{-0.8em}
\begin{center}{\small\itshape Tableau~14.1. Extrait de la DdA de la mairie d'Eséka}\end{center}

Chaque ligne relie une mesure à un risque (inclusion) ou explique pourquoi la
mesure ne s'applique pas (exclusion). Rien n'est laissé implicite.

% ============================================================================
\section{DdA et Plan de Traitement}
\label{sec:c14-plan}
% ============================================================================

\begin{boxdef}
Le \termecle{plan de traitement des risques}~: document qui, pour chaque risque,
précise la stratégie, les mesures, le responsable et l'échéance. La \sigle{DdA}
dit \emph{quelles} mesures~; le plan dit \emph{comment et quand} les mettre en
\oe uvre.
\end{boxdef}

Les deux documents sont complémentaires~: la \sigle{DdA} est la photographie
des mesures~; le plan est la feuille de route. Ensemble, ils constituent le
« Plan » du \sigle{PDCA}.

% ============================================================================
\section{Inclure, Exclure}
\label{sec:c14-justifier}
% ============================================================================

\begin{boxprocedure}
\textbf{Constatation.} La mairie exclut la mesure « développement sécurisé » au
motif qu'elle n'écrit aucun logiciel elle-même.

\medskip
\textbf{Analyse.} L'exclusion est légitime et justifiée~: la mesure ne
correspond à aucun risque de la mairie. Mais la justification doit être écrite~:
une exclusion non justifiée est un écart d'audit. On ne retire une mesure que
si le risque associé n'existe pas.
\end{boxprocedure}

\begin{boxmetier}{Auditeur SMSI}
\textbf{Lire la DdA à rebours.} L'auditeur part de la \sigle{DdA} et remonte~:
cette mesure retenue répond-elle à un risque réel~? cette mesure exclue l'est-
elle à bon droit~? La \sigle{DdA} est le fil qui relie tout le \sigle{SMSI}.
\end{boxmetier}

% ============================================================================
\section{Acter le Risque Résiduel}
\label{sec:c14-acter}
% ============================================================================

Une fois la \sigle{DdA} et le plan établis, il reste un risque résiduel. La
clause~6.1 exige que les \termecle{propriétaires des risques} approuvent le plan
de traitement et \termecle{acceptent} le risque résiduel. À la mairie, c'est le
maire qui signe cette acceptation.

\begin{boxjuris}
Faire acter le risque résiduel par la direction n'est pas une formalité~: c'est
l'acte par lequel la responsabilité est assumée au bon niveau. En cas
d'incident, cette trace montre que le risque avait été identifié, pesé et
décidé, et non ignoré.
\end{boxjuris}

\separateur

\begin{boxresume}
\textbf{Ce qu'il faut retenir de ce chapitre~:}
\begin{itemize}
    \item La \textbf{DdA} indique, pour chaque mesure, si elle est retenue, sa
          justification et son état.
    \item Elle se \textbf{construit} à partir des risques et de l'Annexe~A.
    \item La \textbf{DdA} dit quelles mesures~; le \textbf{plan de traitement}
          dit comment et quand.
    \item \textbf{Inclusion} et \textbf{exclusion} se justifient par écrit.
    \item La direction \textbf{accepte} formellement le risque résiduel.
\end{itemize}
\end{boxresume}

\begin{boxinfo}
\textbf{Points de contrôle de l'élève.} Je sais~: définir la \sigle{DdA}~; la
construire~; la distinguer du plan de traitement~; justifier une inclusion ou
une exclusion~; expliquer l'acceptation du risque résiduel.
\end{boxinfo}

\bigskip

\begin{boxexo}[title={\ding{45}\quad Exercice~14.1. Rôle de la DdA \niveaubase}]
Qu'indique la \sigle{DdA} pour chaque mesure de l'Annexe~A~?
\end{boxexo}

\begin{boxexo}[title={\ding{45}\quad Exercice~14.2. DdA ou plan \niveaubase}]
Laquelle dit « quelles mesures », lequel dit « comment et quand »~?
\end{boxexo}

\begin{boxexo}[title={\ding{45}\quad Exercice~14.3. Exclusion \niveaubase}]
Pourquoi la mairie peut-elle exclure la mesure « développement sécurisé »~?
\end{boxexo}

\begin{boxexo}[title={\ding{45}\quad Exercice~14.4. Une ligne \niveaumoyen}]
Rédigez une ligne de \sigle{DdA} pour la mesure « journalisation ».
\end{boxexo}

\begin{boxexo}[title={\ding{45}\quad Exercice~14.5. Exclusion non justifiée \niveaumoyen}]
Pourquoi une exclusion non justifiée est-elle un écart d'audit~?
\end{boxexo}

\begin{boxexo}[title={\ding{45}\quad Exercice~14.6. Lire à rebours \niveaumoyen}]
Comment un auditeur se sert-il de la \sigle{DdA} pour vérifier le \sigle{SMSI}~?
\end{boxexo}

\begin{boxexo}[title={\ding{45}\quad Exercice~14.7. Acceptation \niveaumoyen}]
Qui, à la mairie, accepte le risque résiduel, et pourquoi la trace importe-t-
elle~?
\end{boxexo}

\begin{boxexo}[title={\ding{45}\quad Exercice~14.8. Complémentarité \niveaumoyen}]
Montrez que DdA et plan de traitement sont complémentaires.
\end{boxexo}

\begin{boxexo}[title={\ding{45}\quad Exercice~14.9. DdA de la mairie \niveauexpert}]
Construisez un extrait de \sigle{DdA} de huit mesures pour la mairie d'Eséka.
\end{boxexo}

\begin{boxexo}[title={\ding{45}\quad Exercice~14.10. Plan complet \niveauexpert}]
Rédigez le plan de traitement des risques de la mairie (cinq risques, stratégie,
mesures, responsable, échéance, acceptation).
\end{boxexo}

\bigskip

\begin{boxinfo}
\textbf{Pour aller plus loin.}
\begin{itemize}
    \item ISO/IEC 27001, clause~6.1.3~\cite{iso27001}~; ISO/IEC 27005~\cite{iso27005}.
\end{itemize}
\end{boxinfo}

\begin{boxinfo}
\textbf{Transition.} Le plan décidé, place à la mise en \oe uvre~: la partie V
s'ouvre au Chapitre~\ref{chap:orga} sur les mesures organisationnelles et
humaines.
\end{boxinfo}
}{}

    \part{Mettre en Œuvre les Mesures}
    \label{part:mesures}

        \IfFileExists{partie5/P05_C15_mesures_organisationnelles.tex}{% ============================================================================
% Fichier  : partie5/P05_C15_mesures_organisationnelles.tex
% Titre    : Mesures Organisationnelles et Humaines
% Auteur   : MINKA MI NGUIDJOÏ Thierry Emmanuel
% ============================================================================
% Sources : ISO/IEC 27002:2022 (thèmes organisationnel et personnes). Fictif.
% ============================================================================

\chapter{Mesures Organisationnelles et Humaines}
\label{chap:orga}
\chaptermark{Mesures organisationnelles}

\epigraph{%
    « La technique protège les machines~;\\
    la sensibilisation protège les décisions des gens. »%
}{Principe du cours}

\niveauChapitre{Intermédiaire}{La partie IV.}

\begin{boxobjectifs}
\begin{itemize}
    \item Identifier les mesures organisationnelles clés.
    \item Comprendre le poids du facteur humain dans la sécurité.
    \item Concevoir une campagne de sensibilisation.
    \item Sécuriser le cycle de vie de l'agent (arrivée, mobilité, départ).
    \item Encadrer les fournisseurs et le télétravail.
\end{itemize}
\end{boxobjectifs}

\bigskip

% ============================================================================
\section*{Question directrice : comment agir sur ce qui n'est ni un logiciel ni un câble ?}
% ============================================================================

L'incident de la mairie est né d'un geste humain~: brancher une clé. Les
mesures techniques ne suffisent pas~: il faut agir sur l'organisation et sur les
personnes. Ce chapitre traite ces mesures.

\medskip
\textbf{Plan du chapitre.}
\begin{enumerate}
    \item Les mesures organisationnelles (section~\ref{sec:c15-orga}).
    \item Le facteur humain (section~\ref{sec:c15-humain}).
    \item La sensibilisation (section~\ref{sec:c15-sensibilisation}).
    \item Le cycle de vie de l'agent (section~\ref{sec:c15-cyclevie}).
    \item Fournisseurs et télétravail (section~\ref{sec:c15-tiers}).
\end{enumerate}

\separateur

% ============================================================================
\section{Les Mesures Organisationnelles}
\label{sec:c15-orga}
% ============================================================================

\begin{tableaucours}{L{4.2cm}L{8.0cm}}{%
    \tableauHeader{Mesure} &
    \tableauHeader{Ce qu'elle installe}
}
Politiques et procédures & Le cadre écrit des règles \\
Rôles et responsabilités & Qui fait quoi (matrice RACI) \\
Classification de l'information & Des niveaux et des règles de manipulation \\
Gestion des incidents & Un processus de signalement et de traitement \\
Relations avec les fournisseurs & Des exigences de sécurité contractuelles \\
\end{tableaucours}
\vspace{-0.8em}
\begin{center}{\small\itshape Tableau~15.1. Mesures organisationnelles (ISO/IEC 27002)}\end{center}

Ces mesures ne coûtent souvent que du temps et de la rigueur~: elles sont les
premières à déployer, avant tout investissement technique.

% ============================================================================
\section{Le Facteur Humain}
\label{sec:c15-humain}
% ============================================================================

\begin{boxdef}
Le \termecle{facteur humain}~: part des incidents liée aux comportements des
personnes (erreur, négligence, manipulation). Il est, dans la plupart des
études, la première cause d'incident.
\end{boxdef}

L'agent n'est pas le maillon faible par nature~: il le devient faute de
formation et de règles claires. Sensibiliser transforme le maillon faible en
première ligne de défense.

\begin{boxscenario}{Opération MAKAK. La clé piégée}
\textbf{Constatation.} L'agent qui a branché la clé ignorait le risque~: il
n'avait reçu ni règle écrite, ni formation.

\medskip
\textbf{Analyse} (confiance élevée). Le reproche ne vise pas l'agent mais
l'organisation~: sans règle ni sensibilisation, le geste était prévisible.
La mesure corrective est double~: interdire les clés personnelles \emph{et}
former les agents.
\end{boxscenario}

% ============================================================================
\section{La Sensibilisation}
\label{sec:c15-sensibilisation}
% ============================================================================

\begin{tableaucours}{L{4.0cm}L{8.2cm}}{%
    \tableauHeader{Thème} &
    \tableauHeader{Message clé}
}
Mots de passe & Individuels, robustes, jamais partagés \\
Hameçonnage & Vérifier l'expéditeur, ne pas cliquer dans le doute \\
Supports amovibles & N'utiliser que des supports maîtrisés \\
Données personnelles & Ne traiter que le nécessaire, protéger \\
Signalement & Prévenir vite en cas de doute, sans crainte \\
\end{tableaucours}
\vspace{-0.8em}
\begin{center}{\small\itshape Tableau~15.2. Thèmes d'une campagne de sensibilisation}\end{center}

Une campagne efficace est régulière, concrète et bienveillante~: on apprend à
signaler sans être puni. Hameçonnage simulé, affiches et courtes sessions
valent mieux qu'un règlement illisible.

\begin{boxmetier}{Responsable conformité}
\textbf{Prouver la sensibilisation.} La formation doit être tracée~: feuilles
d'émargement, attestations, résultats des tests. En audit, « nous sensibilisons
nos agents » ne vaut rien sans preuve.
\end{boxmetier}

% ============================================================================
\section{Le Cycle de Vie de l'Agent}
\label{sec:c15-cyclevie}
% ============================================================================

La sécurité suit l'agent de son arrivée à son départ~: à l'\termecle{arrivée},
charte signée et accès minimal~; en cas de \termecle{mobilité}, droits
réajustés~; au \termecle{départ}, accès révoqués et matériel restitué. Un
compte d'agent parti et toujours actif est une porte ouverte.

% ============================================================================
\section{Fournisseurs et Télétravail}
\label{sec:c15-tiers}
% ============================================================================

\begin{boxprocedure}
\textbf{Constatation.} L'éditeur du logiciel d'état civil accède à distance aux
données de la mairie, sans clause de sécurité au contrat.

\medskip
\textbf{Analyse.} C'est un risque de tiers~: la mairie reste responsable des
données que son prestataire manipule. La mesure impose des \emph{exigences de
sécurité contractuelles} et un contrôle des accès distants. Le télétravail des
agents suit la même logique~: règles écrites et moyens maîtrisés.
\end{boxprocedure}

\separateur

\begin{boxresume}
\textbf{Ce qu'il faut retenir de ce chapitre~:}
\begin{itemize}
    \item Les \textbf{mesures organisationnelles} posent le cadre écrit, à
          faible coût.
    \item Le \textbf{facteur humain} est la première cause d'incident~; la
          formation en fait une défense.
    \item Une \textbf{campagne de sensibilisation} est régulière, concrète,
          tracée.
    \item La sécurité suit le \textbf{cycle de vie de l'agent}, du recrutement
          au départ.
    \item \textbf{Fournisseurs} et \textbf{télétravail} exigent des règles et
          des contrats.
\end{itemize}
\end{boxresume}

\begin{boxinfo}
\textbf{Points de contrôle de l'élève.} Je sais~: citer des mesures
organisationnelles~; expliquer le facteur humain~; concevoir une campagne~;
sécuriser le cycle de vie de l'agent~; encadrer fournisseurs et télétravail.
\end{boxinfo}

\bigskip

\begin{boxexo}[title={\ding{45}\quad Exercice~15.1. Mesures organisationnelles \niveaubase}]
Citez trois mesures organisationnelles de l'ISO/IEC 27002.
\end{boxexo}

\begin{boxexo}[title={\ding{45}\quad Exercice~15.2. Facteur humain \niveaubase}]
Pourquoi dit-on que l'humain est la première cause d'incident~?
\end{boxexo}

\begin{boxexo}[title={\ding{45}\quad Exercice~15.3. Thèmes \niveaubase}]
Citez quatre thèmes d'une campagne de sensibilisation.
\end{boxexo}

\begin{boxexo}[title={\ding{45}\quad Exercice~15.4. Blâmer qui \niveaumoyen}]
Faut-il blâmer l'agent qui a branché la clé~? Argumentez.
\end{boxexo}

\begin{boxexo}[title={\ding{45}\quad Exercice~15.5. Prouver la formation \niveaumoyen}]
Comment la mairie prouve-t-elle, en audit, qu'elle sensibilise ses agents~?
\end{boxexo}

\begin{boxexo}[title={\ding{45}\quad Exercice~15.6. Départ d'un agent \niveaumoyen}]
Quelles actions de sécurité au départ d'un agent~?
\end{boxexo}

\begin{boxexo}[title={\ding{45}\quad Exercice~15.7. Risque de tiers \niveaumoyen}]
Pourquoi la mairie reste-t-elle responsable des données manipulées par son
prestataire~?
\end{boxexo}

\begin{boxexo}[title={\ding{45}\quad Exercice~15.8. Télétravail \niveaumoyen}]
Citez trois règles de sécurité pour le télétravail d'un agent de la mairie.
\end{boxexo}

\begin{boxexo}[title={\ding{45}\quad Exercice~15.9. Plan de campagne \niveauexpert}]
Concevez un plan de campagne de sensibilisation sur un an pour la mairie
d'Eséka.
\end{boxexo}

\begin{boxexo}[title={\ding{45}\quad Exercice~15.10. Clause fournisseur \niveauexpert}]
Rédigez une clause de sécurité à insérer au contrat de l'éditeur du logiciel
d'état civil.
\end{boxexo}

\bigskip

\begin{boxinfo}
\textbf{Pour aller plus loin.}
\begin{itemize}
    \item ISO/IEC 27002, thèmes organisationnel et personnes~\cite{iso27002}~;
          Guide d'hygiène ANSSI~\cite{anssihygiene2017}.
\end{itemize}
\end{boxinfo}

\begin{boxinfo}
\textbf{Transition.} Les mesures humaines posées, le Chapitre~\ref{chap:technique}
traite les mesures techniques~: accès, réseau, chiffrement, journalisation.
\end{boxinfo}
}{}
        \IfFileExists{partie5/P05_C16_mesures_techniques.tex}{% ============================================================================
% Fichier  : partie5/P05_C16_mesures_techniques.tex
% Titre    : Mesures Techniques
% Auteur   : MINKA MI NGUIDJOÏ Thierry Emmanuel
% ============================================================================
% Sources : ISO/IEC 27002:2022 (thème technologique) ; ANSSI. Fictif : Eséka.
% ============================================================================

\chapter{Mesures Techniques}
\label{chap:technique}
\chaptermark{Mesures techniques}

\epigraph{%
    « Une bonne mesure technique applique une décision de sécurité~;\\
    elle ne la remplace pas. »%
}{Principe du cours}

\niveauChapitre{Intermédiaire}{Le Chapitre~\ref{chap:orga}~; notions de réseaux.}

\begin{boxobjectifs}
\begin{itemize}
    \item Distinguer les modèles de contrôle d'accès (\sigle{DAC}, \sigle{MAC},
          \sigle{RBAC}).
    \item Expliquer l'authentification forte (\sigle{MFA}).
    \item Situer pare-feu, \sigle{VPN} et segmentation réseau.
    \item Comprendre le rôle du chiffrement et de la journalisation.
    \item Relier chaque mesure technique à un critère \sigle{DIC}.
\end{itemize}
\end{boxobjectifs}

\bigskip

% ============================================================================
\section*{Question directrice : quelles mesures techniques pour quels risques ?}
% ============================================================================

Les mesures techniques sont nombreuses~: le danger est de les empiler sans
logique. Ce chapitre les présente en les rattachant aux risques et aux critères
de sécurité de la mairie.

\medskip
\textbf{Plan du chapitre.}
\begin{enumerate}
    \item Le contrôle d'accès (section~\ref{sec:c16-acces}).
    \item L'authentification forte (section~\ref{sec:c16-mfa}).
    \item La sécurité réseau (section~\ref{sec:c16-reseau}).
    \item Chiffrement et journalisation (section~\ref{sec:c16-crypto}).
    \item Rattacher aux critères (section~\ref{sec:c16-dic}).
\end{enumerate}

\separateur

% ============================================================================
\section{Le Contrôle d'Accès}
\label{sec:c16-acces}
% ============================================================================

\begin{tableaucours}{L{2.2cm}L{4.4cm}L{5.6cm}}{%
    \tableauHeader{Modèle} &
    \tableauHeader{Principe} &
    \tableauHeader{Exemple}
}
DAC & Fondé sur l'identité, propriétaire & Permissions d'un fichier \\
MAC & Fondé sur des étiquettes & Politique centralisée (SELinux) \\
RBAC & Fondé sur les rôles & Un agent d'état civil voit l'état civil \\
\end{tableaucours}
\vspace{-0.8em}
\begin{center}{\small\itshape Tableau~16.1. Les modèles de contrôle d'accès}\end{center}

Le \termecle{moindre privilège} commande~: chacun n'a que les droits
nécessaires à sa tâche. À la mairie, le modèle \sigle{RBAC} convient~: on donne
des droits par rôle (état civil, taxes), non par personne.

% ============================================================================
\section{L'Authentification Forte}
\label{sec:c16-mfa}
% ============================================================================

\begin{boxdef}
\termecle{Authentification multifacteur} (\sigle{MFA})~: vérification de
l'identité par au moins deux facteurs parmi~: ce que l'on sait (mot de passe),
ce que l'on possède (jeton, téléphone), ce que l'on est (biométrie).
\end{boxdef}

La \sigle{MFA} réduit fortement le risque lié aux mots de passe volés ou
faibles. Pour les accès sensibles de la mairie (administration du serveur
d'état civil), elle devrait être la règle.

% ============================================================================
\section{La Sécurité Réseau}
\label{sec:c16-reseau}
% ============================================================================

\begin{tableaucours}{L{3.4cm}L{8.6cm}}{%
    \tableauHeader{Mesure} &
    \tableauHeader{Rôle}
}
Pare-feu & Filtrer les flux, bloquer par défaut \\
Segmentation & Isoler les réseaux (postes publics / serveurs) \\
VPN & Chiffrer les accès distants \\
IDS/IPS & Détecter ou bloquer les intrusions \\
\end{tableaucours}
\vspace{-0.8em}
\begin{center}{\small\itshape Tableau~16.2. Mesures de sécurité réseau}\end{center}

La \termecle{segmentation} aurait limité l'incident d'Eséka~: si les postes du
public avaient été isolés du serveur d'état civil, la clé n'aurait pas tout
atteint. Cloisonner, c'est contenir.

% ============================================================================
\section{Chiffrement et Journalisation}
\label{sec:c16-crypto}
% ============================================================================

\begin{boxdef}
\termecle{Chiffrement}~: rendre une information inintelligible sans la clé,
pour protéger sa \emph{confidentialité}. \termecle{Journalisation}~:
enregistrer les événements du système, pour la \emph{traçabilité} et la
détection.
\end{boxdef}

Chiffrer le fichier des contribuables protège sa confidentialité même en cas de
vol du support. Journaliser les accès à l'état civil permet de savoir qui a
consulté quoi, et soutient l'enquête en cas d'incident (lien avec le cours de
forensique, \sigle{LPI342}).

\begin{boxmetier}{DSI}
\textbf{Déployer ce que la politique décide.} Le \sigle{DSI} met en \oe uvre,
mais ne décide pas seul~: il déploie les mesures que l'appréciation des risques
a justifiées et que la politique a actées. La technique sert la décision, pas
l'inverse.
\end{boxmetier}

% ============================================================================
\section{Rattacher aux Critères}
\label{sec:c16-dic}
% ============================================================================

\begin{boxprocedure}
\textbf{Constatation.} La mairie déploie~: \sigle{RBAC}, \sigle{MFA},
segmentation, chiffrement, journalisation, sauvegarde.

\medskip
\textbf{Analyse.} Chaque mesure sert un critère~: contrôle d'accès et
chiffrement servent la \emph{confidentialité}~; sauvegarde et segmentation, la
\emph{disponibilité}~; journalisation, la \emph{traçabilité}~; signature et
contrôle d'intégrité, l'\emph{intégrité}. Rattacher chaque mesure à un critère
évite d'en oublier et d'en empiler d'inutiles.
\end{boxprocedure}

\separateur

\begin{boxresume}
\textbf{Ce qu'il faut retenir de ce chapitre~:}
\begin{itemize}
    \item \textbf{DAC, MAC, RBAC}~: trois modèles de contrôle d'accès~; le
          \textbf{moindre privilège} commande.
    \item La \textbf{MFA} réduit le risque des mots de passe volés.
    \item \textbf{Pare-feu, segmentation, VPN} cloisonnent et protègent le
          réseau.
    \item \textbf{Chiffrement} (confidentialité) et \textbf{journalisation}
          (traçabilité) sont essentiels.
    \item Chaque mesure technique se rattache à un \textbf{critère DIC}.
\end{itemize}
\end{boxresume}

\begin{boxinfo}
\textbf{Points de contrôle de l'élève.} Je sais~: distinguer DAC, MAC, RBAC~;
expliquer la MFA~; situer les mesures réseau~; décrire chiffrement et
journalisation~; rattacher une mesure à un critère.
\end{boxinfo}

\bigskip

\begin{boxexo}[title={\ding{45}\quad Exercice~16.1. Les modèles \niveaubase}]
Distinguez DAC, MAC et RBAC en une phrase chacun.
\end{boxexo}

\begin{boxexo}[title={\ding{45}\quad Exercice~16.2. MFA \niveaubase}]
Qu'est-ce que la MFA, et quels sont les trois types de facteurs~?
\end{boxexo}

\begin{boxexo}[title={\ding{45}\quad Exercice~16.3. Rôle d'une mesure \niveaubase}]
À quoi sert la segmentation réseau~?
\end{boxexo}

\begin{boxexo}[title={\ding{45}\quad Exercice~16.4. Moindre privilège \niveaumoyen}]
Appliquez le moindre privilège aux agents d'état civil et des taxes.
\end{boxexo}

\begin{boxexo}[title={\ding{45}\quad Exercice~16.5. Segmentation \niveaumoyen}]
Comment la segmentation aurait-elle limité l'incident de la clé \texttt{USB}~?
\end{boxexo}

\begin{boxexo}[title={\ding{45}\quad Exercice~16.6. Chiffrer quoi \niveaumoyen}]
Quelles informations de la mairie chiffreriez-vous en priorité, et pourquoi~?
\end{boxexo}

\begin{boxexo}[title={\ding{45}\quad Exercice~16.7. Journaliser \niveaumoyen}]
Quels événements journaliser sur le serveur d'état civil~?
\end{boxexo}

\begin{boxexo}[title={\ding{45}\quad Exercice~16.8. Mesure et critère \niveaumoyen}]
Reliez chaque mesure à un critère DIC~: chiffrement, sauvegarde,
journalisation, RBAC.
\end{boxexo}

\begin{boxexo}[title={\ding{45}\quad Exercice~16.9. Architecture \niveauexpert}]
Proposez une architecture réseau segmentée pour la mairie d'Eséka (public,
agents, serveurs).
\end{boxexo}

\begin{boxexo}[title={\ding{45}\quad Exercice~16.10. Plan technique \niveauexpert}]
Dressez le plan des mesures techniques de la mairie, chacune rattachée à un
risque et à un critère.
\end{boxexo}

\bigskip

\begin{boxinfo}
\textbf{Pour aller plus loin.}
\begin{itemize}
    \item ISO/IEC 27002, thème technologique~\cite{iso27002}~; Guide d'hygiène
          ANSSI~\cite{anssihygiene2017}.
\end{itemize}
\end{boxinfo}

\begin{boxinfo}
\textbf{Transition.} Reste à se préparer au pire~: le Chapitre~\ref{chap:continuite}
traite la continuité d'activité, \sigle{PCA} et \sigle{PRA}.
\end{boxinfo}
}{}
        \IfFileExists{partie5/P05_C17_continuite.tex}{% ============================================================================
% Fichier  : partie5/P05_C17_continuite.tex
% Titre    : Continuité d'Activité : PCA et PRA
% Auteur   : MINKA MI NGUIDJOÏ Thierry Emmanuel
% ============================================================================
% Sources : ISO 22301:2019 ; ISO/IEC 27002:2022. Mairie d'Eséka : cas fictif.
% ============================================================================

\chapter{Continuité d'Activité : PCA et PRA}
\label{chap:continuite}
\chaptermark{PCA et PRA}

\epigraph{%
    « Une sauvegarde que l'on n'a jamais restaurée\\
    n'est pas une sauvegarde~: c'est un espoir. »%
}{Principe du cours}

\niveauChapitre{Intermédiaire}{Le Chapitre~\ref{chap:technique}.}

\begin{boxobjectifs}
\begin{itemize}
    \item Distinguer plan de continuité (\sigle{PCA}) et plan de reprise
          (\sigle{PRA}).
    \item Conduire une analyse d'impact sur l'activité (\sigle{BIA}).
    \item Définir \sigle{RTO}, \sigle{RPO}, \sigle{MTTR}, \sigle{MTBF}.
    \item Appliquer la règle de sauvegarde \mbox{3-2-1} et tester les
          restaurations.
    \item Construire une stratégie de continuité adaptée à la mairie.
\end{itemize}
\end{boxobjectifs}

\bigskip

% ============================================================================
\section*{Question directrice : comment tenir le service quand le pire arrive ?}
% ============================================================================

L'incident a montré qu'une panne peut tout arrêter. La continuité prépare la
mairie à tenir et à reprendre. Ce chapitre construit le \sigle{PCA} et le
\sigle{PRA}.

\medskip
\textbf{Plan du chapitre.}
\begin{enumerate}
    \item PCA et PRA (section~\ref{sec:c17-pcapra}).
    \item L'analyse d'impact (section~\ref{sec:c17-bia}).
    \item Les indicateurs de continuité (section~\ref{sec:c17-indic}).
    \item Sauvegarde et restauration (section~\ref{sec:c17-sauvegarde}).
    \item La stratégie de la mairie (section~\ref{sec:c17-strategie}).
\end{enumerate}

\separateur

% ============================================================================
\section{PCA et PRA}
\label{sec:c17-pcapra}
% ============================================================================

\begin{boxdef}
\termecle{Plan de continuité d'activité} (\sigle{PCA})~: dispositif qui maintient
les activités critiques pendant un sinistre, éventuellement en mode dégradé.
\termecle{Plan de reprise d'activité} (\sigle{PRA})~: dispositif qui rétablit le
système après le sinistre. Le \sigle{PCA} tient~; le \sigle{PRA} répare.
\end{boxdef}

Pour la mairie, le \sigle{PCA} permet de continuer à délivrer des actes (même
sur registre papier de secours)~; le \sigle{PRA} remet en service le système
numérique.

% ============================================================================
\section{L'Analyse d'Impact}
\label{sec:c17-bia}
% ============================================================================

\begin{boxdef}
\termecle{Analyse d'impact sur l'activité} (\sigle{BIA})~: étude qui identifie
les activités critiques, évalue l'impact de leur interruption et fixe les délais
tolérables. Elle fonde toute la stratégie de continuité.
\end{boxdef}

Le \sigle{BIA} répond à~: quelles activités ne peuvent s'arrêter~? combien de
temps peut-on tenir sans elles~? quelle perte de données est tolérable~? À
Eséka, l'état civil est l'activité la plus critique.

% ============================================================================
\section{Les Indicateurs de Continuité}
\label{sec:c17-indic}
% ============================================================================

\begin{tableaucours}{L{2.0cm}L{10.2cm}}{%
    \tableauHeader{Indicateur} &
    \tableauHeader{Définition}
}
RTO & Délai maximal de remise en service visé \\
RPO & Perte de données maximale acceptable (point de reprise) \\
MTTR & Temps moyen de réparation \\
MTBF & Temps moyen entre deux pannes \\
\end{tableaucours}
\vspace{-0.8em}
\begin{center}{\small\itshape Tableau~17.1. Indicateurs de continuité}\end{center}

\sigle{RTO} et \sigle{RPO} commandent la stratégie~: un \sigle{RPO} d'une heure
impose des sauvegardes au moins horaires~; un \sigle{RTO} de quatre heures
impose des moyens de reprise rapides.

% ============================================================================
\section{Sauvegarde et Restauration}
\label{sec:c17-sauvegarde}
% ============================================================================

\begin{boxdef}
La règle \termecle{3-2-1}~: conserver \textbf{3} copies des données, sur
\textbf{2} supports différents, dont \textbf{1} hors site. Elle protège contre
la panne, l'erreur et le sinistre local.
\end{boxdef}

\begin{boxprocedure}
\textbf{Constatation.} Avant la démarche, la mairie avait une seule copie de
l'état civil, sur le même serveur que l'original, jamais testée.

\medskip
\textbf{Analyse} (confiance élevée). Cette « sauvegarde » violait la règle
3-2-1 (une copie, un support, sur site) et n'avait jamais été restaurée~: elle
ne protégeait rien. La mesure corrective impose trois copies, dont une hors
site, et un \emph{test de restauration} mensuel. Une sauvegarde non testée est
réputée inexistante.
\end{boxprocedure}

\begin{boxmetier}{DSI}
\textbf{Tester, pas seulement sauvegarder.} Le \sigle{DSI} planifie des tests
de restauration réguliers~: c'est le seul moyen de savoir que la sauvegarde
fonctionne \emph{avant} d'en avoir besoin. Le jour du sinistre est le pire
moment pour découvrir qu'elle est illisible.
\end{boxmetier}

% ============================================================================
\section{La Stratégie de la Mairie}
\label{sec:c17-strategie}
% ============================================================================

La mairie d'Eséka adopte une stratégie proportionnée~: sauvegarde quotidienne
de l'état civil selon la règle 3-2-1, copie hors site chez la tutelle, registre
papier de secours pour le mode dégradé, et test de restauration mensuel. Le
\sigle{PCA} et le \sigle{PRA} sont écrits, diffusés et exercés.

\separateur

\begin{boxresume}
\textbf{Ce qu'il faut retenir de ce chapitre~:}
\begin{itemize}
    \item Le \textbf{PCA} maintient l'activité~; le \textbf{PRA} rétablit le
          système.
    \item Le \textbf{BIA} identifie les activités critiques et fixe les délais
          tolérables.
    \item \textbf{RTO} et \textbf{RPO} commandent la stratégie de continuité.
    \item La règle \textbf{3-2-1} et le \textbf{test de restauration} font
          d'une copie une vraie sauvegarde.
    \item La mairie adopte une stratégie \textbf{proportionnée} et
          \textbf{exercée}.
\end{itemize}
\end{boxresume}

\begin{boxinfo}
\textbf{Points de contrôle de l'élève.} Je sais~: distinguer PCA et PRA~;
conduire un BIA~; définir RTO, RPO, MTTR, MTBF~; appliquer la règle 3-2-1~;
bâtir une stratégie de continuité.
\end{boxinfo}

\bigskip

\begin{boxexo}[title={\ding{45}\quad Exercice~17.1. PCA ou PRA \niveaubase}]
« Continuer à délivrer des actes pendant la panne » relève-t-il du PCA ou du
PRA~?
\end{boxexo}

\begin{boxexo}[title={\ding{45}\quad Exercice~17.2. Les indicateurs \niveaubase}]
Définissez RTO et RPO.
\end{boxexo}

\begin{boxexo}[title={\ding{45}\quad Exercice~17.3. La règle 3-2-1 \niveaubase}]
Énoncez la règle 3-2-1 de sauvegarde.
\end{boxexo}

\begin{boxexo}[title={\ding{45}\quad Exercice~17.4. BIA \niveaumoyen}]
Quelle est l'activité la plus critique de la mairie, et pourquoi~?
\end{boxexo}

\begin{boxexo}[title={\ding{45}\quad Exercice~17.5. RPO et fréquence \niveaumoyen}]
Un RPO d'une heure impose quelle fréquence de sauvegarde~? Justifiez.
\end{boxexo}

\begin{boxexo}[title={\ding{45}\quad Exercice~17.6. Fausse sauvegarde \niveaumoyen}]
Pourquoi la « sauvegarde » initiale de la mairie ne protégeait-elle rien~?
\end{boxexo}

\begin{boxexo}[title={\ding{45}\quad Exercice~17.7. Tester \niveaumoyen}]
Pourquoi faut-il tester les restaurations, et à quelle fréquence~?
\end{boxexo}

\begin{boxexo}[title={\ding{45}\quad Exercice~17.8. Mode dégradé \niveaumoyen}]
Décrivez un mode dégradé pour l'état civil en cas de panne prolongée.
\end{boxexo}

\begin{boxexo}[title={\ding{45}\quad Exercice~17.9. Stratégie 3-2-1 \niveauexpert}]
Concevez la stratégie de sauvegarde 3-2-1 complète de la mairie d'Eséka.
\end{boxexo}

\begin{boxexo}[title={\ding{45}\quad Exercice~17.10. PRA rédigé \niveauexpert}]
Rédigez un plan de reprise d'activité d'une page pour le serveur d'état civil
(RTO, RPO, étapes, responsables).
\end{boxexo}

\bigskip

\begin{boxinfo}
\textbf{Pour aller plus loin.}
\begin{itemize}
    \item ISO 22301~\cite{iso22301}~; ISO/IEC 27002~\cite{iso27002}.
\end{itemize}
\end{boxinfo}

\begin{boxinfo}
\textbf{Transition.} Les mesures en place, il faut les piloter~: la partie VI
s'ouvre au Chapitre~\ref{chap:indicateurs} sur la surveillance et les
indicateurs.
\end{boxinfo}
}{}

    \part{Contrôler, Auditer, Améliorer}
    \label{part:controler}

        \IfFileExists{partie6/P06_C18_indicateurs.tex}{% ============================================================================
% Fichier  : partie6/P06_C18_indicateurs.tex
% Titre    : Surveillance, Mesure et Indicateurs
% Auteur   : MINKA MI NGUIDJOÏ Thierry Emmanuel
% ============================================================================
% Sources : ISO/IEC 27004:2016 ; ISO/IEC 27001:2022 (clause 9.1). Fictif : Eséka.
% ============================================================================

\chapter{Surveillance, Mesure et Indicateurs}
\label{chap:indicateurs}
\chaptermark{Mesure et indicateurs}

\epigraph{%
    « Ce qui ne se mesure pas ne se pilote pas~;\\
    ce qui ne se pilote pas se dégrade. »%
}{Principe du cours}

\niveauChapitre{Intermédiaire}{La partie V.}

\begin{boxobjectifs}
\begin{itemize}
    \item Comprendre l'exigence de surveillance et de mesure (clause~9.1).
    \item Distinguer indicateurs de performance et de risque.
    \item Construire un indicateur utile (mesurable, suivi, décidable).
    \item Bâtir un tableau de bord de sécurité pour la direction.
    \item Relier la mesure à la décision et à l'amélioration.
\end{itemize}
\end{boxobjectifs}

\bigskip

% ============================================================================
\section*{Question directrice : comment savoir si le SMSI fonctionne vraiment ?}
% ============================================================================

La mairie a déployé des mesures. Fonctionnent-elles~? Sans mesure, nul ne le
sait. Ce chapitre construit les indicateurs et le tableau de bord qui
répondent.

\medskip
\textbf{Plan du chapitre.}
\begin{enumerate}
    \item L'exigence de mesure (section~\ref{sec:c18-exigence}).
    \item Performance et risque (section~\ref{sec:c18-kpikri}).
    \item Un bon indicateur (section~\ref{sec:c18-bon}).
    \item Le tableau de bord (section~\ref{sec:c18-tableau}).
    \item De la mesure à la décision (section~\ref{sec:c18-decision}).
\end{enumerate}

\separateur

% ============================================================================
\section{L'Exigence de Mesure}
\label{sec:c18-exigence}
% ============================================================================

\begin{boxdef}
La clause~9.1 de \norme{ISO/IEC 27001} exige de déterminer \emph{quoi}
surveiller et mesurer, \emph{comment}, \emph{quand} et \emph{qui} analyse les
résultats. La norme \norme{ISO/IEC 27004} en donne les lignes directrices.
\end{boxdef}

Mesurer n'est pas accumuler des chiffres~: c'est choisir les quelques
indicateurs qui éclairent une décision. Un bon dispositif de mesure est
sélectif.

% ============================================================================
\section{Performance et Risque}
\label{sec:c18-kpikri}
% ============================================================================

\begin{boxdef}
Un \termecle{indicateur de performance} (\sigle{KPI}) mesure l'efficacité d'une
mesure (taux d'agents formés). Un \termecle{indicateur de risque} (\sigle{KRI})
alerte sur une exposition croissante (nombre de comptes sans \sigle{MFA}).
\end{boxdef}

Les deux se complètent~: les \sigle{KPI} disent « où en est-on », les
\sigle{KRI} disent « où le danger monte ». Un tableau de bord équilibré
comporte les deux.

% ============================================================================
\section{Un Bon Indicateur}
\label{sec:c18-bon}
% ============================================================================

\begin{tableaucours}{L{3.6cm}L{8.6cm}}{%
    \tableauHeader{Qualité} &
    \tableauHeader{Signification}
}
Mesurable & On peut le calculer objectivement \\
Pertinent & Il éclaire une décision réelle \\
Suivi & Il est relevé à intervalle régulier \\
Doté d'une cible & On sait ce qui est bon ou mauvais \\
Attribué & Quelqu'un en répond \\
\end{tableaucours}
\vspace{-0.8em}
\begin{center}{\small\itshape Tableau~18.1. Qualités d'un bon indicateur}\end{center}

Un indicateur que personne ne regarde ou qui ne change aucune décision est à
supprimer~: il coûte sans servir.

% ============================================================================
\section{Le Tableau de Bord}
\label{sec:c18-tableau}
% ============================================================================

\begin{tableaucours}{L{4.6cm}L{2.2cm}L{2.2cm}L{2.0cm}}{%
    \tableauHeader{Indicateur} &
    \tableauHeader{Type} &
    \tableauHeader{Cible} &
    \tableauHeader{État}
}
Agents formés & KPI & 100 \% & 80 \% \\
Sauvegardes testées & KPI & Mensuel & À jour \\
Comptes sans MFA & KRI & 0 & 5 \\
Incidents du mois & KRI & -- & 2 \\
\end{tableaucours}
\vspace{-0.8em}
\begin{center}{\small\itshape Tableau~18.2. Extrait du tableau de bord de la mairie}\end{center}

\begin{boxmetier}{RSSI}
\textbf{Un tableau de bord pour décider.} La RSSI présente au comité un tableau
d'une page, lisible par un non-technicien. Chaque indicateur hors cible
appelle une décision~: former les agents restants, activer la MFA manquante.
Le tableau de bord n'est pas un rapport~: c'est un outil de pilotage.
\end{boxmetier}

% ============================================================================
\section{De la Mesure à la Décision}
\label{sec:c18-decision}
% ============================================================================

\begin{boxprocedure}
\textbf{Constatation.} Le tableau de bord montre 5~comptes sans \sigle{MFA} et
un taux de formation de 80~\%.

\medskip
\textbf{Analyse.} Deux écarts à la cible appellent deux décisions~: activer la
MFA sur les cinq comptes, achever la formation des agents restants. La mesure
débouche sur l'action~: c'est le lien entre « Check » et « Act » du
\sigle{PDCA} (Chapitre~\ref{chap:pdca}).
\end{boxprocedure}

La mesure alimente la revue de direction et l'audit (Chapitre~\ref{chap:audit})~:
elle fournit les preuves chiffrées que le \sigle{SMSI} fonctionne et s'améliore.

\separateur

\begin{boxresume}
\textbf{Ce qu'il faut retenir de ce chapitre~:}
\begin{itemize}
    \item La \textbf{clause~9.1} exige de mesurer quoi, comment, quand et par
          qui.
    \item Les \textbf{KPI} mesurent la performance~; les \textbf{KRI} alertent
          sur le risque.
    \item Un \textbf{bon indicateur} est mesurable, pertinent, suivi, doté
          d'une cible, attribué.
    \item Le \textbf{tableau de bord} est un outil de pilotage, lisible par la
          direction.
    \item La \textbf{mesure débouche sur la décision} et nourrit l'audit.
\end{itemize}
\end{boxresume}

\begin{boxinfo}
\textbf{Points de contrôle de l'élève.} Je sais~: énoncer l'exigence de la
clause~9.1~; distinguer KPI et KRI~; reconnaître un bon indicateur~; bâtir un
tableau de bord~; relier mesure et décision.
\end{boxinfo}

\bigskip

\begin{boxexo}[title={\ding{45}\quad Exercice~18.1. Quoi mesurer \niveaubase}]
Que doit déterminer la clause~9.1 au sujet de la mesure~?
\end{boxexo}

\begin{boxexo}[title={\ding{45}\quad Exercice~18.2. KPI ou KRI \niveaubase}]
« Taux d'agents formés » et « comptes sans MFA »~: lequel est un KPI, lequel un
KRI~?
\end{boxexo}

\begin{boxexo}[title={\ding{45}\quad Exercice~18.3. Bon indicateur \niveaubase}]
Citez trois qualités d'un bon indicateur.
\end{boxexo}

\begin{boxexo}[title={\ding{45}\quad Exercice~18.4. Indicateur inutile \niveaumoyen}]
Pourquoi supprimer un indicateur que personne ne regarde~?
\end{boxexo}

\begin{boxexo}[title={\ding{45}\quad Exercice~18.5. Construire un KPI \niveaumoyen}]
Construisez un KPI pour la mesure « sauvegarde de l'état civil ».
\end{boxexo}

\begin{boxexo}[title={\ding{45}\quad Exercice~18.6. Hors cible \niveaumoyen}]
Pour chaque indicateur hors cible du tableau de bord, proposez une décision.
\end{boxexo}

\begin{boxexo}[title={\ding{45}\quad Exercice~18.7. Lisible \niveaumoyen}]
Pourquoi le tableau de bord doit-il être lisible par un non-technicien~?
\end{boxexo}

\begin{boxexo}[title={\ding{45}\quad Exercice~18.8. Check et Act \niveaumoyen}]
Montrez que la mesure relie les étapes Check et Act du PDCA.
\end{boxexo}

\begin{boxexo}[title={\ding{45}\quad Exercice~18.9. Tableau de bord \niveauexpert}]
Construisez un tableau de bord de six indicateurs pour la mairie d'Eséka (type,
cible, responsable).
\end{boxexo}

\begin{boxexo}[title={\ding{45}\quad Exercice~18.10. Note de pilotage \niveauexpert}]
Rédigez la note mensuelle de pilotage que la RSSI présente au comité de
sécurité.
\end{boxexo}

\bigskip

\begin{boxinfo}
\textbf{Pour aller plus loin.}
\begin{itemize}
    \item ISO/IEC 27004~\cite{iso27004}~; ISO/IEC 27001, clause~9.1~\cite{iso27001}.
\end{itemize}
\end{boxinfo}

\begin{boxinfo}
\textbf{Transition.} Mesurer soi-même ne suffit pas~: il faut un regard
extérieur. Le Chapitre~\ref{chap:audit} traite l'audit et la certification.
\end{boxinfo}
}{}
        \IfFileExists{partie6/P06_C19_audit_certification.tex}{% ============================================================================
% Fichier  : partie6/P06_C19_audit_certification.tex
% Titre    : L'Audit du SMSI et la Certification
% Auteur   : MINKA MI NGUIDJOÏ Thierry Emmanuel
% ============================================================================
% Sources : ISO 19011:2018 ; ISO/IEC 27006:2015 ; ISO/IEC 27001 clause 9.2-9.3.
% Mairie d'Eséka : cas fictif.
% ============================================================================

\chapter{L'Audit du SMSI et la Certification}
\label{chap:audit}
\chaptermark{Audit et certification}

\epigraph{%
    « L'audit ne cherche pas le coupable~:\\
    il cherche l'écart entre ce qui est dit et ce qui est fait. »%
}{Principe du cours}

\niveauChapitre{Intermédiaire}{Le Chapitre~\ref{chap:indicateurs}.}

\begin{boxobjectifs}
\begin{itemize}
    \item Distinguer audit interne et audit de certification.
    \item Décrire la méthodologie d'audit (ISO 19011).
    \item Caractériser constat, preuve, non-conformité et écart.
    \item Expliquer la revue de direction.
    \item Décrire le cycle de certification \norme{ISO/IEC 27001}.
\end{itemize}
\end{boxobjectifs}

\bigskip

% ============================================================================
\section*{Question directrice : comment faire reconnaître que la sécurité est maîtrisée ?}
% ============================================================================

La mairie pense son \sigle{SMSI} solide. Un audit le vérifie, et la
certification l'atteste. Ce chapitre présente l'audit, la revue de direction
et le chemin vers la certification.

\medskip
\textbf{Plan du chapitre.}
\begin{enumerate}
    \item Les types d'audit (section~\ref{sec:c19-types}).
    \item La méthodologie (section~\ref{sec:c19-methodo}).
    \item Constat, preuve, écart (section~\ref{sec:c19-constat}).
    \item La revue de direction (section~\ref{sec:c19-revue}).
    \item Le cycle de certification (section~\ref{sec:c19-cycle}).
\end{enumerate}

\separateur

% ============================================================================
\section{Les Types d'Audit}
\label{sec:c19-types}
% ============================================================================

\begin{tableaucours}{L{3.2cm}L{9.0cm}}{%
    \tableauHeader{Type} &
    \tableauHeader{Qui et pourquoi}
}
Audit interne & Mené par l'organisation, exigé par la clause~9.2 \\
Audit de seconde partie & Mené par un client ou une tutelle \\
Audit de certification & Mené par un organisme accrédité (tierce partie) \\
\end{tableaucours}
\vspace{-0.8em}
\begin{center}{\small\itshape Tableau~19.1. Les types d'audit}\end{center}

L'\termecle{audit interne} prépare~: il détecte les écarts avant l'auditeur
externe. L'\termecle{audit de certification} atteste, auprès des tiers, que le
\sigle{SMSI} satisfait la norme.

% ============================================================================
\section{La Méthodologie}
\label{sec:c19-methodo}
% ============================================================================

La norme \norme{ISO 19011} guide l'audit des systèmes de management~:
planification, revue documentaire, entretiens et observations sur site,
rédaction du rapport, suivi des actions. L'auditeur est \termecle{indépendant}
de ce qu'il audite et procède par \termecle{échantillonnage}.

\begin{boxprocedure}
\textbf{Constatation.} Pour l'opération \textsc{Matomb}, le cabinet
\textsc{Matomb Conseil} conduit l'audit interne de la mairie avant la
certification.

\medskip
\textbf{Analyse.} L'audit interne est une répétition générale~: il cherche les
écarts pendant qu'il est encore temps de les corriger. Un écart trouvé en
interne se corrige~; le même écart trouvé en certification coûte un report.
\end{boxprocedure}

% ============================================================================
\section{Constat, Preuve, Écart}
\label{sec:c19-constat}
% ============================================================================

\begin{boxdef}
\termecle{Preuve d'audit}~: enregistrement vérifiable (document, observation,
déclaration). \termecle{Constat}~: résultat de l'évaluation des preuves face à
un critère. \termecle{Non-conformité}~: écart entre l'exigence et la réalité,
\emph{majeure} (défaillance du système) ou \emph{mineure} (manquement isolé).
\end{boxdef}

L'auditeur sépare le \emph{fait} (ce qu'il constate, preuve à l'appui) de son
\emph{appréciation}. Cette rigueur, qui distingue la constatation de l'analyse,
est au cœur du métier d'auditeur~\cite{minka2026auditologie}.

\begin{boxmetier}{Auditeur SMSI}
\textbf{Pas de constat sans preuve.} L'auditeur ne note jamais une impression~:
il adosse chaque constat à une preuve (un document, un relevé, une observation).
Un écart sans preuve n'est pas un écart~: c'est un avis.
\end{boxmetier}

% ============================================================================
\section{La Revue de Direction}
\label{sec:c19-revue}
% ============================================================================

\begin{boxdef}
\termecle{Revue de direction} (clause~9.3)~: examen périodique du \sigle{SMSI}
par la direction, qui passe en revue les résultats des audits, des mesures et
des incidents, et décide des améliorations et des ressources.
\end{boxdef}

La revue de direction est le moment où la direction \emph{reprend la main}~:
elle transforme les constats en décisions. Sans elle, l'audit resterait sans
suite. À Eséka, le maire la préside deux fois par an.

% ============================================================================
\section{Le Cycle de Certification}
\label{sec:c19-cycle}
% ============================================================================

\begin{tableaucours}{L{3.6cm}L{8.6cm}}{%
    \tableauHeader{Étape} &
    \tableauHeader{Contenu}
}
Audit initial (étape 1) & Revue documentaire, préparation \\
Audit initial (étape 2) & Audit sur site, décision de certification \\
Audits de surveillance & Chaque année, maintien du certificat \\
Renouvellement & Au bout de trois ans, nouvel audit complet \\
\end{tableaucours}
\vspace{-0.8em}
\begin{center}{\small\itshape Tableau~19.2. Le cycle de certification ISO/IEC 27001}\end{center}

La certification se \emph{gagne} en deux étapes et se \emph{maintient} par des
audits de surveillance. Elle n'est pas acquise une fois pour toutes~: le
\sigle{SMSI} doit continuer de tourner (Chapitre~\ref{chap:pdca}).

\separateur

\begin{boxresume}
\textbf{Ce qu'il faut retenir de ce chapitre~:}
\begin{itemize}
    \item L'\textbf{audit interne} prépare~; l'\textbf{audit de certification}
          atteste.
    \item La \textbf{méthodologie} (ISO 19011) procède par preuves,
          échantillonnage et indépendance.
    \item Un \textbf{constat} s'adosse à une \textbf{preuve}~; une
          \textbf{non-conformité} est majeure ou mineure.
    \item La \textbf{revue de direction} transforme les constats en décisions.
    \item La \textbf{certification} se gagne en deux étapes et se maintient par
          surveillance.
\end{itemize}
\end{boxresume}

\begin{boxinfo}
\textbf{Points de contrôle de l'élève.} Je sais~: distinguer les types
d'audit~; décrire la méthodologie~; définir preuve, constat, non-conformité~;
expliquer la revue de direction~; décrire le cycle de certification.
\end{boxinfo}

\bigskip

\begin{boxexo}[title={\ding{45}\quad Exercice~19.1. Types d'audit \niveaubase}]
Distinguez audit interne et audit de certification.
\end{boxexo}

\begin{boxexo}[title={\ding{45}\quad Exercice~19.2. Preuve \niveaubase}]
Qu'est-ce qu'une preuve d'audit~? Donnez deux exemples.
\end{boxexo}

\begin{boxexo}[title={\ding{45}\quad Exercice~19.3. Le cycle \niveaubase}]
Énoncez les étapes du cycle de certification \norme{ISO/IEC 27001}.
\end{boxexo}

\begin{boxexo}[title={\ding{45}\quad Exercice~19.4. Majeure ou mineure \niveaumoyen}]
Distinguez une non-conformité majeure d'une mineure, par un exemple à la
mairie.
\end{boxexo}

\begin{boxexo}[title={\ding{45}\quad Exercice~19.5. Répétition générale \niveaumoyen}]
Pourquoi l'audit interne est-il une « répétition générale »~?
\end{boxexo}

\begin{boxexo}[title={\ding{45}\quad Exercice~19.6. Constat sans preuve \niveaumoyen}]
Pourquoi un écart sans preuve n'est-il qu'un avis~?
\end{boxexo}

\begin{boxexo}[title={\ding{45}\quad Exercice~19.7. Indépendance \niveaumoyen}]
Pourquoi l'auditeur doit-il être indépendant de ce qu'il audite~?
\end{boxexo}

\begin{boxexo}[title={\ding{45}\quad Exercice~19.8. Revue de direction \niveaumoyen}]
Que décide la direction lors de la revue, et à partir de quoi~?
\end{boxexo}

\begin{boxexo}[title={\ding{45}\quad Exercice~19.9. Plan d'audit \niveauexpert}]
Rédigez le plan d'audit interne du \sigle{SMSI} de la mairie d'Eséka (périmètre,
critères, étapes).
\end{boxexo}

\begin{boxexo}[title={\ding{45}\quad Exercice~19.10. Fiche d'écart \niveauexpert}]
Rédigez une fiche de non-conformité (constat, preuve, exigence, gravité, action
corrective) pour un écart de la mairie.
\end{boxexo}

\bigskip

\begin{boxinfo}
\textbf{Pour aller plus loin.}
\begin{itemize}
    \item ISO 19011~\cite{iso19011}~; ISO/IEC 27006~\cite{iso27006}~;
          auditologie~\cite{minka2026auditologie}.
\end{itemize}
\end{boxinfo}

\begin{boxinfo}
\textbf{Transition.} Reste le dernier maillon du cycle~: apprendre des
incidents. Le Chapitre~\ref{chap:incidents} clôt le cours.
\end{boxinfo}
}{}
        \IfFileExists{partie6/P06_C20_incidents_progres.tex}{% ============================================================================
% Fichier  : partie6/P06_C20_incidents_progres.tex
% Titre    : Gérer les Incidents et Progresser
% Auteur   : MINKA MI NGUIDJOÏ Thierry Emmanuel
% ============================================================================
% Sources : ISO/IEC 27035 ; ISO/IEC 27001 clause 10. Mairie d'Eséka : fictif.
% ============================================================================

\chapter{Gérer les Incidents et Progresser}
\label{chap:incidents}
\chaptermark{Incidents et progrès}

\epigraph{%
    « Un incident bien géré n'est pas un échec~:\\
    c'est le moment où le système apprend. »%
}{Principe du cours}

\niveauChapitre{Intermédiaire}{La partie VI.}

\begin{boxobjectifs}
\begin{itemize}
    \item Décrire le processus de gestion des incidents.
    \item Classer et prioriser un incident.
    \item Conduire un retour d'expérience (\sigle{RETEX}).
    \item Relier l'incident à l'amélioration continue (clause~10).
    \item Mesurer les progrès du \sigle{SMSI} dans la durée.
\end{itemize}
\end{boxobjectifs}

\bigskip

% ============================================================================
\section*{Question directrice : comment transformer un incident en progrès ?}
% ============================================================================

L'histoire de la mairie a commencé par un incident. Ce dernier chapitre montre
comment un \sigle{SMSI} gère les incidents et s'en sert pour progresser,
bouclant le cycle du cours.

\medskip
\textbf{Plan du chapitre.}
\begin{enumerate}
    \item Le processus de gestion (section~\ref{sec:c20-processus}).
    \item Classer et prioriser (section~\ref{sec:c20-classer}).
    \item Le retour d'expérience (section~\ref{sec:c20-retex}).
    \item L'incident, moteur d'amélioration (section~\ref{sec:c20-amelioration}).
    \item Mesurer les progrès (section~\ref{sec:c20-progres}).
\end{enumerate}

\separateur

% ============================================================================
\section{Le Processus de Gestion}
\label{sec:c20-processus}
% ============================================================================

\begin{boxdef}
La \termecle{gestion des incidents}~: processus qui détecte, signale, qualifie,
traite et clôt les incidents de sécurité, en tirant les leçons de chacun
(\norme{ISO/IEC 27035}).
\end{boxdef}

\begin{tableaucours}{L{3.0cm}L{9.2cm}}{%
    \tableauHeader{Étape} &
    \tableauHeader{Contenu}
}
Détecter et signaler & Repérer et alerter, sans crainte de sanction \\
Qualifier & Confirmer l'incident, en mesurer la portée \\
Traiter & Contenir, éradiquer, rétablir \\
Clore et apprendre & Documenter et conduire le RETEX \\
\end{tableaucours}
\vspace{-0.8em}
\begin{center}{\small\itshape Tableau~20.1. Le processus de gestion des incidents}\end{center}

Le signalement est la clé~: un incident tu ne peut être traité. La culture du
signalement sans blâme, posée à la sensibilisation (Chapitre~\ref{chap:orga}),
porte ici ses fruits.

% ============================================================================
\section{Classer et Prioriser}
\label{sec:c20-classer}
% ============================================================================

\begin{boxprocedure}
\textbf{Constatation.} La mairie reçoit deux signalements~: un courriel
d'hameçonnage non ouvert, et un poste d'état civil qui redémarre seul.

\medskip
\textbf{Analyse.} On classe selon l'\emph{impact} et l'\emph{urgence}~: le
courriel non ouvert est mineur (surveiller)~; le poste d'état civil qui
redémarre touche un actif critique (priorité haute, isoler et investiguer).
Prioriser, c'est mettre les moyens là où l'enjeu est le plus fort.
\end{boxprocedure}

% ============================================================================
\section{Le Retour d'Expérience}
\label{sec:c20-retex}
% ============================================================================

\begin{boxdef}
\termecle{Retour d'expérience} (\sigle{RETEX})~: analyse, après un incident, de
ce qui s'est passé, de ce qui a bien ou mal fonctionné, et des améliorations à
apporter. Il vise la cause, non le coupable.
\end{boxdef}

Le \sigle{RETEX} de l'incident initial de la mairie a produit des mesures
durables~: interdiction des clés personnelles, sauvegardes 3-2-1, comptes
nominatifs, sensibilisation. L'incident, bien analysé, a été le point de départ
de tout le \sigle{SMSI}.

\begin{boxmetier}{RSSI}
\textbf{Chercher la cause, pas la faute.} La RSSI conduit le \sigle{RETEX} dans
un esprit d'amélioration~: on cherche pourquoi le système a permis l'incident,
non qui a cliqué. Un RETEX accusateur tarit les signalements~; un RETEX
constructif les encourage.
\end{boxmetier}

% ============================================================================
\section{L'Incident, Moteur d'Amélioration}
\label{sec:c20-amelioration}
% ============================================================================

L'incident alimente la clause~10 (amélioration)~: chaque \sigle{RETEX} génère
des actions correctives qui éliminent des causes, et parfois des améliorations
qui renforcent au-delà. C'est le « Act » du \sigle{PDCA}~: on referme la boucle
et l'on relance un cycle plus mûr.

% ============================================================================
\section{Mesurer les Progrès}
\label{sec:c20-progres}
% ============================================================================

\begin{boxobjet}{Le SMSI de la mairie, un an après}
\textbf{Constatation.} Un an après l'incident, la mairie a~: une politique
signée, des risques appréciés, des mesures déployées, un tableau de bord, un
audit interne passé, une certification en vue.

\medskip
\textbf{Analyse.} Le progrès se mesure~: délai de reprise réduit, sauvegardes
testées, agents formés, incidents mieux traités. La mairie ne promet pas le
risque zéro~; elle démontre un risque \emph{maîtrisé}. C'est l'aboutissement du
cours~: non pas un système parfait, mais un système qui apprend.
\end{boxobjet}

\separateur

\begin{boxresume}
\textbf{Ce qu'il faut retenir de ce chapitre~:}
\begin{itemize}
    \item La \textbf{gestion des incidents} détecte, signale, qualifie, traite
          et clôt.
    \item On \textbf{classe} et \textbf{priorise} selon l'impact et l'urgence.
    \item Le \textbf{RETEX} cherche la cause, non le coupable.
    \item L'incident nourrit l'\textbf{amélioration continue} (clause~10).
    \item Le progrès se \textbf{mesure}~: un \sigle{SMSI} mûr maîtrise le
          risque, sans l'annuler.
\end{itemize}
\end{boxresume}

\begin{boxinfo}
\textbf{Points de contrôle de l'élève.} Je sais~: décrire le processus de
gestion des incidents~; classer et prioriser~; conduire un RETEX~; relier
l'incident à l'amélioration~; mesurer les progrès du \sigle{SMSI}.
\end{boxinfo}

\bigskip

\begin{boxexo}[title={\ding{45}\quad Exercice~20.1. Le processus \niveaubase}]
Énoncez les quatre étapes de la gestion des incidents.
\end{boxexo}

\begin{boxexo}[title={\ding{45}\quad Exercice~20.2. Signaler \niveaubase}]
Pourquoi le signalement sans blâme est-il essentiel~?
\end{boxexo}

\begin{boxexo}[title={\ding{45}\quad Exercice~20.3. RETEX \niveaubase}]
Qu'est-ce qu'un retour d'expérience, et que vise-t-il~?
\end{boxexo}

\begin{boxexo}[title={\ding{45}\quad Exercice~20.4. Prioriser \niveaumoyen}]
Entre le courriel non ouvert et le poste qui redémarre, lequel traiter d'abord~?
Pourquoi~?
\end{boxexo}

\begin{boxexo}[title={\ding{45}\quad Exercice~20.5. Cause ou faute \niveaumoyen}]
Pourquoi un RETEX accusateur tarit-il les signalements~?
\end{boxexo}

\begin{boxexo}[title={\ding{45}\quad Exercice~20.6. Incident et clause 10 \niveaumoyen}]
Comment un incident alimente-t-il l'amélioration continue~?
\end{boxexo}

\begin{boxexo}[title={\ding{45}\quad Exercice~20.7. L'incident fondateur \niveaumoyen}]
Montrez comment l'incident initial de la mairie a produit des mesures durables.
\end{boxexo}

\begin{boxexo}[title={\ding{45}\quad Exercice~20.8. Mesurer le progrès \niveaumoyen}]
Citez trois indicateurs montrant que le \sigle{SMSI} de la mairie a progressé.
\end{boxexo}

\begin{boxexo}[title={\ding{45}\quad Exercice~20.9. Procédure d'incident \niveauexpert}]
Rédigez la procédure de gestion des incidents de la mairie d'Eséka (détection à
clôture).
\end{boxexo}

\begin{boxexo}[title={\ding{45}\quad Exercice~20.10. RETEX complet \niveauexpert}]
Conduisez le RETEX écrit de l'incident initial de la mairie (déroulé, causes,
actions correctives, améliorations).
\end{boxexo}

\bigskip

\begin{boxinfo}
\textbf{Pour aller plus loin.}
\begin{itemize}
    \item ISO/IEC 27001, clause~10~\cite{iso27001}~; auditologie~\cite{minka2026auditologie}.
\end{itemize}
\end{boxinfo}

\begin{boxinfo}
\textbf{Clôture du cours.} De l'incident à la certification, la mairie d'Eséka
a appris à \emph{décider} sa sécurité, à la \emph{tenir} dans la durée et à la
\emph{prouver}. Tel est le management de la sécurité de l'information~: non pas
un état que l'on atteint, mais une roue que l'on fait tourner. Planifier,
déployer, contrôler, agir~: et recommencer, un cran plus haut.
\end{boxinfo}
}{}

% --------------------------------------------------------------------------
% ANNEXES
% --------------------------------------------------------------------------
\appendix

    %% Annexe A : Glossaire
\chapter{Glossaire}
\label{ann:glossaire}

Ce glossaire reprend les notions définies dans les chapitres. Le renvoi
indique le chapitre où la notion est présentée.

\begin{description}
    \item[Actif] Élément de valeur pour l'organisation, qu'il faut protéger (Ch.~\ref{chap:cia}).
    \item[Appétence au risque] Niveau de risque qu'une organisation est prête à accepter (Ch.~\ref{chap:vocabrisque}).
    \item[Appréciation des risques] Démarche d'identification, d'analyse et d'évaluation des risques (Ch.~\ref{chap:methodes}).
    \item[Déclaration d'applicabilité (DdA)] Document indiquant, pour chaque mesure de l'Annexe~A, si elle est retenue, sa justification et son état (Ch.~\ref{chap:dda}).
    \item[Disponibilité, Intégrité, Confidentialité (DIC)] Les trois critères fondamentaux de la sécurité de l'information (Ch.~\ref{chap:cia}).
    \item[Gouvernance de la sécurité] Responsabilités et pratiques par lesquelles la direction oriente la sécurité (Ch.~\ref{chap:gouvernance}).
    \item[Mesure de sécurité] Moyen (organisationnel, humain, physique, technique) de traiter un risque (Ch.~\ref{chap:iso27002}).
    \item[Non-conformité] Écart entre une exigence et la réalité, majeure ou mineure (Ch.~\ref{chap:audit}).
    \item[Partie intéressée] Personne ou organisme ayant une attente vis-à-vis du SMSI (Ch.~\ref{chap:clauses46}).
    \item[PCA / PRA] Plan de continuité d'activité~; plan de reprise d'activité (Ch.~\ref{chap:continuite}).
    \item[PDCA] Cycle Planifier-Déployer-Contrôler-Agir, moteur de l'amélioration continue (Ch.~\ref{chap:pdca}).
    \item[Politique de sécurité (PSSI)] Document fondateur fixant objectifs, règles et responsabilités (Ch.~\ref{chap:politique}).
    \item[Revue de direction] Examen périodique du SMSI par la direction, qui décide des améliorations (Ch.~\ref{chap:audit}).
    \item[Risque] Combinaison de la probabilité d'un événement et de son impact (Ch.~\ref{chap:vocabrisque}).
    \item[Risque résiduel] Risque subsistant après traitement, accepté par la direction (Ch.~\ref{chap:traitement}).
    \item[RTO / RPO] Délai de reprise visé~; perte de données maximale acceptable (Ch.~\ref{chap:continuite}).
    \item[SMSI] Système de management de la sécurité de l'information (Ch.~\ref{chap:prologue}).
    \item[Vraisemblance] Probabilité qu'un risque se réalise (Ch.~\ref{chap:vocabrisque}).
\end{description}

    %% Annexe B : Référentiel des outils
\chapter{Référentiel des Outils}
\label{ann:outils}

Les outils et modèles cités sont des exemples. Le cours insiste sur la
méthode~: l'outil sert la décision, il ne la remplace pas.

\begin{tableaucours}{L{4.2cm}L{5.4cm}L{2.8cm}}{%
    \tableauHeader{Outil ou modèle} &
    \tableauHeader{Usage} &
    \tableauHeader{Chapitre, lab}
}
Inventaire des actifs & Recenser et classer l'information & Ch.~\ref{chap:cia}, SGSI-LAB-01 \\
Modèle de politique & Rédiger la PSSI & Ch.~\ref{chap:politique}, SGSI-LAB-02 \\
Grille d'écart ISO 27001 & Mesurer la maturité (gap analysis) & Ch.~\ref{chap:clauses46}, SGSI-LAB-03 \\
Méthode EBIOS RM & Apprécier les risques & Ch.~\ref{chap:methodes}, SGSI-LAB-04 \\
Matrice des risques & Hiérarchiser et traiter & Ch.~\ref{chap:traitement}, SGSI-LAB-05 \\
Modèle de DdA & Relier risques et mesures & Ch.~\ref{chap:dda}, SGSI-LAB-06 \\
Matrice RBAC & Contrôler les accès par rôle & Ch.~\ref{chap:technique}, SGSI-LAB-07 \\
Plan de continuité & Sauvegarde 3-2-1, RTO/RPO & Ch.~\ref{chap:continuite}, SGSI-LAB-08 \\
Tableau de bord & Suivre les indicateurs & Ch.~\ref{chap:indicateurs}, SGSI-LAB-09 \\
Fiche de non-conformité & Conduire un audit interne & Ch.~\ref{chap:audit}, SGSI-LAB-10 \\
\end{tableaucours}

\begin{boxalert}
\textbf{Rappel.} Une mesure ne se retient que si un risque la justifie. Le
catalogue de l'Annexe~A n'est pas une liste à appliquer en bloc~: c'est une
réponse au risque apprécié.
\end{boxalert}

    %% Annexe C : Modèles et procédures
\chapter{Modèles et Procédures}
\label{ann:templates}

\section*{Fiche d'actif}
\begin{tableaucours}{L{3.4cm}L{8.8cm}}{%
    \tableauHeader{Rubrique} & \tableauHeader{Contenu}
}
Actif & Information et support \\
Propriétaire & Qui en répond \\
Valeur / classification & Niveau de sensibilité \\
Critères prioritaires & D, I, C \\
\end{tableaucours}

\section*{Ligne de déclaration d'applicabilité}
\begin{tableaucours}{L{3.4cm}L{8.8cm}}{%
    \tableauHeader{Rubrique} & \tableauHeader{Contenu}
}
Mesure (Annexe A) & Intitulé de la mesure \\
Retenue & Oui / non \\
Justification & Risque couvert ou motif d'exclusion \\
État & En place, en cours, planifiée \\
\end{tableaucours}

\section*{Fiche de risque}
\begin{tableaucours}{L{3.4cm}L{8.8cm}}{%
    \tableauHeader{Rubrique} & \tableauHeader{Contenu}
}
Actif et menace & Ce qui est visé, par quoi \\
Vulnérabilité & La faiblesse exploitée \\
Probabilité / impact & Niveaux estimés \\
Traitement & Stratégie, mesure, responsable \\
Risque résiduel & Niveau, acceptation \\
\end{tableaucours}

\section*{Plan d'une politique de sécurité (PSSI)}
\begin{enumerate}
    \item Introduction et contexte (objectifs, portée, références).
    \item Gouvernance (organisation, rôles, comité).
    \item Classification et manipulation de l'information.
    \item Contrôle d'accès.
    \item Sécurité des opérations (sauvegarde, journalisation, incidents).
    \item Conformité (audits, revue, sanctions).
\end{enumerate}

    %% Annexe D : Textes juridiques et normatifs
\chapter{Textes Juridiques et Normatifs}
\label{ann:textes}

Seuls les éléments vérifiés dans le texte officiel sont cités dans les
chapitres. Ce tableau indique où les trouver.

\section*{Droit camerounais et africain}
\begin{tableaucours}{L{5.0cm}L{4.6cm}L{2.6cm}}{%
    \tableauHeader{Texte} & \tableauHeader{Objet} & \tableauHeader{Chapitre}
}
Loi n° 2010/012 du 21 déc. 2010 & Cybersécurité et cybercriminalité & \ref{chap:enjeux} \\
Loi n° 2024/017 du 23 déc. 2024 & Protection des données personnelles & \ref{chap:conformite} \\
Convention de Malabo (UA, 2014) & Cybersécurité et données, Afrique & \ref{chap:conformite} \\
\end{tableaucours}

\section*{Normes et référentiels internationaux}
\begin{tableaucours}{L{4.4cm}L{5.2cm}L{2.6cm}}{%
    \tableauHeader{Référence} & \tableauHeader{Objet} & \tableauHeader{Chapitre}
}
ISO/IEC 27000 & Vocabulaire et vue d'ensemble & \ref{chap:iso27000} \\
ISO/IEC 27001 & Exigences du SMSI & \ref{chap:iso27000} \\
ISO/IEC 27002 & Mesures de sécurité & \ref{chap:iso27002} \\
ISO/IEC 27005 & Gestion des risques & \ref{chap:methodes} \\
ISO 31000 & Management du risque & \ref{chap:vocabrisque} \\
ISO 22301 & Continuité d'activité & \ref{chap:continuite} \\
ISO 19011 & Audit des systèmes de management & \ref{chap:audit} \\
NIST CSF 2.0 & Cadre de cybersécurité & \ref{chap:iso27000} \\
EBIOS Risk Manager (ANSSI) & Appréciation des risques & \ref{chap:methodes} \\
\end{tableaucours}

    %% Annexe E : Bibliographie commentée
\chapter{Bibliographie Commentée}
\label{ann:bibliographie}

Pour chaque partie, quelques lectures. Les références complètes figurent dans
la bibliographie du livre.

\section*{Partie I · Fondements du management}
\begin{itemize}
    \item ISO/IEC 27000~\cite{iso27000}~: vocabulaire et vue d'ensemble.
    \item ISO/IEC 27001~\cite{iso27001}~: l'approche processus et le PDCA.
    \item ISO 31000~\cite{iso31000}~: le management du risque.
\end{itemize}

\section*{Partie II · Gouvernance et politique}
\begin{itemize}
    \item COBIT 2019~\cite{cobit2019}~: la gouvernance du SI.
    \item Lois 2010/012~\cite{loi2010012} et 2024/017~\cite{loi2024017}~; Malabo~\cite{malabo2014}.
    \item Guide d'hygiène ANSSI~\cite{anssihygiene2017}.
\end{itemize}

\section*{Partie III · Le SMSI ISO/IEC 27001}
\begin{itemize}
    \item ISO/IEC 27001~\cite{iso27001}~: les clauses 4 à 10.
    \item ISO/IEC 27002~\cite{iso27002}~: les mesures de l'Annexe~A.
    \item NIST CSF 2.0~\cite{nistcsf2024}~: un cadre complémentaire.
\end{itemize}

\section*{Partie IV · Les risques}
\begin{itemize}
    \item ISO/IEC 27005~\cite{iso27005}~: la gestion des risques.
    \item EBIOS Risk Manager~\cite{ebios2018}~: la méthode par scénarios.
\end{itemize}

\section*{Parties V et VI · Mise en œuvre, contrôle et audit}
\begin{itemize}
    \item ISO 22301~\cite{iso22301}~: la continuité d'activité.
    \item ISO 19011~\cite{iso19011}~; ISO/IEC 27006~\cite{iso27006}~: l'audit et la certification.
    \item ISO/IEC 27004~\cite{iso27004}~: la mesure~; auditologie~\cite{minka2026auditologie}.
\end{itemize}

    %% Annexe F : Contacts et ressources
\chapter{Contacts et Ressources}
\label{ann:contacts}

Les adresses sont celles des sites institutionnels à la date de rédaction~; à
vérifier.

\begin{tableaucours}{L{4.4cm}L{4.8cm}L{3.2cm}}{%
    \tableauHeader{Organisation} & \tableauHeader{Rôle} & \tableauHeader{Site}
}
ANTIC & Sécurité électronique, CERT national (Cameroun) & \texttt{www.antic.cm} \\
Union africaine & Convention de Malabo & \texttt{au.int} \\
ISO & Normes ISO et ISO/IEC & \texttt{www.iso.org} \\
NIST & Cadre de cybersécurité, SP 800 & \texttt{www.nist.gov} \\
ANSSI & EBIOS RM, guide d'hygiène & \texttt{cyber.gouv.fr} \\
ISACA & COBIT, certifications & \texttt{www.isaca.org} \\
\end{tableaucours}

\begin{boxinfo}
\textbf{La plateforme du cours.} Les fiches d'auto-évaluation et les dix labs
de NDAP SUKULU pour ce cours sont accessibles par le code \texttt{IS498}.
\end{boxinfo}


% --------------------------------------------------------------------------
% ÉLÉMENTS FINAUX
% --------------------------------------------------------------------------
\backmatter

    \nocite{*}
    \printbibliography[heading=bibintoc, title={Références bibliographiques}]

    \printindex

% ==============================================================================
% 4ÈME DE COUVERTURE
% ==============================================================================
\clearpage
\thispagestyle{empty}
\pagecolor{CouleurPrimaire}

\begin{tikzpicture}[remember picture, overlay]
    \fill[CouleurAccent]
        (current page.south west) rectangle
        ([yshift=1.8cm] current page.south east);
    \fill[CouleurAccent]
        ([yshift=-1.4cm] current page.north west) rectangle
        ([yshift=-1.7cm] current page.north east);
\end{tikzpicture}

\vspace*{2.5cm}

\begin{center}
    {\color{CouleurAccent}\rule{0.6\linewidth}{1pt}}\\[1.5ex]
    {\color{white}\Large\bfseries À propos de ce cours}\\[1.5ex]
    {\color{CouleurAccent}\rule{0.6\linewidth}{1pt}}
\end{center}

\vspace{1cm}

\begin{center}
\begin{minipage}{0.75\textwidth}
\color{white!85!gray}\setstretch{1.4}
Quand un système d'information s'arrête, la question décisive n'est pas
technique : elle porte sur qui avait décidé, et qui devait rendre compte.
Ce cours apprend à y répondre avant la panne.

\medskip

Il progresse en six parties, précédées d'un prologue : les fondements de la
gouvernance des systèmes d'information, le référentiel COBIT~2019, le cadre
normatif et juridique, ITIL~4 et les processus de gestion des services, les
défis de l'agilité, du nuage et de l'intelligence artificielle, et enfin la
recherche.

\medskip

Chaque notion est ancrée à la fois dans les standards internationaux (ISO,
ISACA, IIA, INTOSAI) et dans le droit camerounais et communautaire (CEMAC,
OHADA). Le cours s'adresse aux étudiants du Master Professionnel en
Management des Systèmes d'Information et aux professionnels qui exercent comme
DSI, responsables GRC, auditeurs des SI, responsables des services
informatiques ou RSSI.
\end{minipage}
\end{center}

\vfill

\begin{center}
    {\color{CouleurAccent}\rule{0.7\linewidth}{1.2pt}}\\[1.2ex]
    {\color{white}\large\bfseries MINKA MI NGUIDJOÏ Thierry Emmanuel}\\[0.4ex]
    {\color{white!80!gray}\small
        CISA~·~CISM~·~CGEIT~·~CRISC~·~CDPSE~·~ISO 27001~·~ISO 22301~·~ISO 9001}\\[0.3ex]
    {\color{white!70!gray}\footnotesize\itshape
        Chercheur au laboratoire d'Ingénierie Mathématiques et des Systèmes d'Information}\\[0.3ex]
    {\color{white!70!gray}\footnotesize\itshape
        Chercheur~|~Cryptologue~|~Auditeur Senior~|~Expert Judiciaire~|~GRC \& Sécurité de l'Information}\\[1.2ex]
    {\color{CouleurAccent}\rule{0.7\linewidth}{1.2pt}}
\end{center}

\vspace{1cm}

\afterpage{\nopagecolor}
\end{document}
